% =====================================================================
%  GENERATED; DO NOT EDIT BY HAND.  Rebuild: sh ebook-method/religionsfilosofiske-arbejder-1903/rebuild.sh
%  Harald Høffding, Om nogle religionsfilosofiske Arbejder fra den nyeste Tid.  Indbydelsesskrift til Kjøbenhavns
%  Universitets Aarsfest i Anledning af Hans Majestæt Kongens Fødselsdag den 8. April 1903.  Kjøbenhavn: Universitets-
%  bogtrykkeriet (J. H. Schultz), 1903.  83 pp. of text + Noter pp. 84-86.  Diplomatic transcription, 2026-10-26.
%  TEXT OF RECORD: the 1903 print.  Reprinted, revised by Høffding, in Mindre Arbejder. Anden Række (1905), pp. 58-140
%  (texts/hoeffding/mindre-arbejder-2); the differences are listed in ebook-method/religionsfilosofiske-arbejder-1903/variants.tsv.
%
%  (1) SOURCE.  ~/bibliotek/Høffding, Harald/1903-om-nogle-religionsfilosofiske-arbejder.pdf (SCANS.tsv, sha256 a76c3050...):
%  Google Books' PDF (id yqYPAAAAQAAJ) of the Taylor Institution (Oxford) copy, Google OCR layer.  Printed p. = PDF - 12
%  for pp. 1-104 (every folio checked in the layer).  PDF 7 the invitation's title page (the University's seal; pencil
%  library marks "Scand Pi 16" not transcribed), PDF 9 the title, PDF 11 the Indhold: transcribed by hand and checked at
%  the image.  The Indhold and the heads differ, both kept as printed: Indhold "a. Gennembrud og Kontinuitet. (Starbuck)"
%  / head "a. Gennembrud og Kontinuitet i den religiøse Udvikling."; "b. Temperamentets Indflydelse. (Coe)" / "...
%  Indflydelse."; "c. Forskelligheder indenfor én Konfession. (Leuba, Arréat)" / "c. Forskelligheder i Forholdet til
%  Konfessionen."; "d. Poesi og Religion. (Santayana)" / "d. Poesi og Religion."; "I. SYSTEMATISKE FORSØG" / "I. Systematiske
%  Forsøg." (and so for II); the Indhold's sections are in capitals, the heads in roman, the lettered heads in italics.
%  NOT TRANSCRIBED: pp. 87-104 of the Indbydelsesskrift are not Høffding's -- the University's reports on the prize essays
%  for 1901-02 (pp. 87-100, "Prisafhandlingerne for Aaret 1901-1902"), its prize questions for 1902-03 (pp. 101-103) and
%  the invitation itself (p. 104: the festival in the University's hall, Tuesday 21 April 1903, the Rector Harald Høffding
%  giving the address; dated 11 April 1903).
%
%  (2) CONVENTIONS.  Page turns "% --- p. N ---" + \opage{N} at the page's first word (a word divided at a turn ends in %).
%  Running heads (the folio alone) and the sheet signatures (a figure at the foot every 8th page) are not set.  The NOTES
%  are endnotes (Noter, pp. 84-86, smaller type), marked in the text by a raised figure and a parenthesis, set
%  \textsuperscript{N}); their labels ("1) S. 1.", raised in the print) are set plain.  Italic -> \textit; names in small
%  capitals -> \textsc.  „ “ as printed; the long dash ---; points of omission spaced as printed (. . . .).  The drop
%  capital V of p. 1 set as a plain letter.  Misprints are kept and marked "% PRINTED AS IS" at the line (8: Erkendslse
%  p. 20, tibage p. 21, frany p. 22, betoner. p. 28, 1 for I p. 49, urandsagelige p. 54, skjulteLiv p. 54, forekomme.. p. 71);
%  d'aprés (p. 18) and "Varieties of religious belief" (note 12, for "experience") are as printed.  The ø rule (Hans,
%  2026-10-08): the damaged type of p. 20 ("religiese") is read religiøse.  "stode" (pp. 55, 76) is right (plural past of staa).
%
%  (3) ACCURACY (2026-10-26).  Method: the draft is Google's layer by rule (ebook-method/religionsfilosofiske-arbejder-1903/
%  draft.py); ebook-method/collate/check.py collated it word by word with punctuation against an independent tesseract
%  reading (dan, 300 dpi): 467 disagreements in 28,054 words, 43 settled by check.py's own rules, 424 decided (decisions.tsv),
%  OPEN 0.  THE 1905 REPRINT AS A THIRD READING (rules.py, w1905.py): where the 1905 text agreed with exactly one of the two
%  OCRs it decided (203 words, 89 dashes, 1 mark); never on a word division or a case difference, which 1905 revised.  An
%  audit sample of 38 of those decisions read at crops: 0 wrong.  dash+ (a long dash tesseract alone has) 15; quote/hyphen/
%  para rules 11.  The remaining 106 word checks were read at crops (9 sheets, main session).  Italics by stroke slant
%  (check.py --italics): 93 mismatches, 39 settled where 1905 sets the same words in italics, 46 read at crops; then the
%  italic EXTENTS against 1905 (italext.py; the slant test starts a run late) read at crops: 20 widened (patches.LAST).
%  THE 1903/1905 DIFFERENCES: every word difference between this text and the 1905 transcription (variants.py) was read
%  at a crop of the 1903 scan (var-NN sheets; 115 of 121 variant rows, the other 6 punctuation on which both OCRs agree);
%  that reading found 12 faults of this transcription that both OCRs shared (points of omission lost 9 times, a dash, ame-|rikanske,
%  Værdi-|erne) and 7 of the 8 misprints marked.  Further checks (lessons of Mindre Arbejder III): 17 paragraph-initial dashes moved
%  to the end of the line before (the dash test; patches.DASH_TO_PREV); 15 pages whose last word is divided at the turn
%  (pagemap.TURN_DIVIDED); 2 false paragraph breaks (p. 1 after the drop capital, p. 44); no suspended compounds, no elided
%  French articles lost, quote balance clean except quotations crossing pages; the layer-against-transcription check
%  found no lost line.  BLIND PAGE CHECK: one Sonnet reader (103k tokens), pages 4, 26, 32, 42, 69 (seed 2026102604, three
%  strata) and the whole Noter pp. 84-86: 2,681 words, 0 word errors, 0 punctuation errors, 10 italic words (6 runs) not
%  italic -- mended.  Endnote marks 6 and 9 the layer gave as bare ")" -- set from the crops.
%  NOT DONE: the remaining italic runs were not read one by one beyond the slant test, the 1905 comparison and the 8-page
%  check (expect a few short italic words missed); points of omission elsewhere were checked only where 1905 differs.
%
%  (4) 1903 AND 1905.  variants.tsv lists 121 differences (1903 page, 1905 page, both readings, kind, how each was read):
%  23 of wording (e.g. kunne > kan p. 5, Nyskabelser > Nydannelser p. 9, the anonymous "teologisk Kritiker, hvis Standpunkt
%  ellers er højst forskelligt fra mit" named as Pastor L. P. Larsen p. 4, p. 173 > § 66 p. 21, var > er p. 57, Intuitioner
%  > Institutioner p. 59, Tilstrømmen > Tilstrømning p. 79), 21 of spelling (særlig > særligt, Juli > July, extreme >
%  ekstreme, hovedkuls > hovedkulds; 1903's misprints corrected in 1905), 31 of word division (selv om > selvom 18 times,
%  ud over > udover, des mere > desmere, for at > forat), 42 of punctuation, 4 of case.  Not listed row by row: the
%  quotation marks („ “ in 1903, » « in 1905), the headings (in capitals in 1905), names in small capitals, the points of
%  omission, and 1903's endnote marks.  THE NOTES: 1903 has 13 endnotes; 1905 has none of them as notes: note 2's
%  reference went into the text (p. 4: "(Pastor L. P. Larsen i Studenterhjemmets Julebog 1901)", its polemic dropped), note 4
%  into the text as "(DOWDEN: Puritan and Anglican. London 1900. p. 234)" (p. 10, the cross-reference dropped); notes 1, 3
%  and 5-13 were left out; 1905 adds one footnote of its own (Moderne Filosofer, 1904, at 1903 p. 11).  The 1905
%  transcription (mindre-arbejder-2) carries 14 faults found by this comparison (variants-1905-faults.tsv), not mended there.
% =====================================================================
\documentclass[12pt,a4paper,openany]{book}
\usepackage[utf8]{inputenc}
\usepackage[T1]{fontenc}
\usepackage{libertinus}
\usepackage{libertinust1math}
\usepackage[danish]{babel}
\usepackage[protrusion=true,expansion=true]{microtype}
\usepackage{setspace}
\usepackage[margin=1.3in]{geometry}
\usepackage{fancyhdr}
\usepackage{hyperref}
\hypersetup{pdftitle={Om nogle religionsfilosofiske Arbejder fra den nyeste Tid}, pdfauthor={Harald Høffding}, hidelinks}
\pagestyle{fancy}
\fancyhf{}
\fancyhead[LE,RO]{\thepage}
\fancyhead[C]{\textit{Høffding: Om nogle religionsfilosofiske Arbejder (1903)}}
\setlength{\headheight}{15pt}
\setstretch{1.3}
\usepackage{marginnote}
\newcommand{\opage}[1]{\marginnote{\footnotesize\textit{[#1]}}}
\newcommand{\whiteline}{\par\vspace{\baselineskip}\par}
\newcommand{\kapitel}[3]{\chapter*{\centering\normalsize\if\relax\detokenize{#1}\relax\else #1\\[1ex]\fi\Large #2}\addcontentsline{toc}{chapter}{#3}}   % [numeral]{title}{toc}
\newcommand{\afsnit}[2]{\section*{\centering\normalsize\bfseries #1}\addcontentsline{toc}{section}{#2}}   % a numbered section, bold
\newcommand{\underafsnit}[2]{\subsection*{\centering\normalsize\normalfont\itshape #1}\addcontentsline{toc}{subsection}{#2}}   % a lettered subsection, all italic
\newcommand{\streg}{\begin{center}\rule{3em}{0.4pt}\end{center}}
\newcommand{\ix}[3]{\par\noindent\hspace*{#1}#2\dotfill #3}
\begin{document}

\frontmatter
% PDF 7: the invitation's title page; PDF 9: the title page; PDF 11: the Indhold (PDF 8, 10, 12 blank).
\begin{titlepage}\centering
{\Large INDBYDELSESSKRIFT}\\[1ex]{\small TIL}\\[1ex]{\Large KJØBENHAVNS UNIVERSITETS}\\[1ex]{\Large AARSFEST}\\[1ex]
{\small I ANLEDNING AF}\\[1ex]{\large HANS MAJESTÆT KONGENS}\\[1ex]{\large FØDSELSDAG}\\[1ex]{\small DEN 8. APRIL 1903}\par\vspace{2cm}
{\large HARALD HØFFDING}\\[1ex]{\large OM NOGLE RELIGIONSFILOSOFISKE ARBEJDER}\\[0.5ex]{\large FRA DEN NYESTE TID}\par\vfill
{\small KJØBENHAVN}\\[0.5ex]{\small TRYKT I UNIVERSITETSBOGTRYKKERIET (J. H. SCHULTZ)}\\[0.5ex]1903
\end{titlepage}
\begin{titlepage}\centering\vspace*{2cm}
{\Large OM NOGLE}\\[1ex]{\Huge RELIGIONSFILOSOFISKE}\\[1ex]{\Large ARBEJDER}\\[1ex]{\large FRA DEN NYESTE TID}\par\vspace{1.5cm}
{\small AF}\par\vspace{1ex}{\large HARALD HØFFDING}\par\vfill
{\small KJØBENHAVN}\\[0.5ex]{\small TRYKT I UNIVERSITETSBOGTRYKKERIET (J. H. SCHULTZ)}\\[0.5ex]1903
\end{titlepage}
\chapter*{INDHOLD}
\noindent\hfill Side
\ix{0em}{INDLEDNING}{1}
\ix{0em}{I. SYSTEMATISKE FORSØG}{12}
\ix{1.5em}{1. Hermann Siebeck}{12}
\ix{1.5em}{2. Auguste Sabatier}{18}
\ix{1.5em}{3. Rudolph Eucken}{21}
\ix{0em}{II. RELIGIONSPSYKOLOGISKE UNDERSØGELSER}{32}
\ix{1.5em}{1. Om religiøse Domme. (Max Reischle).}{35}
\ix{1.5em}{2. Den religiøse Følelses Patologi. (Murisier).}{38}
\ix{1.5em}{3. Bidrag til den konfessionelle Religiøsitets Psykologi.}{41}
\ix{3em}{a. Gennembrud og Kontinuitet. (Starbuck).}{46}
\ix{3em}{b. Temperamentets Indflydelse. (Coe).}{51}
\ix{3em}{c. Forskelligheder indenfor én Konfession. (Leuba, Arréat).}{54}
\ix{3em}{d. Poesi og Religion. (Santayana)}{65}
\ix{3em}{e. Sammenligning mellem forskellige Typerækker}{67}
\ix{1.5em}{4. Komparativ Religionspsykologi. (James).}{69}
\ix{3em}{a. De tre Synspunkter}{71}
\ix{3em}{b. Den religiøse Erfarings Natur}{75}
\ix{3em}{c. De vigtigste Typer}{76}
\ix{3em}{d. Den religiøse Grundtanke.}{79}
\ix{3em}{e. Bemærkninger om James' Religionsfilosofi}{80}
\par\medskip\noindent Noter S. 84.
\clearpage


\mainmatter
% --- p. 1 ---
\kapitel{}{Indledning.}{Indledning}\opage{1}

Ved Religionsfilosofi forstaar jeg ikke et Forsøg paa at sætte
Spekulation i Religionens Sted, endnu mindre et Forsøg
paa ad Spekulationens Vej at begrunde religiøse Antagelser, men
et Forsøg paa at belyse Religionens Stilling i Tidens Aandsliv.
Religionen er her Genstand for Undersøgelse som en faktisk
given Form for aandeligt Liv, hvis Beskaffenhed, Betydning og
Berettigelse det gælder at bestemme. Spørger man nu, hvilke
Synspunkter Filosofien raader over til Brug ved en saadan Undersøgelse,
saa kan der sikkert kun være Tale om tre Synspunkter,
--- men det er egentligt ogsaa alle de Synspunkter, Filosofien
idethele raader over: det logiske (eller erkendelsesteoretiske), det
psykologiske og det etiske\textsuperscript{1}). Da enhver Religion --- hvor vid
eller hvor snever en Betydning man saa tager Ordet i --- faar
Udtryk i visse Forestillinger, opstaar naturligt Spørgsmaalet om
disse Forestillingers Gyldighed, deres Erkendelsesværdi. Indeholde
de en konsekvent og gyldig Erkendelse af Tilværelsen, og
kunne de bringes i Overensstemmelse med de Forestillinger, som
Videnskaben ad sine egne Veje er naat til at fastholde som
gyldige? Selv om nu en saadan \textit{erkendelsesteoretisk} Undersøgelse
skulde føre til det Resultat, at de religiøse Forestillingers Betydning
ikke kan bestaa i at skaffe os en videnskabelig Forstaaelse
af Tilværelsen, er dermed Spørgsmaalet om deres Værdi
og Berettigelse ikke afgjort. Den kan bero paa deres Forhold
til andre Sider af Sjælelivet end den rent intellektuelle. Og
% --- p. 2 ---
\opage{2}direkte Iagttagelse viser jo ogsaa, at Følelse, Trang og Villie
spille nok saa stor en Rolle i Religionen som Erkendelse. Der
er da al Anledning til at foretage en \textit{psykologisk} Undersøgelse
af det religiøse Liv, saaledes som Iagttagelse lærer os det at
kende i Samtiden og historisk Overlevering i Fortiden. Denne
Undersøgelse gaar for det Første ud paa at finde de stadigt virkende
Motiver i Religionen og at udskille dem af de skiftende
Former og Forbindelser, i hvilke de under forskellige Forhold
fremtræde. Kun i disse Motiver kan det Væsentlige i Religionen
være at søge, --- det, der, saalænge det besidder Livskraft,
stedse vil kunne ytre sig i nye Skikkelser. Der existerer ingen
Religion i al Almindelighed; men Erfaringen viser os visse Grundmotiver
i en Rigdom af forskellige Forbindelser eller Variationer.
Derfor er det for det Andet Religionspsykologiens Opgave
at opsøge og beskrive de vigtigste Typer for Religion og religiøst
Liv. Dersom nu saaledes Motiver og Typer ere paaviste, er det
den \textit{etiske} Undersøgelses Sag at afgøre, hvilken Værdi Religionen
som ejendommelig Form for menneskeligt Aandsliv har for Menneskelivet
i dets hele, store Sammenhæng. Maalestokken er her
den samme, som vi ogsaa ellers anvende, naar vi drøfte en menneskelig
Egenskabs eller Handlings Værdi. Vi undersøge deres
Forhold til visse Grundværdier, i sidste Instans til en eneste
Grundværdi, som for os er den afgørende. Hvilken denne Grundværdi
er, er det Etikens Sag at udfinde. Hvis der er Uenighed
om den Grundværdi, der bestemmer, hvad vi i de enkelte Tilfælde
kalde værdifuldt, saa vil der ogsaa være Uenighed om Vurderingen
af Religionen. Spørgsmaalet vilde naturligvis være
afgjort paa Forhaand, naar man gik ud fra, at Religionen (d. v. s.
den Religion, man selv har eller mener at have) selv er Grundværdien,
efter hvilken Alt skal maales. Men saa vilde der heller
ikke bestaa noget religiøst Problem eller være Mulighed for
eller Trang til nogen Religionsfilosofi, ligesaa lidt som der vilde
være Mulighed for eller Trang til nogen Erkendelsesteori, naar
man gik ud fra, at en eller anden Religion indeholdt den absolute
Sandhed, efter hvilken Alt, hvad man ellers vilde kalde
sandt, skulde maales. Et Problem opstaar der derimod her, saasnart
man bliver sig bevidst, at man ved Anerkendelsen af Re%
% --- p. 3 ---
\opage{3}ligionen forudsætter en vis Maalestok for gyldig, en Maalestok,
som ikke er hentet fra selve Religionen alene, og især, naar
man opdager, at man paa det religiøse Omraade anvender en
Maalestok, der er forskellig fra den, man anvender paa andre
Omraader.

Af de tre Synspunkter, fra hvilke det religiøse Problem kan
og maa behandles, er det det psykologiske, der i den nyeste Tid
især tildrager sig Opmærksomheden. Det er naturligvis ikke
noget Nyt, at Religionen opfattes som et psykologisk Fænomen.
Enhver dybere gaaende religiøs Bevægelse maatte føre dertil,
især ved at gøre det tydeligt, at det kommer an paa ikke blot
at staa i Forestillingsforhold til det Guddommelige, men ogsaa
at bevæges af det i inderlig Følelse. Mere og mere henpeges
der fra religiøs Side til, at Religionens sidste Grundlag er menneskelig
Livserfaring. Religionspsykologien anerkender fuldtud
dette, og den stiller sig da den Opgave at beskrive denne Erfaring.
Den gaar herved ud fra, at skal Beraabelsen paa Erfaring
have nogen Mening, maa man ved Erfaring forstaa umiddelbar
Oplevelse, og det kommer da an paa at bestemme, hvad
og hvor Meget der umiddelbart kan opleves. Det gælder at
skelne imellem, hvad der virkeligt er umiddelbart givet, og hvad
der føjes til som Forklaring. Forklaringen hentes fra en Forestillingskreds,
som ligger Mennesket nær. Og fra Forklaringen
maa igen skelnes Udtryksmaaden. Men disse tre Ting, Oplevelse,
Udtryksmaade og Forklaring, er det særligt ved den religiøse
Erfaring, der væsentligt er Følelseserfaring, overordentligt vanskeligt
at holde ude fra hverandre. Religionspsykologisk kommer
det da an paa at opsøge simple Exempler, hvor Forskellen
er tydelig, og ved deres Hjælp at forklare sig de mere indviklede
Tilfælde. Men saa Meget vil strax være klart, at der er en
lang Vej fra en umiddelbar religiøs Erfaring til de ofte meget
subtile Dogmer eller til de digterisk udformede Myter og Legender,
som de historisk foreliggende Religioner fastholde i deres
Overleveringer. Og der maa ganske særligt skelnes mellem
Oplevelser, som fremkaldes ved saadanne overleverede Forestillinger,
under Forsøget paa at omsætte dem til selvvunden Erfaring,
og Oplevelser, som paa første Haand føre til visse bestemte
% --- p. 4 ---
\opage{4}Forestillingers Dannelse. Kun denne sidste Art er oprindelig religiøs
Erfaring. Og selv ved den maa der endnu spørges, om
ikke den ene og samme Oplevelse kunde udtrykkes og forklares
ved andre Forestillinger og i andre Former, end det faktisk sker.
Kun hvis en saadan Kontraprøve kunde gøres, og hvis den saa
førte til et negativt Resultat, vilde Erfaringsbeviset være ført.

Jeg tror, at denne psykologiske Metode vil faa en meget
stor Betydning under den fortsatte religiøse Diskussion. Det er
med fuld Ret, man fra religiøs Side henviser til personlig Livserfaring
som sidste Kriterium overfor den ofte meget overfladiske
Beraabelse paa Fornuftens eller Videnskabens formentlige Resultater.
Men hvis man ikke behandler selve Begrebet Erfaring
med streng Kritik, sondrende bestemt mellem det Givne og dets
Fortolkning, saa indfører man en Dogmatisme, der er langt værre
end den gamle rationalistiske Dogmatisme, man fornemt ser ned
paa, saavist som det er værre at bedrage sig selv og Andre med
Hensyn til sit Livs inderligste Forhold end at begaa logiske
Fejl i et indviklet Ræsonnement. I hvert Tilfælde kan Psykologien
ikke lade sig hindre i at forsøge at analysere de ved
første Øjekast umiddelbare og usammensatte Fænomener, man
beraaber sig paa. Der ligger her et stort Materiale til Rede i
den religiøse Bevidstheds Udtalelser til forskellige Tider og under
højst forskellige Kulturforhold, tiltrækkende og forfriskende
især for dem, der have faaet nok af den hidtil gængse Maade at
behandle de religiøse Spørgsmaal paa, baade hos Angriberne og
hos Forsvarerne.

Det er dette, jeg efter bedste Evne har søgt at indskærpe og
gennemføre i min Religionsfilosofi, hvis mest omfattende Afsnit
netop er psykologisk. Jeg har da ogsaa havt den Tilfredsstillelse,
at en teologisk Kritiker, hvis Standpunkt ellers er højst forskelligt
fra mit, har erklæret, at min Metode er den ene rette\textsuperscript{2}).
Men des større maa da min Forundring være, naar jeg ser, at
denne samme Kritiker bestrider mine Resultater --- ved Hjælp
af en ganske anden Metode end den, han saaledes har anerkendt!
Ti han bruger Begrebet Erfaring netop i den vage og vide Betydning,
i hvilken det ofte bruges i daglig Tale. Han stiller
gennemgaaende Spørgsmaalet saaledes: have visse overleverede
% --- p. 5 ---
\opage{5}Dogmer Værdi for Livet? Men en gennemført religionspsykologisk
Undersøgelse gaar ogsaa den omvendte Vej og spørger:
forsvinder der virkeligt Livsværdier, fordi man ikke antager visse
Dogmer, eller fordi man sætter poetiske Symboler i Stedet for
Dogmerne? Jeg selv tør (uagtet jeg ikke lægger Skjul paa mit
eget personlige Standpunkt i denne Sag) ikke bestemt besvare
dette Spørgsmaal; jeg har udtrykkeligt fremhævet Muligheden
af, at der her kunde være et Værditab. Denne Mulighed kan
ingen Videnskab modbevise, og derfor grunder Antagelsen af
Værdiens Bestaaen sig paa en Tro, selv om man ikke udstrækker
denne Antagelse ud over den menneskelige Verdens Grænser. ---

Paa et Omraade, hvor endnu --- og maaske til evig Tid
--- forskellige Retninger brydes, vil det være af stor Betydning at
undersøge, om den anvendte Fremgangsmaade er den eneste,
der kan være Tale om. Derfor har jeg ønsket at fremdrage en
Række religionsfilosofiske Forsøg, af hvilke de fleste først ere
fremkomne efter Afslutningen af min Bog, for at se, om der er
andre Synspunkter at anlægge end de tre nævnte (det logiske,
det psykologiske og det etiske), og specielt for at se, om den
forsøgte Bestemmelse af religiøs Erfaring og af Religionens
Grundtanke skulde trænge til Supplering eller Berigtigelse.

Ikke blot hvad Metoden, men ogsaa hvad Resultaterne angaar,
vil en saadan Sammenligning med andre Forskeres samtidige
Arbejder kunne være af Betydning. Jeg tillader mig da
at minde om, at de Resultater, til hvilke jeg mener at være
--- Det religiøse Problem
naat, hovedsageligt ere følgende tre.
gaar ind under det almindelige Vurderingsproblem, og Metoden
for dets Behandling vil derfor mere og mere blive etisk Vurdering
ved Siden af psykologisk Beskrivelse. Ingen af disse to
Fremgangsmaader kunne undværes, da man jo kun kan vurdere
nøjagtigt, hvad der foreligger beskrevet nøjagtigt. --- Den religiøse
Trang kan betegnes som en Trang til at hævde Værdiens
Bestaaen, og den religiøse Erfaring bestaar af, hvad Mennesket
har oplevet angaaende Forholdet mellem de Værdier, det kender,
og den Virkelighed, det kender. De specielle religiøse Antagelser
angive --- hvad enten de fremtræde i Mytens, Dogmets,
Legendens eller Symbolets Form --- de Maader, paa hvilke Men%
% --- p. 6 ---
\opage{6}neskene have søgt dels at udtrykke, dels at forklare, hvad de
saaledes have oplevet. --- Der er i Aandslivets Vilkaar i den
nyere Tid Intet til Hinder for, at en saadan Trang kan vedblive
at ytre sig og at finde Tilfredsstillelse, selv om det mere og
mere maa ske i andre Former end i tidligere Tider, --- og selv
om det ikke er nødvendigt, at den ytrer sig hos Alle. Det er
en religionshistorisk Hypotese, at Sætningen om Værdiens Bestaaen
er den religiøse Grundtanke, d. v. s. at det er den, der udtrykker
den gennemgaaende Tendens i de historisk foreliggende
Former for Religion. Men Sætningen om Værdiens Bestaaen
kan ogsaa opstilles som en filosofisk Hypotese. Den betegner en
Livsanskuelse, der kan gøre Fordring paa at være \textit{forenelig med}
en strengt gennemført videnskabelig Verdensanskuelse, selv
om den ikke kan \textit{udledes} af denne. Livsanskuelsen staar hos Alle
med en vis Selvstændighed overfor Verdensanskuelsen. Der sker
en stadig Indvirkning af denne paa hin; men ingensinde er
Livsanskuelsen et blot Produkt af den videnskabeligt begrundede
Verdensanskuelse. Ti Livsanskuelsen indeholder et Plus,
der skriver sig fra Personlighedens Vurdering af den Tilværelse,
den søger at forstaa ved Videnskabens Hjælp. Nu kan der vel
indenfor Vurderingen findes større eller mindre Konsekvens,
og de Synspunkter, der lægges til Grund ved Vurderingen,
kunne være mere eller mindre naturlige og væsentlige. Men
paa Vurderingens Omraade ere vi i langt højere Grad nødte til
at være Empirikere, end hvor det gælder Forstaaelse og
Forklaring. Det Irrationale (i matematisk Betydning), som vor
Tænkning stedse maa lade staa tilbage, har et større Omfang
paa Vurderingens end paa den rene Teoris Omraade. Ud over
\textit{Muligheden} af en Livsanskuelse, der bygger paa Grundtanken
om Værdiens Bestaaen, kan Filosofien ikke føre os\textsuperscript{3}). Grundmotivet
til enhver Livsanskuelse udspringer af den praktiske Erfaring
om Værdiernes Skæbne i Virkelighedens Verden; og dersom
denne Erfaring fører til Tro paa Værdiens Bestaaen, saa
vil denne Tro ikke behøve at staa som en mystisk Hypotese,
blot Genstand for passiv Kontemplation, men den vil naturligt
blive en Arbejdshypotese, idet den ansporer til at finde nye
Værdier, naar de gamles Tid er forbi, eller til maaske endog selv
% --- p. 7 ---
\opage{7}at frembringe saadanne. Den giver da Mod til at arbejde baade
paa Tankens og paa Virkelighedens Omraade for at vinde fremskridende
Bekræftelse. Maaske ville ikke alle Individer af deres
Erfaringer føres til at nære en saadan Tro; det vil kunne være
tilstrækkeligt, om den holder sig i Slægten. Ti Slægtens Liv
behøver for at være rigt og kraftigt mange forskellige Individualiteter,
maaske endog en stadig Kamp mellem aandelige Retninger.
Mange Individer kunne være saa optagne af Tankens
og Virkelighedens specielle Arbejde, at de ikke faa Tid, Kraft
og Interesse tilovers til at danne sig nogen i bestemte Tanker
udtrykt Livsanskuelse. Der kan her, som i andre Henseender,
finde en Arbejdsdeling Sted. Hvorfor skulde Arbejdsdelingen
være indskrænket til Livets rent ydre Forhold? Er der ikke
Grund til at antage, at den griber ind i Menneskenes inderligste
Anliggender? Livet vil blive rigere derved, ligesom Kunstens
Verden bliver rigere, jo mere den enkelte Kunstner følger
sin individuelle Trang og Evne. Selve Sætningen om Værdiens
Bestaaen kræver maaske endog, at ikke Alle tro paa den eller
tænke paa den.

Disse sidste Bemærkninger lægger jeg særlig Vægt paa
overfor en Indvending, der er bleven rettet mod min Bog: at
den fremstiller en rent subjektiv Opfattelse. Jeg har ikke lagt
Skjul paa mit eget personlige Stade overfor de Spørgsmaal, jeg
drøfter; men Hovedsagen for mig har været at undersøge de
Muligheder, der under Nutidens Kulturforhold frembyde sig i
religiøs Henseende, og her var det ikke nok at holde sig til et
enkelt Individ. Det er en sammenlignende Undersøgelse, jeg
har villet anstille. Jeg har villet foretage en Sammenligning
mellem det Religiøses Stilling i Fortid og i Nutid, og mellem
dets Betydning for forskellige Personligheder og Standpunkter
i Nutiden. Det drejer sig i en Religionsfilosofi ikke om vedkommende
Filosofs egen Tro, men om de Muligheder for Tro,
han finder i Slægtens Aandsliv, saaledes som det nu kan føres
i Klarhed og i Dybde. Mindre og mindre har en Filosof Ret
til at antage, at Slægtens Psykologi er udtømt med hans egen
individuelle Psykologi.

En beslægtet Indvending er rejst mod min Fremstilling af
% --- p. 8 ---
\opage{8}en energisk tysk Forsker, som jeg har den Ære at kunne kalde
min personlige Ven. I en udførlig Anmeldelse af min Bog i
Tidsskriftet „Das freie Wort“ (1902) har \textsc{Ferdinand Tönnies}
kritiseret min Opfattelse af det religiøse Problem paa en Maade,
der forekommer mig at kunne belyse Forholdet mellem Religionshistorie
og Religionsfilosofi.

„For Guyau“, siger han, „og jeg tilstaar, at jeg deri slutter
mig til ham, er Religionen i første Linie Genstand for Socialpsykologi
eller Sociologi. For Høffding hører den næsten fuldstændigt
til Individualpsykologien, der jo falder sammen med
det, man plejer at forstaa ved Psykologi overhovedet. Han kalder
Buddhismen og Kristendommen „de højere Folkereligioner“:
i dem skal det indre psykologiske Drama fremtræde som det
egentlige Verdensdrama . . . . Nu er det jo umiskendeligt, at i
begge disse højst vigtige religiøse Systemer har den individuelle
Trang og den individuelle Tænkning faaet en overvældende Betydning.
Men hvorfra stammer denne Sjælens Kamp? Er den
en oprindelig og væsentlig Bestanddel af Religionen som Folketro
og Folkesæd, eller dog af den religiøse Følelse, der „bærer“
Tro og Sæd? \textit{Eller} hidrører den, netop hvad dens renere og
højere Indhold angaar, fra en anden Kilde, nemlig fra en Tænkning,
en Filosoferen, der efter sit Væsen er fremmed for Folketro
og Folkesæd, . . . skønt den ganske vist i hine to Folkereligioner
har opnaaet en Udligning med populære og vildtvoxende
(naturwüchsige) Forestillinger og Følelser, ja, i Kristendommen
endog har sluttet et evigt Forbund med dem?“

Konsekvensen af denne Tankegang synes at maatte blive
den, at Buddhismen og Kristendommen ikke egentligt ere Religioner.
Ja, man maatte vel egentligt gaa videre. Gives der
nogensomhelst positiv Religion, i hvilken der ikke kan paavises
Elementer af Tænkning? De mest elementære, for ikke at sige
de mest primitive Religioner bygge paa Animismen, Menneskets
mest elementære Filosofi, og saaledes vil paa alle Stadier af Udviklingen
et Folks Religion bære Præget af dets forudgaaende
„Filosofi.“ Historikerne have undertiden temmeligt romantiske
Forestillinger om den Umiddelbarhed og „Naturwüchsigkeit“,
hvormed Menneskene gaa op i deres Kultus, Myter og Dogmer.
% --- p. 9 ---
\opage{9}Der foregaar, som Religionshistorien selv viser, stadige Kombinationer,
Suppleringer og Metamorfoser, hvorved et mere eller
mindre dybtgaaende Reflexionsarbejde medvirker. Men hvad der
for mig er Hovedsagen: det er misvisende her blot at tale om
Tænkning og Filosofi. Det er først og fremmest personlig Livserfaring,
der øver sin Indflydelse ved disse Ændringer, der ofte
ere Nyskabelser. De Personligheder, som have været afgørende
for de højere Folkereligioners Tilblivelse og Udvikling, have
først og fremmest lagt deres egne Erfaringer om det Værdifulde
og dets Skæbne i Verden til Grund ved de Ændringer, de førtes
til at foretage i det Overleverede. Ingen Religion laver sig
selv. Religionen opstaar og existerer som Religion kun i Menneskers
Sind og Tanke. Det maa have været visse bestemte Mennesker,
der have formet Kultus, Myte og Dogme, ligesom det maa
have været visse bestemte Mennesker, der have digtet „Folkepoesien“,
og ligesom „Samfundet“ til enhver Tid bestaar af individuelle
Mennesker. Forsaavidt maa al Psykologi være individualistisk,
selv om den med Iver anvender den sociologiske
Metode (ved Siden af andre psykologiske Metoder), hvad der
ganske særligt i Religionspsykologien er Anledning til at gøre.
--- I Slutningen af sin Anmeldelse udtaler Tönnies iøvrigt, efterat
have gentaget, at han i Religionen især ser en Genstand for
Socialpsykologi: „Ich bekenne meinerseits, dass mich Høffding
gelehrt hat, die Tragweite der individuellen Entscheidung besser
als bisher zu würdigen“. ---

Kun naar man lægger Vægt paa det subjektive Element i
Religionen, bliver det muligt at finde en indre Kontinuitet mellem
de skiftende Religioner og Livsanskuelser. Og paa dette
Element beror ogsaa den Betydning for Slægtens fortsatte Liv,
en Religion kan have, selv om dens Tid engang er forbi. Det
er de store menneskelige Skatte, den religiøse Aand har formaaet
at hæve i Menneskenes Indre, der udgør den kosteligste
Arv, Fortidens Religioner efterlade. Et Digt som \textsc{Bunyan’s}
„Pilgrims Progress“, der er blevet kaldet Protestantismens divina
comedia, har saaledes faaet en blivende Betydning (idetmindste
for den engelske Race) ud over den bestemte Konfessions Grænser,
indenfor hvilken det blev til. Og Grunden er den, at det
% --- p. 10 ---
\opage{10}udtrykker, hvad der bevægede sig i et af Livets Modsætninger
dybt grebet Menneskes Sjæl. En udmærket Literaturhistoriker
har med Rette sagt om det: „Alt hvad der er bedst og mest
karakteristisk i Bunyans Skrifter, fremgaar af det indre Drama,
i hvilket de Optrædende vare Gud, Satan og en ensom Menneskesjæl.
Institutioner, Kirker, Ordinanser, Ritus og Ceremonier
hjælpe ham kun lidet eller slet ikke . . . . Lad Enhver med
Troskab fastholde sine inderligste Erfaringer i Livets store Anliggender,
og hans Ord ville vække mangefold Genlyd i andre
Sjæle; \textit{det dybeste Samfund findes ikke i Institutioner eller Korporationer
eller Kirker, men i det ensomme Hjertes Hemmelighed}“\textsuperscript{4}). ---
I min Religionsfilosofi har jeg ogsaa peget hen paa, hvorledes
Sansen for de rent subjektive Elementers Betydning paa det religiøse
Omraade kan bane Vejen for en. dybere Fællesskabsfølelse
end den, der betinges ved Sammenslutning om fælles
Dogmer. Til et og samme Dogme eller til en og samme Kultus
kan der svare højst forskellige personlige Stemninger og Tilstande.
Denne Forskelligartethed er et Bevis paa det subjektive
Elements Selvstændighed. Forholdet mellem de objektive og de
subjektive Elementer (eller maaske bedre: mellem de sociale og
de individuelle Elementer) i Religionen er stedse varierende,
saalænge der er Liv i Religionen. I de historisk udformede Religioner
er det de objektive, sociale Elementer, der stærkest falde
i Øjnene, og hvis Udtrykket „positiv Religion“ er at forstaa i
Analogi med Udtryk som „positiv Ret“, menes derved en Religion,
der optræder som Institution, med en udformet Kultus
og en fastslaaet Overlevering. Psykologisk set maa det dog
altid være det Væsentlige, hvad der motiverer Menneskets Tilslutning
til og Fastholden af det saaledes Fastslaaede, og hvad
der under visse Forhold fører til at foretage Ændringer deri.
Dette subjektive Elements Betydning kan stige saaledes, at det
objektive Element saa godt som elimineres. Vi komme da til
en Grænse, hvor der kan være Strid om, hvorvidt Begrebet Religion
endnu kan anvendes. Det kan let blive en Strid om Ord.
Det kan føre til at opkaste dybsindige Spørgsmaal i Lighed med
det, en Mand stillede: „Mon nu den Stjerne, Astronomerne kalde
Uranus, ogsaa virkeligt er Uranus?“ --- Dersom man lægger
% --- p. 11 ---
\opage{11}Vægten paa Grundmotiverne, ikke paa de ydre Kendemærker,
vil det være berettiget at udstrække Begrebet Religion ud over
de saakaldte positive Religioners Sfære. Religion bliver da et
Slægtsbegreb, indenfor hvilket Religion som positiv Religion,
bygget paa Overlevering og forsynet med faste Kultusformer,
kun betegner en særegen Art. Derved undervurderes Modsætningerne
ikke; ti indenfor et og samme Slægtsbegreb kunne de
stærkeste Modsætninger fremtræde, ligesom der i Virkeligheden
kan føres en Kamp paa Liv og Død mellem Arter, der høre til
samme Slægt. ---

Efter disse indledende Bemærkninger gaar jeg nu over til
at omtale nogle religionsfilosofiske Arbejder fra den nyeste Tid
for derigennem at give et kritisk Supplement til mit eget Forsøg.
Jeg vil vælge saadanne Arbejder, der baade kunne siges
at være typiske for vor Tids Religionsfilosofi og at være lærerige
for dem, der interessere sig for det religiøse Problem. ---
Der er dog nogle Forfattere, som det i og for sig kunde være
naturligt at fremdrage her, men som jeg hellere vil omtale ved
en anden Lejlighed, fordi de bedst ville forstaas i Sammenhæng
med den nyeste Tids Filosofi idethele. Dertil høre især Guyau
og Nietzsche, som jeg om ikke lang Tid vil faa Lejlighed til at
omtale i et Skrift om moderne Filosofer, jeg har under Arbejde,
og til hvilket nærværende Afhandling ogsaa kan betragtes som
Supplement.

I Afhandlingens første Afsnit behandler jeg nogle systematiske
Forsøg; i dens andet Afsnit fremdrager jeg en Række
speciellere religionspsykologiske Undersøgelser, især fra den amerikanske
Literatur.
% --- p. 12 ---
\kapitel{I.}{Systematiske Forsøg.}{I. Systematiske Forsøg}\opage{12}
\afsnit{1. Hermann Siebeck.}{1. Hermann Siebeck}

\textsc{Siebeck} (f. 1842), der virker som Professor i Filosofi i
Giessen, har gjort sig bekendt ved et udmærket Arbejde over
Psykologiens Historie i Oldtid og Middelalder og ved sine æstetiske
Studier. Hans Hovedværk er den i Aaret 1893 udkomne
Religionsfilosofi. Oprindeligt udgaaet fra Herbarts Skole har
Siebeck senere modtaget Paavirkninger af Lotze; men især mærkes
Virkningen af ivrige Kantstudier i hans Hovedværk.

Religionsfilosofiens Opgave er for Siebeck stillet ved Forholdet
mellem Religion og Kultur. Den moderne Kultur er tilbøjelig
til at betragte sig selv som Formaal; Religionen kan
den, naar det kommer højt, kun anerkende som en af sine ufuldkomne
Forgængere. Fra religiøs Side paastaas derimod, at
Kulturlivet --- ligesom Livet og Arbejdet i Erfaringens Verden
idethele ---
kun kan være et Middel for et højere, „oververdsligt“
Formaal. Den Strid, der saaledes opstaar, er det Religionsfilosofiens
Sag at afgøre. Den er altsaa ikke en religiøst
stemt Filosofi, som Nyplatonismen i Oldtiden eller som Schellings
senere Lære; den er heller ikke, som Middelalderens Skolastik,
en Fremstilling af en vis bestemt Religion i filosofisk Form.
Den betragter Religionen i dens forskellige Skikkelser som en
Kendsgerning i det menneskelige Liv og undersøger, hvilken
Betydning denne Kendsgerning har for det menneskelige Livs
Udvikling. Den søger at bestemme, paa hvilket Omraade af
menneskeligt Aandsliv Religionen særligt hører hjemme, og hvil%
% --- p. 13 ---
\opage{13}ken Indflydelse den ud derfra øver paa andre Sider af Livet.
Filosofien undersøger her Religionen paa lignende Maade, som
den i andre af sine Discipliner undersøger Videnskaben, Kunsten,
Retten og Moralen.

Det maa ifølge Siebeck væsentligt være en indirekte Metode,
Religionsfilosofien kan anvende. Den maa prøve, om selve
Kulturlivets Bestaaen og Udvikling ikke blive værdiløse og selvmodsigende,
naar det afskærer sig fra de supplerende Synspunkter
og Motiver, Religionen yder. En positiv Bestemmelse af
den højere Verden, til hvilken Religionen vil aabne Adgang,
kan Filosofien ikke give. Længere end til at paavise Nødvendigheden
af en Supplering kan den ikke naa; men den kan prøve,
om de Former, i hvilke den religiøse Bevidsthed har søgt at udtrykke
det Guddommelige, egne sig til at yde den Supplering,
en Kritik af Kulturudviklingen fører til at fordre. Paa Forhaand
udtaler Siebeck som sin Overbevisning, at der i den angivne
Betydning af Ordet kan gives en Paavisning af „Normaliteten
af den religiøse Bevidstheds væsentlige Indhold.“

Det vil for Religionsfilosofien, naar dens Opgave bestemmes
paa denne Maade, være naturligt at holde sig til de højeste
Former for Religion, da kun de blive af Betydning for Forholdet
mellem Religion og Kultur i det moderne Aandsliv. Siebeck
skelner mellem tre Hovedformer for Religion: Naturreligion,
etisk Religion („Moralitätsreligion“) og Forløsningsreligion.
I den første Form opfattes Guderne som Naturmagter,
og Forholdet til dem bestemmes væsentligt ved Selvopholdelsestrangen.
I den anden Form have Guderne et Kald at udøve,
og deres Dyrkelse motiveres ved den Værdighed, dette Kald
giver dem. Den tredie Form, Forløsningsreligionen, der fremtræder
typisk i Buddhismen og i Kristendommen, er karakteriseret
ved den Betydning, profetiske Personligheder faa. Myte og
Kultus, som raade i de tidligere Former, træde her tilbage for
de oprindelige Personligheder. Der er i al selvstændig Personlighed,
især i den geniale Personlighed, noget Irrationelt, noget
Nyt, der ikke finder sin fulde Forklaring ved, hvad der gaar
forud, og som maaske kun langsomt finder Anerkendelse og Tilslutning.
Men allermest gælder dette om de store Personligheder,
% --- p. 14 ---
\opage{14}fra hvem de nye Værdibegreber udgaa, hvis Gyldighed gennem
lange Tider skal prøves i Slægtens Livserfaringer. Paa dem
passer \textit{Begrebet Aabenbaring.} Dette Begreb er, ifølge Siebeck,
selv et Værdibegreb, d. v. s. et Begreb, hvis Gyldighed ikke
beror paa den Maade, paa hvilken dets Indhold kan paavises i
Virkeligheden, men paa Indholdets praktiske Betydning for Livet.
Dette ses ved at sammenholde forskellige Udtalelser af
Siebeck. Et Sted (p. 192) hedder det, at det religiøse Indhold,
der kundgøres af de store Personligheder, kaldes Aabenbaring
formedelst dets „oververdslige“ Karakter. Og hvad Begrebet
„det Oververdslige“ angaar, hedder det et tidligere Sted (p. 18):
„Begrebet det Oververdslige kan først udspringe paa Grundlag
af en Verdensanskuelse, der ikke blot tager Verden med dens
Tilværelse og Indhold som given Kendsgerning, men i Kraft af
en af det Givne uafhængig, spontan aandelig Virken underkaster
denne Kendsgerning en \textit{Vurdering,} hvis Princip er hentet fra
selve Aandens Indre.“

Det er egentligt først indenfor denne tredie Form for Religion,
at Religion og Kultur træde i Modsætning til hinanden.
Forløsningen, der i Buddhismen forkyndes paa mere negativ, i
Kristendommen paa mere positiv Maade, bestaar i Personlighedens
Myndighed overfor Verden; den er given med Tilegnelsen
af den Tanke, at Verden kun har Værdi paa Grund af det personlige
Livs Existens. Personligheden er paa én Gang fremgaaet
af Verden og selvstændig overfor Verden; den vil ikke
blot forstaa Verden og hævde sig i den og overfor den, men
ogsaa arbejde i den for det, der fra det personlige Livs Standpunkt
staar som det Højeste. Deraf kommer det, at det falder
den religiøse Tro saa naturligt at opfatte selve Verdensgrunden
i Analogi med det Personlige. For Troen er Personligheden det
første Objekt, medens den er Videnskabens allersidste Objekt,
fordi den er den mest sammensatte og tillige mest centrale Genstand,
Erkendelsen kan beskæftige sig med.

Paa dette højeste Standpunkt fremkommer der en Modsætning
mellem Tro og Viden. I Kraft af visse Fænomeners Værdi
er den religiøse Bevidsthed tilbøjelig til at stille dem udenfor
den naturlige Aarsagssammenhæng, medens paa den anden Side
% --- p. 15 ---
\opage{15}Videnskaben er tilbøjelig til at erklære sig selv for at være det
tilstrækkelige Grundlag for Livsanskuelse og Livsførelse. Siebeck
hævder overfor den første Paastand, at Undertro ikke er et nødvendigt
Element i den religiøse Tro. Jo klarere den religiøse
Bevidsthed bliver over sig selv, des mere vil den --- mener han
--- indse, at Undertroen ikke forherliger, men netop begrænser Gud,
idet den sætter Naturen som en Skranke, der skal gennembrydes
af Gud. Overfor Videnskabens Paastand om at være tilstrækkeligt
Grundlag for Livet hævder han, at selve denne Paastand
kun kan være en Trossag, da et videnskabeligt Bevis for den
ikke er muligt, og da Afvisningen af Religionens Nødvendighed
sker ved en Personlighedsakt, der motiveres ved den Forudsætning,
at med selve Videnskaben er det højeste Gode givet. Selve
Trangen til videnskabelig Forsken hviler paa denne ubeviselige
Forudsætning\textsuperscript{5}).

Der er i al Religion et subjektivt og et objektivt Element.
Det subjektive Element, som man kunde kalde Religiøsiteten,
ytrer sig i Stemning, det objektive Element, Religion i Modsætning
til blot Religiøsitet, er Sag, Værk, Ydelse, synlig og
haandgribelig Skikkelse for den subjektive Stemning. Religiøsitetens
Dybde er Værdimaaler for det objektive Religionsvæsens
Tilstrækkelighed og Livskraft; men den subjektive Stemning
vilde aldrig have udviklet sig uden under de objektive Elementers
Indflydelse, saa at Religiøsitet aldrig i Tiden gaar forud
for Religion. Der er et stadigt Vexelvirkningsforhold mellem
dem. Men jo mere et Kultussystem organiseres, des mere udelukkes
nye Stemninger og nye Tanker fra det religiøse Liv.
Det er Kultus, der stedse gør den største Modstand mod Religionsfornyelse.
Jo højere et Udviklingstrin Religionen naar, des
mere udprægede individuelle Variationer vil det religiøse Liv
frembyde, saa at de store menneskelige Forskelligheder ikke
blot kunne gøre sig gældende i Forholdet mellem Religionerne,
men ogsaa indenfor en og samme Religion. ---

Dette er det væsentligste Indhold af den første Afdeling af
Siebecks Værk, den beskrivende Del, som afhandler Religionens
Væsen og Udvikling. Idet han nu gaar over til at drøfte Religionens
Sandhed, udtaler han, at det her ikke drejer sig om
% --- p. 16 ---
\opage{16}nogen enkelt Religions Sandhed, men om Gyldigheden af det
religiøse Princip overhovedet. Denne Gyldighed kan kun begrundes
ved, at Aabenbaringen --- i den ovenfor angivne Betydning
af Ordet --- griber ind i den menneskelige Kultur og, skønt
den ikke selv udspringer af denne, dog yder et nødvendigt Supplement
til den.

Paa dette Punkt er det nu ikke ganske let at forstaa Siebecks
Tankegang. I sin beskrivende Del (som iøvrigt ikke er
blot beskrivende) har han bestemt skelnet mellem Vurdering
og Aarsagsforklaring, saaledes at disse ikke behøve at træde i
Modsætning til hinanden, men betegne forskellige Synsmaader
for Tilværelsen. Men nu opstiller han den Paastand, at en
etisk Handling, en Handling, der gaar ud paa at hævde og gennemføre
det Værdifulde i Virkelighedens Verden, kun bliver
mulig, dersom denne i sidste Instans er underlagt en etisk
Magt, og han finder heri en Begrundelse for Gudsbegrebet. Enhver
etisk Handling hviler, mener han, paa den Forudsætning,
at det er muligt at virkeliggøre de etiske Ideers Indhold indenfor
Natursammenhængen, og denne Forudsætning er en religiøs
Forudsætning. Om den etiske Magt, der her forudsættes, kan
man kun ad Analogiens Vej gøre sig en Forestilling, men en
saadan Brug af Analogien anser Siebeck for berettiget, naar den
undgaar baade „Deismens“ dualistiske Forhold mellem Gud og
Verden og „Panteismens“ Ophævelse af enhver Forskel mellem
dem; derved naar man til en „Teisme“. --- Der savnes her hos
Siebeck en nærmere Undersøgelse baade af Analogiens Betydning
og Berettigelse og af Forholdet mellem Vurdering og Aarsagsforklaring.
Hvad det første Punkt angaar, er dette saa meget
des mærkeligere, som Siebeck selv tidligere i sin interessante
Afhandling \textit{Die metaphysischen Systeme in ihrem gemeinsamen Verhältnisse}
\textit{zur Erfahrung} (i andet Bind af Vierteljahrsschrift für
wissenschaftliche Philosophie) har paavist Analogiens Betydning
for enhver Spekulation, der gaar ud over Erfaringen. Det stod
nu blot tilbage at anvende denne Betragtning paa de religiøse
Forestillinger. Hvad det andet Punkt angaar, fremhæver han
vel Vanskelighederne ved at forene den teleologiske og den kau%
% --- p. 17 ---
\opage{17}sale Betragtningsmaade i fuld indbyrdes Harmoni, men han
mener, at den religiøse Bevidsthed ikke betragter Verden som
en afsluttet Totalitet, men som Del af en højere Tingenes Sammenhæng,
fra hvilken stedse nye skabende Kræfter kunne
strømme ind i den. Jeg er ikke sikker paa, at den religiøse
Bevidsthed her tydes rigtigt; den har i hvert Tilfælde meget
ofte en Tendens til netop at betragte „Verden“ som et Hele for
sig, og Siebeck har selv fremhævet dette i sin Kritik af Mirakelforestillingen.
Det er snarere for den videnskabelige Bevidsthed,
at Verden ikke danner nogen afsluttet Totalitet. „Verden“
er ikke noget videnskabeligt Begreb.

Idet Siebeck saaledes anerkender Berettigelsen af at tænke
sig den naturlige Aarsagssammenhæng underordnet et etisk Princip,
overser han ikke de Indvendinger, der stadigt rejse sig mod
denne Betragtningsmaade ved Existensen af Smerte og Ondskab
i en Tilværelse, der skulde have sin Grund i et saadant Princip.
Han ser ganske vist ikke, at denne gamle Indvending
særligt skarpt maa ramme „Teismen“, der tænker sig hint Princip
i personlig Form. Men han indrømmer, at Vanskeligheden
er uovervindelig for den religiøse Bevidsthed, saa meget end
denne trøster sig med, at det Onde maatte være muligt i en
Verden, hvor frie Væsener skulde udvikle sig. Han har givet
det Kapitel i sin Bog, i hvilket dette Spørgsmaal drøftes, Overskriften
„Theodice“ og bemærker, at teoretisk er det i et Forsøg
paa Theodice Gud, der skal forsvares, men at hvad det for
den religiøse Bevidsthed gælder om, er at holde sit Haab oppe
trods Anfægtelserne. Man kunde hertil bemærke, at forud for
al „Theodice“ gaar et Stadium, hvor det er Mennesket, der skal
retfærdiggøres for Guds Domstol. Saa kommer først senere et
Stadium, hvor Gud stedes for Mennesketankens Domstol; da
begynder man at skrive Theodiceer. Derpaa kommer det tredie
Stadium, hvor ingen af Delene sker, men hvor det tydeligt og
klart gælder en Kamp for at holde Haabet oppe, og hvor de
Forestillinger, man danner sig om, hvad der ligger ud over Erfaringen,
udtrykkeligt staa som Symboler for dette Haab.

Men naar det saaledes viser sig, at de Problemer, Viden%
% --- p. 18 ---
\opage{18}skaben ikke løser, heller ikke løses i Religionen, ja netop fremtræde
i denne i den skarpeste Form, saa er den Supplering af
Kulturen, den Afslutning af Verdensanskuelsen, den religiøse
Ide skulde bringe, ikke lykkedes. Nødvendigheden af at optage
dette Supplement kan ikke paavises.

\whiteline

\afsnit{2. Auguste Sabatier.}{2. Auguste Sabatier}

Det er et Tidens Tegn, at man ogsaa fra teologisk Side begynder
paany at interessere sig for Religionsfilosofi. Denne
Interesse har ellers været saa godt som forsvunden, efterat de
Forhaabninger, man havde næret til den romantiske Spekulation,
vare bristede. Den for faa Aar siden afdøde protestantisk-teologiske
Professor i Paris \textsc{Auguste Sabatier} har i sit Værk
\textit{Esquisse d’une philosophie de la religion d’aprés la psychologie et}
\textit{l’histoire} (5 éd. Paris 1898) sat sig til Opgave at behandle det
religiøse Problem paa psykologisk og historisk Grundlag.

I psykologisk Henseende lægger Sabatier Vægten paa, at
hvad der tilstræbes i Religionen ikke er en Løsning af „Verdensgaaden“,
men af „Livsgaaden“. Det er det Gaadefulde og Modsigende
i vort eget Væsen, der afføder den religiøse Trang, og
Verdensforholdene interessere paa det religiøse Standpunkt kun,
forsaavidt de antages at have Indflydelse paa Løsningen af det
indre Livs Gaader. Den religiøse Trang er en Art højere Selvopholdelsestrang,
som vækkes ved Disharmonierne i vort Indre,
især ved Modsætningen mellem Afhængighed og Frihed, eller
mellem Passivitet og Aktivitet. Vort Livs Strøm, der udspringer
fra det Inderste i os, brydes og splittes som en afmægtig Bølge
mod de ydre Tings Skær og Klipper. Det er Grunden til al
Smerte, og da især til den højeste Smerte af alle, den, der betinges
ved Disharmonien mellem vort Ideal og hvad vi kunne
gennemføre i Virkeligheden. En Harmoni mellem disse Modsætninger
naas ifølge Sabatier i Bønnen, den egentlige religiøse
Akt. Her fremtræder nemlig paa én Gang Afhængighedsfølelsen
overfor den hemmelighedsfulde Magt, vi føle os i Forhold til,
og Frihedsfølelsen formedelst den Kraft, der fra denne Magt
tilstrømmer os. Det Væsentlige i Bønnen er ikke de Ord, den
% --- p. 19 ---
\opage{19}iklædes, eller de Forestillinger, den knyttes til, men den indre
Forbindelse af Hengivelse og Selvvirksomhed, der her gensidigt
støtte hinanden og føre over i hinanden. Religionens Nødvendighed
beror paa, at den udspringer af vore inderste Livsvilkaar.
Derfor vil den vise sig selv der, hvor der er Antipati mod enhver
ydre Religion. Der er et indre Alter, som kan bevares og
holdes i Hævd, selv om man omstyrter alle ydre Altre. Der
gives kun én sand Ateisme: Frivolitet og Blaserethed.

Naar Trangen er stærk og inderlig, er Tilfredsstillelsen allerede
given; den kommer ikke udefra. For Sabatier ligger hele
Religionens Mysterium i Pascals Ord: „Du vilde ikke søge mig,
hvis du ikke havde fundet mig“. Svaret, Aabenbaringen er til
Stede i selve Bønnens Mulighed, i den indre Kraft og Higen,
hvoraf Bønnen udspringer\textsuperscript{6}). I samme Øjeblik Mennesket søger
Gud, er det Gud, der søger Mennesket. Aabenbaringen er ikke
en Meddelelse én Gang for alle af uforanderlige Lærdomme, som
det saa blot gælder at holde fast, men den bestaar i den Maade,
paa hvilken Gudsbevidstheden hos den Enkelte og i Menneskeheden
opstaar og udvikler sig.

Dette Aabenbaringsbegreb betegner Sabatier som det psykologiske
eller det kritiske, i Modsætning dels til det mytologiske,
dels til det dogmatiske Aabenbaringsbegreb. Det har sine store
historiske Repræsentanter i de israelitiske Profeter og i Jesus,
men det finder Anvendelse i enhver religiøs Personlighed. Den
religiøse Inspiration er ligesaa hemmelighedsfuld som den poetiske,
men ikke mere mirakuløs end denne. Den er en indre
befrugtende Kraft, der virker i Harmoni med alle aandelige
Kræfter og Love. Den profetiske Inspiration er kun en højere
Grad af, hvad der rører sig i enhver stærkt bevæget religiøs
Følelsestilstand. Maalestokken for Aabenbaringens Værdi ligger
i den Berigelse, Frigørelse og Forhøjelse af det personlige Liv,
den bevirker.

De Forestillinger og Handlinger, i hvilke Religionen finder
Udtryk, variere med Tiderne. Det er Ortodoxiens, selv Reformatorernes
store Fejl, at de ikke have skelnet mellem den religiøse
Erfaring og dens Udtryk. Man har sammenblandet de
psykologiske Fænomener og de teologiske Forklaringer. Ved
% --- p. 20 ---
\opage{20}ethvert Dogme maa der skelnes mellem den religiøse Erfaring.
der ligger til Grund, og det mere eller mindre ufuldkomne Udtryk,
Reflexionen har kunnet give den. Dette sidste Element
er det mest variable Element i Religionen. Der kan gives tre
Arter af dogmatiske Variationer. Dogmer kunne aftage og forsvinde
af Mangel paa Anvendelse (par désuétude, idet den religiøse Erfaring
ikke længere yder det Indhold, de skulle udtrykke.
Saaledes gaar det med Antagelsen af dæmonisk Besættelse og af
en personlig Djævel. Dogmer kunne dernæst faa et andet Indhold
end det oprindelige, idet de fortolkes efter nye Erfaringer.
En saadan Forskydning foregaar uafbrudt. Ved Inspiration,
Inkarnation og Treenighed forstaa senere Slægter ikke det Samme
som tidligere Slægter. Der heldes stadigt ny Vin paa de gamle
Læderflasker. Endeligt kan der opstaa nye Dogmer som Udtryk
for nye Erfaringer. Saaledes fremkom indenfor Protestantismen
Retfærdiggørelsen ved Tro og det almindelige Præstedømme.

At proklamere Uforanderligheden af Kirkens Dogmer vilde
være at proklamere Uforanderligheden af den Filosofi, der anvendes
til at udforme dem. Dogmatiken maa holde sig til de
religiøse Erfaringer, hvis Udlægning den er. Og naar der anvendes
visse Begreber for at kunne foretage denne Udlægning,
saa maa den ogsaa anerkende de Resultater, Erkendelsesteorien
kommer til angaaende vore Begrebers Betydning og Gyldighed.

Den religiøse Erkendslse er symbolsk Erkendelse. Til at % PRINTED AS IS: Erkendslse
udtrykke Følelserne af Tillid, Frygt og Kærlighed benytter den
Billeder, der ere hentede fra menneskelige Forhold. Symbolets
Magt beror paa den nøje Sammenhæng mellem Følelse og Billede.
Den religiøse Følelse skelner fra først af ikke mellem Indhold
og Form; den er sig ikke bevidst, at den udtrykker sig
symbolsk. Det er den kritiske Filosofi, der indskærper Symbolismens
Uundgaaelighed, som baade Ortodoxien og Rationalismen
hver paa sin Maade miskende. --- Symbolerne kunne blive fælles
for store Kredse, og der er Symboler, som maaske aldrig ville
forsvinde fra det religiøse Liv.

Ikke blot Erkendelsesteorien fordrer Sabatier respekteret af
Dogmatiken, ogsaa Erfaringsvidenskabens almindelige Metode.
Han skelner, som vi have hørt, mellem Verdensgaaden og Livs%
% --- p. 21 ---
\opage{21}gaaden. Verdensgaaden er det Videnskabens Sag at drøfte.
Verden maa tages, som Videnskaben viser os, at den er, og
Livsgaaden stilles saa ved Spørgsmaalet om, hvorledes Videnskabens
Verdensbillede forholder sig til det aandelige Livs inderste
Trang, --- om det indeholder Midler eller Hindringer for Tilfredsstillelsen
af denne Trang.

Erfaringsvidenskabens Selvstændighed maa ogsaa komme
Psykologien til Gode, og her staa vi (som jeg allerede har gjort
opmærksom paa i min Religionsfilosofi p. 173) ved det vanskeligste
Punkt i Sabatiers Religionsfilosofi. Han har ført alle det
religiøse Livs Fænomener tibage til religiøs Erfaring, saaledes % PRINTED AS IS: tibage
som han nærmere beskriver denne, og han har hævdet, at alle religiøse
Forestillinger ere Symboler, som udtrykke denne Erfaring.
Men alligevel antager han, at Gud er Aarsag til den religiøse
Erfaring, og her vilde da foreligge et Tilfælde, hvor et Erfaringsfænomen
fik en teologisk Forklaring, hvad Sabatier selv
udtrykkeligt har afvist Muligheden af, da Gud ikke selv er Led
i Fænomenernes Række og derfor ikke kan bruges til videnskabelig
Forklaring. Tilmed maa jo selve Gudsforestillingen som
alle andre religiøse Forestillinger have en symbolsk Karakter\textsuperscript{7}).

Sabatiers teologiske Udgangspunkt viser sig dog ikke blot
i hans Forsøg paa at forklare den religiøse Erfaring, men ogsaa
i selve hans Beskrivelse af den. Han fremdrager kun en eneste
Type for saadan Erfaring og lægger ikke Vægt paa Variationerne,
skønt han i Principet anerkender deres Mulighed og deres
Berettigelse. Ligesom hos Siebeck føle vi ogsaa her Savnet af
et bredere psykologisk Materiale. Et saadant vilde maaske
have ført til en mindre snever Bestemmelse af Religionens Natur.
Dog hindre disse Indvendinger ikke, at man skilles fra
den franske Teologs filosofiske Arbejde med Erindringen om en
ædel, fin og begejstret Mand.

\whiteline

\afsnit{3. Rudolph Eucken.}{3. Rudolph Eucken}

Efter en Række Forarbejder (især \textit{Die Einheit des Geisteslebens}
1888, og \textit{Der Kampf um einen geistigen Lebensinhalt} 1896)
har denne aandfulde Tænker, der virker i Idealismens gamle
% --- p. 22 ---
\opage{22}Arnested Jena, udtalt sig om det religiøse Problem i sit Skrift
\textit{Der Wahrheitsgehalt der Religion} (1901).

Det er ikke et religionsfilosofisk System, Eucken her har
villet give. Et saadant er efter hans Opfattelse foreløbigt umuligt
paa Grund af de uklare og usikre Forhold paa det religiøse
Omraade: „Under vor Tids aandelige Anarki gives der
ikke noget fast og anerkendt Tilknytningspunkt; al dybere
gaaende Drøftelse maa gaa tilbage til det første Grundlag og
bygge op frany derpaa. Saaledes maatte ogsaa jeg gennem en % PRINTED AS IS: frany
almindelig Undersøgelse af den menneskelige Tilværelse Skridt
for Skridt arbejde mig hen til det Sted, hvor Religionens Problem
bryder frem for saa strax at fremtræde som Midtpunktet
for al Søgen efter vor Tilværelses Sjæl og Betydning“. Forfatteren
erklærer sig selv for en Søgende, og han henvender sig
til Søgende. Han finder i Tiden en stærk Længsel efter Religion
forbunden med en tydelig Bevidsthed om Utilstrækkeligheden
af Religionens nuværende Form. Den nye Livstype, som
Renaissancen hævdede, og det nye Tankearbejde, som derved
indlededes, førte bort fra Kristendommen. Men den storartede
Kultur, i hvilken denne Livstype og dette Tankearbejde udformedes,
har medført indre Forviklinger, og Troen paa dens Tilstrækkelighed
er rokket. Derved kommer da Tanken frem, om
der ikke i Religionen skulde være Noget, der ikke kan og maa
gaa tabt. Der føles Savnet af et stort Maal, som ikke blot
kunde kaste Lys over Menneskelivets store Kreds, men ogsaa
hæve den Enkelte op over hans ofte smaalige Forhold. Og tillige
trænge de gamle Livsgaader, man betragtede som løste
eller dog som fortrængte, sig frem paany. Især føles der en
skarp Modsigelse mellem Menneskets aandelige Anlæg og dets
virkelige Tilstand. Alt dette fører til en Kamp mellem Religion
og Kultur.

Men hverken Religionen alene eller Kulturen alene raader
over Midler til at afgøre Kampen. Afgørelsen maa ske ved, at
der baade indenfor Religionen og indenfor Kulturen skelnes
mellem det Evige og det Forgængelige, og mellem det Uundværlige
og det Uudholdelige. Der maa foretages en Revision
af selve Grundbetingelserne for vort aandelige Liv, og at fore%
% --- p. 23 ---
\opage{23}tage en saadan er Filosofiens Sag. Ikke i den Forstand, at
denne kan konstruere en Religion. Ti Religion er noget Faktisk,
Udtryk for en real Existens, som kan paavises, naar den
er der, men ikke kunstigt frembringes. Men da det her gælder
Paavisning af en aandelig Kendsgerning, er Paavisningen ikke
en ganske simpel Sag; der behøves aandeligt Arbejde for at
finde Religionens ejendommelige Plads i Livets hele Sammenhæng.
Filosofien kan ikke naa videre her end til at paavise Muligheder,
som først gennem store Personligheder og ofte under store
Omvæltninger i de menneskelige Forhold kunne omsættes til
Virkelighed. Men dersom det kan godtgøres, at Religionen
hænger sammen med vor Tilværelses inderste Grund, saa kan
ethvert Angreb paa den og enhver Kritik af den kun tjene til,
at dens inderste, evige Kærne mere og mere lægger sig for
Dagen.

Problemet stilles ved det aandelige Livs Stilling i Tilværelsen.
Dets Bestaaen og dets Udvikling forudsætter, at der bagved
og ud over de aandelige Processer, Erfaringen viser os,
ligger en evig og sammenhængende aandelig Verden (eine geistige
Substanz, ein seelischer Bestand), der indeholder \textit{Mulighederne}
for nye Dannelser og \textit{Udbyttet} af hvad der allerede har
udviklet sig. Aandslivet maa være Mere end et blot Fænomen,
altsaa Mere, end Psykologien kan paavise. Det maa have et
evigt Princip i sig. Ikke blot en psykologisk, men en metapsykisk
Betragtning maa anlægges. For en saadan Betragtning
ere Væsen og Værdi Et, idet dens højeste Genstand er selve
Tilværelsens værdifulde Kærne. Kun ved en saadan Betragtning
bliver et centralt Livsindhold muligt.

Den antike Kultur overvurderede Formen; den ved Renaissancen
grundlagte Kultur overvurderede Kraften. Den sande
Kultur, som har en religiøs Karakter, stiller \textit{Væsensdannelse} (die
Wesensbildung) i Modsætning baade til Formdannelse (die Formbildung)
og til Kraftdannelse (die Kraftbildung). Baade Formen
og Kraften skulle tjene et værdifuldt Indhold. Formdannelsen
alene fører til Stivnen i plastiske Kunstværker, i hvilke det
Evige antages at være udtrykt for alle Tider. Kraftdannelsen
fører til rastløs Bevægelse, saa at al Væren opløses. Væsens%
% --- p. 24 ---
\opage{24}dannelsen derimod hævder en stadig Vexelvirkning af Evighed
og Tid, idet de evige Muligheder gennem Arbejdet i Tiden omsættes
til Virkelighed, og idet paa den anden Side dette Arbejdes
Udbytte bevares i Evighedens Form. Bevarelse af det
Givne og Frembringelse af nyt Indhold maa gaa Haand i Haand,
for at vi kunne trænge frem til den sande Virkelighed.

Den Tankegang, der hævder „Væsensdannelsens“ Mulighed og
Nødvendighed, kalder Eucken den \textit{noologiske} (af noos, Fornuft,
Aand, der hos Nyplatonikerne stilles over psyche, Sjæl), og han
stiller den i Modsætning til den psykologiske Betragtning, som
blot beskæftiger sig med det i Erfaringen givne Bevidsthedsliv.
Det er efter hans Overbevisning ikke noget Nyt, han herved
proklamerer, men en gammel Sandhed, som især \textsc{Kant} har indskærpet
ved sin bestemte Skelnen mellem den logiske, den
etiske og den æstetiske Betragtningsmaade paa den ene Side,
den empirisk-psykologiske paa den anden Side. Her staar Gyldighed
eller Værdi i Modsætning til Virkelighed --- men saaledes,
at Gyldigheden tilsidst hævdes som den højeste Virkelighed.
---
Den noologiske Tankegang er forskellig fra en tidligere
Tids spekulative eller metafysiske Tankegang, som troede at
kunne vinde en Forstaaelse af Fænomenerne ved teoretiske Abstraktioner.
Den noologiske Metode bygger paa Erfaring, paa
Faktum. At forklare noologisk vil sige at indordne en aandelig
Livsytring i det aandelige Livs Helhed, paavise dens Stilling
og dens Opgave indenfor denne Helhed, og derved belyse og
befæste den. Denne Forklaring er forskellig fra den psykologiske,
som blot bestaar i Paavisningen af de Processer, gennem
hvilke Mennesket erhverver og tilegner sig et aandeligt Indhold.
Begge Metoder maa anvendes, forat Modsætningen mellem Idealisme
og Realisme kan overvindes. Idealismen har hævdet
Aandsindholdets Selvstændighed og Egenværdi; Realismens Berettigelse
beror paa dens Fremhæven af de naturlige Betingelser.
Den realistiske Betragtningsmaade kommer til sin Ret, hvor det
indenfor den aandelige Helhed, den noologiske Opfattelse hævder,
lykkes at finde de specielle Aarsager og Love; men den
har Uret i at mene, at Aarsagsforklaringen kan afgive den sidste
Maalestok for Sandhed. Idealismen har med Rette fordret et
% --- p. 25 ---
\opage{25}værdifuldt Indhold for Tilværelsen, men har set bort fra de
reale Betingelser.

Den noologiske Betragtning fører nu ifølge Eucken netop
til den Kendsgerning, Religionen hævder: Nærværelsen af et
absolut Aandsliv, ophøjet over og dog virkende i de i Erfaringen
fremtrædende Livsfænomener. Der gives ingen Religion uden
levende Nærværelse af en højere Verden i vor empiriske Verden
--- men saaledes at disse to Verdener staa i et irrationelt
Forhold til hinanden. Religion kan bestaa uden Gudetro; det
er den ægte Buddhisme et Exempel paa. Men uden Udblik ud
over Erfaringsverdenen er Religion et tomt Ord. Motivet til
Religion ligger i Modsætningen mellem Væsen og Existensform:
i den blotte Erfaring staar Aandslivet som et Tillæg til det
Materielle, som en Bisag; men dette strider mod de ubetingede
Krav, Aandslivet stiller, og de ophøjede Værdier, det hævder.

Ud fra Religionens Grundpaastand tager Virkeligheden sig
anderledes ud end før. En fra Religionen udgaaende Verdensbetragtning
vil rette sit Øje mod det Blivende (i Modsætning
til det Omskiftende), mod det Frie (i Modsætning til det blot
Naturlige) og mod det Fornuftige (i Modsætning til det Ufornuftige).
Den vil særligt mærke sig de Punkter, hvor en ny
Sfære begynder, saaledes Overgangen fra det Uorganiske til det
Organiske, fra Naturlivet til Sjælelivet, fra Sjælelivet til Aandslivet,
og den vil opfatte denne Trinrække som en Udfoldelse af
det Inderste i Tilværelsen. Men bevises kan en saadan religiøs
Verdensopfattelse ikke. Den har Karakteren af en Overbevisning,
ikke af en Videnskab. Det er Filosofiens Sag at bestemme
Forholdet og Grænsen mellem den videnskabelige og den religiøse
Verdensbetragtning.

Motivet til Religion ligger i en aandelig Trang. Det drejer
sig ikke om Menneskets empiriske Lyksalighed og dens Bestaaen,
men det Afgørende er Aandslivets Helhed og Enhed,
Grundlaget for alle specielle og individuelle Livsytringer. Det
er dets Bestaaen, der staar paa Spil, idet den blotte Erfaring
ikke bekræfter den. Det kommer da an paa, om der kan vækkes
en stærk Trang til at fastholde en Overbevisning paa dette
Punkt. Empirisk set er der Grunde nok til Tvivl. Erfaringen
% --- p. 26 ---
\opage{26}viser os ingenlunde Aandens Verden som en selvstændig Verden,
end mindre som Hovedsagen i Tilværelsen. Og under Kulturens
Udvikling ofres Personlighederne for Samfundet og Fremskridtet.
Indenfor selve Aandslivet røre disharmoniske Tendenser
sig; Del staar mod Del, og Del staar mod Helhed; aandelige
Retninger bekæmpe hverandre, og det Onde viser sig at være
en virkelig Magt.

Overfor de Tvivl, der saaledes rejse sig, fremhæver Eucken
for det Første, at selve den stærke Følelse af Disharmonien er
en Følge af, at der stilles store Fordringer til Livet. Lidelsens
Bitterhed og Modstandens Haardhed vidne om Dybden af vor
Existens og om højere Kræfters Tilstedeværelse i vort Væsen.
Med Modsætningerne vilde ogsaa Smerten falde bort. Det var
aldrig i Tider med ulykkelige Oplevelser og tunge Forviklinger,
at Menneskeheden tvivlede om sin ideale Opgave, men derimod
trivedes denne Tvivl i saadanne Tider, hvor trægt Velbehag
og indbildt Mættelse raadede.

For det Andet er det Euckens Overbevisning, at hin Tvivl
maa bekæmpes ved en kraftigere og mere koncentreret Form
for Religion end den ubestemte og universelle Religion, som i
det Foregaaende er beskreven. Han henpeger til de store Personligheders
Betydning i Religionernes Historie. Hos dem rørte
det Guddommelige sig paa dybere og kraftigere Vis end hos
Andre, og de formaaede at udtrykke Livets mangfoldige Erfaringer
i klare og glødende Billeder. Et nyt Liv brød her igennem
paa afgørende Maade og greb bestemmende ind i Slægtens historiske
Udvikling. Ved at fremdrage deres Betydning gøre vi
Overgang fra \textit{den universale Religion} til \textit{den karakteristiske Religion.}

„Karakteristisk Religion“ betegner hos Eucken de historiske
eller positive Religioner, forsaavidt de --- hver paa sin ejendommelige
Maade --- give et nyt, i sine Grundtræk klart Totalbillede
af Livet og stille dette Billede i Modsætning til Erfaringens
Verden. Den adskiller sig fra „den universale Religion“
ved den bestemte Udformning, uden hvilken en saadan Modsætning
til Erfaringsverdenen ikke kan opretholdes. Den udspringer
af store Personligheder og hæver sig over al Folkekultur som
en umiddelbar Ytring af Aandslivets inderste Væsen. Uden
% --- p. 27 ---
\opage{27}saadan karakteristisk Religion er selvstændig Inderlighed ikke
mulig. Det drejer sig her ikke om en blot Udsmykning af Livet,
men om et nyt Livs, en ny Virkeligheds Gennembrud. Dog
ere de Gennembrud, de store positive Religioner vidne om, kun
Højde- og Vendepunkter indenfor en stor Proces. Det Guddommelige,
som her bryder igennem, er ikke først opstaaet paa
disse bestemte Steder.

Det højere Liv, der saaledes bryder frem, formaa vore Begreber
ikke at udtrykke. Fantasien maa komme os til Hjælp
med sine Symboler. Og her ligger Religionens saarlige Punkt.
Fantasien bestemmes stedse ved det givne Kulturtrin og dets
Overleveringer, og Symbolerne ville derfor faa en historisk Karakter.
Der vil engang komme en Tid, da det, der maaske for
Aartusinder var Udtryk for den højeste Sandhed, og som for
den menneskelige Bevidsthed havde indgaaet en inderlig Forbindelse
med det religiøse Indhold, rokkes i sin Betydning og
maa vige Pladsen for nye Former --- ofte først efter svare Rystelser
og haarde Kampe, under hvilke Tvivlen fra at være rettet
mod Symbolerne kan gaa over til at angribe Religionens inderste
Kærne.

Den universale Religion er stadigt nødvendig som Korrektiv
for eller Indledning til de karakteristiske Religioner, der overladte
til sig selv let indsnevres og blive til blotte Trøstemidler.
De karakteristiske Religioners Betydning beror ikke paa, at
nogen af dem skulde have fundet afsluttende Udtryk for den
højeste Sandhed, men paa deres levendegørende og forhøjende
Virken under Tidens Uro og Omskiftelser. Det maa stedse
paany prøves, om det historisk Overleverede virkeligt udfører et
saadant Arbejde. Men Kirken har en Tendens til at stille sig
slavisk overfor et en Gang for alle givet Forbillede og opfatte
al Sandhed som bestaaende i Efterligning og Gentagelse. Den
isolerer det Højeste og giver det derved en passiv Stilling overfor
det virkelige Liv. Derved bliver Kamp stedse nødvendig
paa det religiøse Omraade. Det er de store religiøse Personligheders
Betydning, at de nøde os til at vælge. ---

Hvad der især bestemmer Euckens Plads i den moderne
Religionsfilosofi, er hans afgjorte Overbevisning om, at Værdi%
% --- p. 28 ---
\opage{28}erne kun kunne hævdes, naar der bagved og ud over Erfaringens
Verden antages en højere Virkelighed, i hvilken det Aandsliv,
der i Erfaringen fremtræder spredt og i skiftende Former,
bestaar i Enhedens og Evighedens Form. Paa en ofte dybsindig
og stemningsfuld Maade hævder han denne Modsætnings
Nødvendighed, --- snart saaledes, at han minder om den gamle
Mystik, snart saaledes, at han minder om den romantiske Idealisme,
og snart saaledes, at det er den personlige Livskamps ophøjede
Poesi, der faar Udtryk i hans Ord. Dog betoner. han % PRINTED AS IS: the stop after betoner
meget stærkt, at Filosofien her kun kan paavise Muligheder, og
at det er gennem profetiske Personligheder, saadanne Muligheder
vinde Virkelighed og bestemt Form. Naar han --- for
at undgaa Subjektivismen --- hævder Nødvendigheden af en
Metafysik (eller en „Metapsyki“), saa er det dog egentligt i
Kraft af det Sandes, Skønnes og Godes \textit{Værdi,} at denne Metafysik
skal være nødvendig. Hvad han kalder den „noologiske“
Metode, adskiller sig fra den psykologiske Metode derved, at
den vil paavise, at hin Værdi kun kan hævdes ved Antagelse
af en Helhedssammenhæng, der bærer de i vor Erfaring spredt
og flygtigt optrædende værdifulde Livsytringer. Især lægger
den noologiske Betragtningsmaade Vægt paa de store Gennembrud,
de betydningsfulde Begyndelser og de kvalitative Forskelligheder
indenfor Aandslivet, som vi kende dette i Erfaringen.
Han anser Erfaringsvidenskabens Forsøg paa at paavise
Kontinuitet i Fænomenernes Verden for haabløse, men fortolker
Bruddene paa den empiriske Kontinuitet som enkelte Udslag af
en stor transscendent Kontinuitet, et Mulighedernes, Endemaalenes
og Værdiernes Rige, der kun glimtvis aabenbarer sig i
Endelighedens Verden. Derved er det, han er beslægtet med
den romantiske Idealisme. Men Noologien adskiller sig fra den
spekulative Metode derved, at Antagelsen af en saadan Helhedssammenhæng
udtrykkeligt kun grundes paa Erfaringen om det
Værdifulde. Det synes da klart, at Udgangspunktet er empirisk
og personligt. Han faar Fart til at gaa over fra Erfaringen til
„Metapsykien“ gennem den Spænding, som frembringes ved et
stort Dilemma. Enten --- paastaar han --- maa man betragte
Aandslivet som et flygtigt Fænomen, et Vedhæng ved den ma%
% --- p. 29 ---
\opage{29}terielle Natur, --- eller man maa antage, at det i sig selv udgør
en stor Totalitet og er Et med Tilværelsens inderste Kærne.
Dette Dilemma kommer frem hos ham, baade hvor han udvikler
Ideen om „den universale Religion“, altsaa vil gøre Overgangen
fra det empiriske Standpunkt til et religiøst Standpunkt af mere
ubestemt Art, og hvor han gaar over til den „karakteristiske Religion“,
altsaa vil paavise Nødvendigheden af en historisk udformet
Religion. Karakteristiske ere især følgende Udtalelser. „Ved
denne Afgørelse staar ikke dette eller hint, men Aandslivet som
Helhed paa Spil. I Menneskets Erfaring fremtræder Aandslivet i
en uholdbar Mellemstilling; fra denne maa Mennesket enten
vende tilbage til den blotte Natur og bekende, at al ejendommelig
Stræben kun var en tung Drøm, eller modigt gaa fremad
og sikre den nye Stræben en ny Verden. Kun en uklar Tænkemaade
kan søge en Mellemvej mellem disse Veje, og kun et
mat Sindelag kan roligt udholde Dilemmaet. I hvert Tilfælde
staar det fast, at uden Religion gives der ingen Sandhed i
Aandslivet og ingen indre Storhed i Mennesket.“ (Der Wahrheitsgehalt
der Religion. p. 238.) „Er det menneskelige Liv et
blot Tillæg til Naturen, eller er det Begyndelsen til en ny
Verden? Af dette Spørgsmaals Besvarelse afhænger al Udformning
af vore Forhold, hele Retningen af vor Handlen.
Religionen vil hæve den menneskelige Tilværelse op paa en over
Verden hævet Højde og derved redde vort Liv fra den Intethed,
som det ellers sikkert forfalder til. Er dette dens Forsøg
kun en Ikarusflugt, saa falder alt Haab bort, saa staar det
Ædleste og Bedste som tom Indbildning, saa ender det Hele i
Ufornuft.“ (ibid. p. 319). ---

Hvad der skal drive os til at vælge her, det er dog stedse
tilsidst den personlige Livserfaring, særligt den Grad, i hvilken
man betages af Værdiernes Skæbne i Virkeligheden. Der synes
ikke at være nogen Grund til at indføre en ny Terminus;
„Noologi“ er hverken nogen nødvendig eller nogen heldig Berigelse
af den allerede i Forvejen tilstrækkeligt forsynede filosofiske
Terminologi. Af en privat Udtalelse af Eucken, som
jeg haaber at kunne anføre uden at begaa en Indiskretion, ses
det klart, at det er den personlige Erfaring om Livets Mod%
% --- p. 30 ---
\opage{30}sætninger, særligt af den store Modsætning mellem Ideal og
Virkelighed, der har bestemt hans eget Valg og ført ham til
„Noologien“. I et Brev, han skrev til mig efter Læsningen af
min Religionsfilosofi, udtaler han, efter at have dvælet ved den
overvejende psykologiske Karakter af mit Forsøg paa at behandle
Problemet: „Mich treibt es, bei allem Gegensatz zur alten ontologischen
Metaphysik, immer wieder zur Metaphysik und so
auch zu einer metaphysischen Begründung der Religion; ich
fürchte ohne das dem Subjectivismus zu verfallen, und dagegen
sträube ich mich namentlich deswegen, weil ich stärker unter
dem Eindruck der grossen, in der Erfahrung unlösbaren Konflikte
unseres Daseins stehe, kurz gesagt, dualistischer denke
und denken muss.“ I denne Udtalelse peger Eucken direkte
hen til individuel Erfaring som bestemt ved den enkelte Personligheds
Ejendommelighed. Det er den, der bestemmer hans
Standpunkt. Men jo klarere dette Grundlag fremtræder, des
naturligere bliver det dog at anvise Religionsfilosofien den Opgave
at undersøge den psykologiske Jordbund, af hvilken de
forskellige religiøse Standpunkter voxe frem, og dertil knytte
en erkendelsesteoretisk og en etisk Prøvelse af de Tanker og
de Bestræbelser, som disse Standpunkter give sig Udslag i. En
ganske særegen Metode behøver da Religionsfilosofien ikke, og
den saakaldte „Noologi“ er jo ogsaa selv en „Mellemvej“ mellem
Spekulation og praktisk Tro. Som Spekulation henfalder
den til Erkendelsesteoriens, som praktisk Tro til Psykologiens
og Etikens Behandling. Den Opfattelse af Religionsfilosofiens
Opgave og Fremgangsmaade, som heraf følger, synes at være
mest i den kritiske Filosofis Aand, i hvilken jo Eucken --- ligesom
Siebeck og Sabatier --- vil arbejde. Den „metafysiske“ Begrundelse
af Religionen, som Eucken kræver af Filosofien, viser
sig efter hans egne Udtalelser tilsidst at støtte sig til Trangen
til Afgørelse af et stort praktisk Dilemma.

Jeg finder i Euckens Standpunkt --- i den „Noologi“, der
fører ham til hans religiøse Metafysik --- et Exempel paa, hvorledes
personlig Livserfaring kan føre til Tro paa Værdiens Bestaaen.
For mig er enhver Tro, selv om den fremtræder i de
mest filosofiske Former, stedse kun \textit{Genstand} for Filosofi, ikke
% --- p. 31 ---
\opage{31}selv \textit{Produkt} af Filosofi. I hvert Tilfælde er den Filosofi, der
producerer den, selv Produkt af personlig Livserfaring. Filosofien
som Videnskab raader nu engang ikke over andre Metoder
end den logisk sluttende, den psykologisk forklarende og den
etisk vurderende.

Men fra mit Standpunkt faa saadanne Arbejder som Euckens,
og det Samme gælder Sabatiers, en overordentligt stor Interesse
ved den klare Maade, paa hvilken de afgørende Motiver for
religiøs Livsanskuelse træde frem. To idealt sindede, af Livet
grebne og med Tankens Energi udrustede Mænd have her udtalt
sig om et af de største menneskelige Problemer.

% --- p. 32 ---
% CHECK p. 32: centred line: heading or verse
\kapitel{II.}{Religionspsykologiske Undersøgelser.}{II. Religionspsykologiske Undersøgelser}\opage{32}

Et gennemført Forsøg paa at gaa tilbage til umiddelbar Erfaring
ved Behandlingen af det religiøse Problem maa særligt
henvende Opmærksomheden paa det Punkt, hvor Bevidstheden
fra umiddelbar Oplevelse gaar over til at forme en Dom, altsaa
paa den første Opstilling eller uvilkaarlige Tilblivelse af en religiøs
Sætning. Hvad der her gaar for sig, kan ikke være væsensforskelligt
fra, hvad der altid gaar for sig ved Overgangen
fra Iagttagelse til Dømmen. Den religiøse Bevidsthed har hverken
en Psykologi eller en Logik, der er forskellig fra den almindelige
Psykologi og Logik, og de religiøse Domme maa blive
til, som alle andre Domme blive til. Hvad der er ejendommeligt
for de religiøse Domme, er, at den Erfaring, de tilsidst grunde
sig paa, er en Følelseserfaring, en Erfaring om visse Oplevelsers
Værdi som fremhjælpende eller hæmmende, og at de Oplevelser,
der her er Tale om, angaa Forholdet mellem Ideal og Virkelighed,
mellem Værdiernes og Realitetens Verden. Religiøse Værdier
ere derfor Værdier i anden Potens eller sekundære Værdier:
de forudsætte, at Mennesket kender Noget, der giver Livet Værdi
(hvori dette saa bestaar), og at han saa erfarer, at der er Magter
i Verden, som fremme eller hæmme dette Værdifulde. Paa den
religiøse Følelses Højdepunkter falder naturligvis Forskellen
mellem primær og sekundær Værdi bort; der kendes intet Hvorfor
og anvendes ingen Prædikater. I mest udpræget Form sker
dette i Mystiken. Men hvor der kan finde en Dømmen Sted,
altsaa være Tale om Udsagn af religiøs Bevidsthed, der vil ---
% --- p. 33 ---
\opage{33}paastaar jeg --- en nærmere Analyse godtgøre, at hvad der udsiges
Noget om, er Forholdet mellem Værdi og Virkelighed,
saaledes som Individet mener umiddelbart at have erfaret dette.
Religiøse Domme angaa Værdiernes Skæbne i Virkelighedens
Verden. --- Herved maa det dog vel bemærkes, at religiøse Domme
ikke altid, maaske endog saare sjeldent, dannes paa første Haand.

Ved enhver Domsdannelse maa der (som jeg nærmere har
udviklet i min Afhandling \textit{Det psykologiske Grundlag for logiske
Domme)} skelnes mellem Udgangsforestillingen og Slutningsforestillingen.
Ved den sproglige Formulering af Domme optræde
disse to Forestillinger som logisk Subjekt og logisk Prædikat.
Ved Domsprocessens Begyndelse har Udgangsforestillingen en
vis Ubestemthed, som søges hævet ved Opdagelsen af et Prædikat,
en nærmere Bestemmelse, der sker gennem Analyse af
Udgangsforestillingen og af den Sammenhæng, i hvilken den forekommer.
Jo rigere Udgangsforestillingens Indhold er, des vanskeligere
vil det være at finde en definitiv Slutningsforestilling,
et afsluttende og adækvat Prædikat. I denne Henseende frembyder
den religiøse Bevidsthed, saaledes som den fremtræder i
Historien, de største Forskelligheder. Jo snevrere baade Værdiernes
og Virkelighedens Verden er, og jo mere elementært
Forholdet mellem dem fremtræder, des lettere er det for Mennesket
at slaa sig til Ro ved et Prædikat og dermed afslutte
den religiøse Tænkning. Her er alle mulige Grader lige fra
Fetischismens og Animismens begrænsede og barnlige Prædikater
til den højeste Mystiks sublime Grænsebegreber. Og den
Proces, ved hvilken Slutningsforestilling her knyttes til Udgangsforestilling,
er snart uvilkaarlig Ideassociation, snart Efterligning
og Vane, snart spekulativ Grublen, snart bekymret Søgen og
Spørgen. Udgangsforestillingens Indhold er ved religiøs Erfaring
det harmoniske eller disharmoniske Forhold mellem Værdi
og Virkelighed, og Slutningsforestillingen skal angive, hvad det
er, der lægger sig for Dagen i dette Forhold. Dersom en saadan
Slutningsforestilling vindes eller findes, er Overgangen gjort fra
Erfaring til Tro. Gudsbegrebet i dets forskellige Former er det
Grundprædikat, i hvilket den religiøse Bevidsthed søger at faa
Udtryk for og Forklaring af sine dybeste og væsentligste Er%
% --- p. 34 ---
\opage{34}faringer. Ordet Gud skal betyde „den, til hvem der ofres“, altsaa
den, til hvem Mennesket føler sig i Gæld, til hvem det fører
sine Erfaringer om Forholdet mellem Værdi og Virkelighed tilbage.
Men hvem er det? Qvid est quod amo, quum te amo?
spørger Augustinus. Her ligger Mulighed for stedse ny Grublen,
især da Udgangsforestillingen ikke er konstant, eftersom
Forholdet mellem Værdi og Virkelighed stedse ændrer sig.

Der findes i Goethes Faust en Monolog, som egner sig godt
til Exempel paa, hvad her menes med religiøs Erfaring. I Scenen
„Wald und Höhle“ foretager Faust et Opgør. Naturens og
Tankens store Skatte ere aabnede for ham; han kan fordybe sig
i den indre som i den ydre Verden i Forstaaelse og Stemning;
men tillige er der medgivet ham en dæmonisk Ledsager, som
han ikke kan undvære, og hvis Kulde og Frækhed forvandler
hine høje Gaver til et Intet --- samtidigt med, at han vækker
det sanselige Begærs flammende Ild. Til en Slutningsforestilling
kommer det her ikke; vil man søge en saadan i Digtet, maa den
findes dels i Prologen til første Del, dels i Slutningen af anden
Del. --- Den store Modsætning, Goethes Faust her gør Rede for,
ligger --- kun i endnu skarpere Form --- til Grund for Schopenhauers
Livserfaring, og Schopenhauers Filosofi var i sin inderste
Grund et Forsøg paa at finde en Slutningsforestilling som Udtryk
for denne Erfaring. ---

Umiddelbar religiøs Erfaring er et Ideal, da ethvert Menneske,
vi kende, begynder under Indflydelse af en Tradition og
strax tyder sine Oplevelser ved Overleveringens Hjælp, eller
maaske nøjagtigere, strax ser dem i de Former og den Belysning,
Overleveringen har lagt til Rette. Allerede Indlevelsen i
Sproget bevirker, at der her ikke kan være Tale om absolut oprindelig
Erfaring. Enhver religiøs Erfaring er mere eller mindre
konfessionel. Baade Værdierne, Virkeligheden og Forholdet mellem
dem opfatter Mennesket fra først af under Indflydelse af Forestillinger,
som det ikke selv har dannet paa egen Haand, endsige
selv udledet af sine Erfaringer. Det vil da være af stor Interesse
at undersøge, hvorledes de overleverede Forestillinger efterhaanden
ændres ved Individets egne fremskridende Erfaringer, og
dernæst at anstille en Sammenligning mellem religiøse Stand%
% --- p. 35 ---
\opage{35}punkter og Udviklingsprocesser, der bestemmes ved væsentligt
forskellige Udgangspunkter.

Der foreligger nu om alle her berørte Emner --- baade om
religiøse Domme som Værdidomme, om en konfessionel Religiøsitets
Former og Modifikationer og om komparativ Religionspsykologi
--- en Række Arbejder, der for største Delen ere fremkomne
efter Afslutningen af min Religionsfilosofi, og som det
vil have sin Interesse at faa karakteriserede.

\whiteline

\afsnit{1. Om religiøse Domme. {\normalfont\textit{(Max Reischle).}}}{1. Om religiøse Domme. (Max Reischle)}

De Forsøg paa Religionsfilosofi, som vi i det første Kapitel
beskæftigede os med, vare mere eller mindre konsekvent byggede
paa den kritiske Filosofi og dens Sondring mellem rent
teoretiske Domme og Værdidomme. Indenfor Teologien har denne
Distinktion i de senere Aar spillet en ikke ringe Rolle, især hos
\textsc{Albrecht Ritschl} og hans Disciple. Auguste Sabatier staar
denne Retning nær, og han taler da ogsaa om Værdidomme
(jugements d’estimation ou de dignité). Den Vægt, der lægges
paa Forskellen mellem de to Arter af Domme, er karakteristisk
for den nyeste Tids Bestræbelser for at føre Religionen tilbage
til umiddelbar Erfaring og bestemme denne Erfarings egentlige
Beskaffenhed. Det er gennem Undersøgelsen af Værdidommene,
at der i den nyeste Tid er kommet en Berøring mellem Teologi
og Filosofi, som ellers i de sidste halvhundrede Aar --- især paa
Grund af Teologiens mere og mere udpræget konfessionelle Karakter
---
have staaet hinanden fjernt. --- Jeg skal her som
Exempel henlede Opmærksomheden paa en Afhandling af \textsc{Max
Reischle}: \textit{Werturteile und Glaubensurteile.} (Halle 1900).

En Værdidom forudsætter ifølge Reischle en Værdifølelse,
der er forskellig fra den blotte Lystfølelse. Forudsætningen for
en Værdifølelses Opstaaen er, at Tilfredsstillelsen angaar hele
Jeget (das Gesammt-Ich), ikke en enkelt Del eller Tilstand af
Jeget. Værdi tillægger jeg en Ting, om hvilken jeg ad Eftertankens
Vej har indset, at den vil tilfredsstille mit Jeg som Helhed.
Dog indrømmer Reischle, at Grænsen mellem „hele Jeget“ og
dets enkelte Dele eller Tilstande er flydende. Hvad han vil
% --- p. 36 ---
\opage{36}fremhæve, er især det, at der forud for Værdidommen gaar en
sammensat Proces, ved hvilken Værdifølelsen naar en Klarhed
og Sikkerhed, den blotte Lystfølelse mangler.

Værdidomme kunne ifølge Reischle være af to Arter: naturale
og ideale (eller normative).

\textit{Naturale Værdidomme} ere begrundede i Menneskets Natur og
Livsvilkaar, der bestemme den Tilfredsstillelse, der er mulig
for det. Her gøre sig mange Forskelligheder gældende mellem
de forskellige Individer, men da Interesserne dog til en vis Grad
ere solidariske, kan der gives kollektive Værdidomme, og da
Livsvilkaarene til en vis Grad ere fælles, kan der gives generelle
Værdidomme. Hvor en ydre Maalestok (f. Ex. Penge) er mulig,
kan Værdidommen faa en objektiv Karakter, saa at maaske endog
dens subjektive Oprindelse glemmes.

\textit{Ideale eller normative Værdidomme} forudsætte en Ide eller en
Norm, som vi anerkende og vente anerkendt af Andre. Herhen
høre æstetiske, etiske, intellektuelle og religiøse Værdidomme.
Disse bestemmes ikke blot ved den Maade, paa hvilken vi passivt
bestemmes udefra, men væsentligt ved Overensstemmelsen
med en Lov, der gælder for vor egen Aands Virksomhed. Ved
etiske Værdidomme f. Ex. følge vi en Norm eller en Ide, som
formaar at føre vor Villie til Harmoni og Frihed. Ved religiøse
Værdidomme er den tilgrundliggende Ide den Betydning, det
har at underkaste sig en højeste Magt for ved Enhed med den
at naa den højeste Frihed overfor Verden. De ideale Værdidomme
gøre Fordring paa Almengyldighed, fordi Ingen kan fornægte
de Ideer, der betinge dem, uden at opgive sin Personlighed,
sit Væsens Frihed og indre Sammenhæng. ---

Forskellen mellem naturale og ideale Værdidomme er
--- hvad Reischle ikke har bemærket --- aabenbart ligesaa flydende,
som han selv indrømmer, at Forskellen mellem Lystfølelse og
Værdifølelse er. Hint „Jeg som Helhed“, Værdifølelsen skal
forudsætte, indeholder nemlig en Norm, der allerede gælder for
de specielle og skiftende Elementer og Tilstande; disses Værdi
bestemmes ved Forholdet til det. Og desuden maa paa den
anden Side „Jeget som Helhed“ selv godtgøre sin Værdi overfor
de enkelte Elementer og Tilstande, hvis Underordning kræves;
% --- p. 37 ---
\opage{37}det er ingenlunde selvfølgeligt, at Helheden har Ret overfor det
enkelte Element eller den enkelte Tilstand. Endeligt beror jo
ifølge Reischle selv de æstetiske, etiske, intellektuelle og religiøse
Aandsfunktioners Værdi paa, at de skaffe vort Væsen Harmoni
og Frihed; --- men det vil ikke sige Andet, end at de gøre
et „Jeg som Helhed“ muligt, og de ideale Værdidomme blive
fra denne Side set ikke forskellige fra de naturale. Det er derfor
ingenlunde let at opstille Værdimaalere paa det aandelige
Omraade, hvor der er den inderligste Vexelvirkning og gensidige
Bestemmelse af Elementerne indbyrdes og mellem Elementer og
Helhed. Indenfor den individuelle Bevidsthed, ligesom indenfor
Samfundet, kæmpe de forskellige Værdimaalere. Enhver Ides
eller Norms Værdi skal selv godtgøres --- og derfor stilles Problemet
stedse paany. Af alle Værdidomme frembyde netop, som
jeg et andet Sted har vist, de, der føre til Opstilling og Anerkendelse
af en Norm, de største Vanskeligheder for Teorien\textsuperscript{8}).
Overgangen fra det „Naturale“ til det „Normative“ (Ideale) indbefatter
store Problemer, som en blot Klassifikation af de faktisk
givne Værdidomme let tildækker. Reischle har klart Øje
for det analoge Problem, der opstaar, naar de forskellige Religioner,
hver med sine Grundværdier, staa overfor hverandre;
men han har ikke set, at Kampen allerede kan begynde i den Enkeltes
Indre, og at en ledende Ide eller Norm stedse er Resultatet
af en saadan, mere eller mindre heftig Kamp. ---

I den sidste Del af sin Undersøgelse drøfter Reischle, hvorvidt
\textit{Trossætninger} kunne kaldes Værdidomme. Han indrømmer,
at Autoritet og Tradition spille saa stor en Rolle baade i den
protestantiske og i den katolske Kirke, at de Trossætninger, de
Enkelte hylde, ikke kunne kaldes Værdidomme i den Betydning,
at disse Enkelte selv have udledet dem af deres Værdierfaringer.
Derimod anser Reischle det for utvivlsomt, at Trossætningerne
maa kaldes Værdidomme, naar der spørges, hvorpaa
deres Fordring om Gyldighed principielt beror. De ere
Værdidomme af ideal eller normativ Art; deres Betydning bestaar
i, at de angive Betingelser for Saligheden, d. v. s. for
Personlighedens fuldstændige Harmoni og Frihed. Det vilde
være urigtigt at kalde dem Postulater, som om det blot var vor
% --- p. 38 ---
\opage{38}egen Trang og Higen, der i dem formulerede, hvad der kunde
volde Tilfredsstillelse. Der ligger ved de religiøse Domme bestemte
Kendsgerninger til Grund, for den Kristne saaledes Jesu Liv
og Død, og de religiøse Domme udtrykke Tilliden til saadanne
Kendsgerningers Magt og Betydning\textsuperscript{9}). Den Kristne tror paa
Kristus, fordi han kun i Tilslutning til ham kan finde det
sande Liv. Om alle historiske Religioner gælder det, at visse
givne Fænomener --- Personer og Begivenheder --- tydes som
guddommelige Aabenbaringer, og de religiøse Domme udtrykke
Tilliden til disse Aabenbaringer. Paafaldende, især velgørende
eller ødelæggende Fænomener i Naturen, i det bevidste eller
det ubevidste Sjæleliv, i Folket eller i Historien vække den
uvilkaarlige Forestilling om, at højere Magter her virke til
Menneskers Vel eller Ve. Ofte dikteres Dommen af Frygt,
oftest dog af Tillid. Og jo højere en Religion staar, des mere
er det etiske Ideer, der diktere Dommen. Selve det, at Noget
kaldes „Aabenbaring“, er altsaa efter Reischle en Værdidom;
det beror paa, om vedkommende Personligheder eller Begivenheder
staa som Former for store Værdiers Indtræden i Historien,
og om de kunne blive Genstand for tillidsfuld Tilslutning
og Hengivelse. En rent objektiv Afgørelse af Striden mellem
forskellige Aabenbaringer er derfor umulig. Dette er en nødvendig
Konsekvens af Empirismen i Religionsfilosofien, og
Reischle drager ogsaa denne Konsekvens, som han kun fornegter
ved at sætte Ordet „Aabenbaring“ med Anførselstegn,
hvor det ikke drejer sig om hans egen Religion.

\whiteline

\afsnit{2. Den religiøse Følelses Patologi. {\normalfont\textit{(Murisier).}}}{2. Den religiøse Følelses Patologi. (Murisier)}

For at kunne paavise Elementerne i den religiøse Følelse
kan det være hensigtsmæssigt at studere dens extreme Former,
i hvilke den antager en sygelig Karakter. Det Sygelige bestaar
i, at Elementer, der ellers forekomme i harmonisk Sammenhæng
med andre Elementer, udsondre sig og optræde med overvældende
Styrke, saa at de beherske hele Livets Gang. I sit
Skrift \textit{Les maladies du sentiment} religieux (Paris 1901) har
E. \textsc{Murisier} (Professor i Neufchâtel) sat sig den Opgave at
% --- p. 39 ---
\opage{39}undersøge den religiøse Følelses sygelige Former, som han
mener at kunne føre tilbage til to, Extasen og Fanatismen,
hin det extreme Udtryk for den religiøse Følelse i dens individuelle
Form, denne dens extreme Udtryk i social Form. Disse
to Former staa i Modsætning til hinanden: den ene fører til
absolut Isolering, den anden til udadgaaende Virken og til Forbindelse
med Andre. Men der ligger, som Murisier viser, en
og samme Trang til Grund ved begge Typer. Mystikeren ser
i Isolationen og i Bortfjernelsen af al Mangfoldighed fra Bevidstheden
det eneste Middel til at bevare Enheden og Sikkerheden,
til at fastholde Livets totale Retning mod et eneste Maal. Hans
Aand mangler Energi til at magte Livets Mangfoldighed i ham
og udenfor ham, og han samler sig derfor paa et eneste Punkt.
Fanatikeren søger derimod at opnaa Enhed og Sikkerhed i
Livets Gang og Retning ved en praktisk Handlen, som fører
saa Mange som muligt til at dele hans Tro. Han søger at tilvejebringe
en social Enhed, indenfor hvilken ingen disharmoniske
og forstyrrende Tendenser mærkes. Ved social Ensformighed
støttes det egne Jegs Enhed og Identitet. Fanatikeren priser
kirkelig Enhed med ligesaa stor Iver som Mystikeren Tankens
Enhed. Og ligesaa ivrigt som Mystikeren holder alle Indtryk
borte, der kunde forstyrre hans Opgaaen i det Ene, der er det
Højeste for ham, ligesaa ivrigt undertrykker Fanatikeren enhver
Retning, der kunde bryde den sociale Enhed. Hans Forfølgen
svarer til Mystikerens Askese, og hans Proselytisme til Mystikerens
Koncentration. For mange Mennesker er Tilslutningen
til et religiøst Samfund, især et snevert og exklusivt, den eneste
Vej, ad hvilken de kunne finde Ligevægt og Fred i Sindet.

Gennem Studiet af disse modsatte Former for Religiøsitet,
støttet til et historisk, biografisk og medicinsk Materiale, finder
Murisier i de Karaktermærker, der ere fælles for dem, et Bidrag
til Karakteristik af Religionen. Religionen skyldes --- hvad enten
den fremtræder i individuel eller i social Form --- en Enhedstrang,
der hænger nøje sammen med Selvopholdelsestrangen.
Denne Enhedstrang fremtræder i Religionen særlig som Trang
til Ledelse (besoin de direction), og forsaavidt var Schleiermachers
Beskrivelse af Religionen som Afhængighedsfølelse et
% --- p. 40 ---
\opage{40}heldigt Greb. I sygelig Form kan Enhedsretningen, den koncentrerede
Afhængighed kun naas gennem Udelukkelse af forskelligartede
Indtryk og Retninger; i sund Form naas den paa
mere harmonisk Maade gennem „Opbyggelse“ og Organisation.

Religionen gaar i sin sociale Form forud for den individuelle
Religiøsitet; Religionshistorikerne beskæftige sig derfor mest med
hin og ere tilbøjelige til at betragte denne som et Opløsningstegn.
Naturfolkenes Religion bestaar mest i Trosforestillinger
og Ceremonier, som ere fælles for hele Samfundet. Det primitive
Menneskes Jeg har ikke udsondret sig af Familie, Klan og
Folk, ligesom Barnets Bevidsthed stadigt henter sit første Indhold
fra de nærmeste Omgivelser. Selve Extasen hos de Vildes
Troldmænd minder mere om Fanatikernes Aandstilstand end om
Mystikernes Extase; der hører en vis Udvikling til, før den
egentligt individuelle Form kan blive til. Men Religionens
Sociologi kan paa intet Punkt ganske skilles fra Religionens
Psykologi. Selve Trangen til at frembringe social Enhed eller
slutte sig til en saadan har en psykologisk Aarsag, der maa
studeres hos de enkelte Individer. Selv hvor Religionen nærmest
fyldestgør et socialt Krav, bæres den dog væsentligt af
en indre Higen i de Enkeltes Sind. ---

Der kan sikkert indvendes Adskilligt mod Murisiers Opfattelse
af Mystiken som et rent individualistisk Fænomen.
Dels begynder Mystikeren som Regel indenfor en konfessionel
Religion, og dels hører hans eget Jeg til det, han vil befries
fra; han vil ud af de snevre Forhold, ud over alle Grænser, og
Mystiken er ofte en til det Yderste dreven Altruisme. Mystiken
er et Expansionsfænomen, nok saa meget som et Koncentrationsfænomen,
eller rettere, der gives mangfoldige Nuancer
indenfor den mystiske Retning\textsuperscript{10}). --- Der gives ogsaa flere Former
for sygelig religiøs Følelse end de to Hovedformer, til hvilke
Murisier holder sig. Der er saaledes Tilstande af Angst og
Fortvivlelse, hvor hverken Expansion eller Koncentration er
mulig, og der er i Modsætning dertil den overvættes Selvfølelse,
hvor Æren og Magten ved at være knyttet til det Højeste
nydes paa egoistisk Vis. Denne Modsætning mellem potenseret
Afmagtsfølelse og potenseret Magtfølelse er maaske nok saa
% --- p. 41 ---
\opage{41}betydningsfuld for Religionens Patologi, som de af Murisier
beskrevne Former.

Hvad der ikke kommer tilstrækkeligt frem hos Murisier, er
den religiøse Følelses Sammenhæng med Værdierfaringer. Man
vilde kunne have en Ledetraad til at finde de sygelige Former
for Religion ved at gaa ud fra de Forudsætninger, hvortil
Studiet af Religionen atter og atter vil føre tilbage: Erfaringer
om Værdier, om Virkeligheder og om Forholdet mellem Værdierne
og Virkeligheden. Men det mest videnskabelige Synspunkt
vilde dog være det, der betragter Religionens patologiske
Former som specielle Tilfælde af de almindelige Sindssygdomsformer.
Anlæggelsen af dette Synspunkt vilde ogsaa føre til at
indse den Begrænsning, indenfor hvilken man maa bevæge sig,
naar man fra det syge Bevidsthedslivs Fænomener vil drage
Slutninger angaaende det sunde Bevidsthedslivs Natur. --- Disse
Indvendinger hindre dog ikke i at anerkende, at der i Murisiers
lille Bog findes megen god og fin Psykologi, og at den
--- indenfor den Begrænsning, Forfatteren har valgt --- belyser det
religiøse Liv fra en karakteristisk Side.

\whiteline

\afsnit{3. Bidrag til den konfessionelle Religiøsitets Psykologi.}{3. Bidrag til den konfessionelle Religiøsitets Psykologi}

Fra religiøs Side hævdes det mere og mere, at Grundlaget
for Religion er Erfaring, og ethvert Forsøg paa psykologisk
Undersøgelse af Religionen fører til samme Resultat. Men saa
bliver det af afgørende Betydning at undersøge Arten og Beskaffenheden
af den Erfaring, der her kan være Tale om. Ellers
kan man ikke vide, hvor Meget dette Grundlag formaar at bære,
og Dogmatismen holder igen sit Indtog. Men her indtræder
nu en Vanskelighed, som fremkommer ved, at enhver Følelsestilstand
har en Tilbøjelighed til hurtigt at indgaa Forbindelse
med bestemte Forestillinger og føre til bestemte Handlinger.
Og de Forestillinger og Handlinger, der stedse ligge nærmest,
ere de i det omgivende Samfund raadende. I dem lever Individet
sig ind, maaske endog før det har begyndt at gøre sine
egne Livserfaringer. Det arver Slægtens Vurderinger, før det
kan begynde at vurdere Livet paa egen Haand. Og naar det
% --- p. 42 ---
\opage{42}begynder at gøre sine egne Erfaringer, vil det i hvert Tilfælde
strax lægge sig dem til Rette efter de overleverede Former.
For en stor Del, om ikke for største Delen, fremkaldes dets
Erfaringer netop ved det Indtryk, overleverede Forestillinger
og Handlinger gøre paa det, især gennem Omgivelsernes Betagethed
af dem. Kun med stor Vanskelighed og under særegne
Forhold vil den Kontraprøve blive mulig, som behøves til
et virkeligt Erfaringsbevis, --- nemlig Prøvelsen af, om den
Erfaring og den Livsvurdering, Individet selv efter sin Natur og
sine Vilkaar kan naa til, nu ogsaa virkeligt ene og alene kunne
udtrykkes og forklares gennem de overleverede Former, saa at
Indhold og Form ganske svare til hinanden. Muligvis kunde
samme Indhold udtrykkes gennem andre Former, eller andre
Forestillinger og Handlinger kunde endog gøre en rigere og
værdifuldere Livserfaring mulig.

Den religiøse Erfaring bliver derfor, historisk set, stedse
mere eller mindre konfessionel. Den bliver overvejende en Gentagelse
af, hvad Andre forud have erfaret og oplevet; den bevæger
sig i færdige Udtryk og Formler, og det, der fra først af,
nemlig da Udtrykket fandtes, og Formlen skabtes, var Virkning
og Resultat, det virker i den konfessionelle Erfaring som Aarsag,
ofte Hovedaarsag. Dog er der indenfor den konfessionelle
Erfaring alle mulige Grader mellem den passive Indlevelse i
givne Former og den selvstændige Tilegnelse, der paa sit Højdepunkt
kan gaa over til at finde nye Veje og skabe nye Former.

Religionspsykologien maa have denne Kendsgerning for Øje,
naar den indsamler sit Erfaringsstof. Et rigt Materiale have vi
her --- foruden i hvad egen Menneskekundskab yder --- i religiøse
Personligheders Udtalelser, især i Selvbiografier. Det er
især til et saadant Materiale, jeg har støttet mig i det psykologiske
Afsnit af min Religionsfilosofi, idet jeg saa tillige har
sammenlignet de Typer, jeg ad denne Vej fandt, med de Typer,
de store Folkereligioner, især deres Grundlæggere, frembyde.
Men det vil naturligvis være af Betydning at faa en bredere
Basis, og dette vilde kunne opnaas, dersom der kunde samles
Oplysninger og Udtalelser fra større Kredse, ogsaa fra ikkeliterære
Kredse. Forsøg paa at skaffe saadanne støde paa mange
% --- p. 43 ---
\opage{43}Vanskeligheder, dels paa Grund af den Ubehjælpsomhed, som
vil hindre Mange i at udtale sig klart, dels paa Grund af den
Skyhed, der vil hindre Mange fra at udtale sig om det Inderligste
i deres Liv; det blotte Spørgsmaal vil maaske endog undertiden
blive betragtet som en Profanation. Men da psykologisk
Forskning som al anden Forskning øves for Sandhedens og Klarhedens
Skyld, vil det være fuldt berettiget at gaa denne Vej og
prøve, hvor langt den kan føre. Det er især i Nordamerika, man
har forsøgt det. I Nordamerika rører det religiøse Liv sig fra
den praktiske Side, og den herskende Metodisme med den dertil
knyttede Revivalbevægelse har naturligt henledet Opmærksomheden
paa de sjælelige Processer, der foregaa under religiøse
Oplevelser. Det er da ogsaa Omvendelsens Psykologi, de amerikanske
Psykologer især have beskæftiget sig med. Allerede
Revivalbevægelsens Ledere vare begyndte paa en Statistik over
Omvendelsernes Hyppighed, og de første Forsøg fra rent psykologisk
Side bleve til Dels gjorte i Samvirken med dem. Ved
udsendte Spørgsmaalslister har man da søgt at faa Oplysninger
om Tilstandene før, under og efter det religiøse Gennembrud,
saavel som om andre Typer for religiøs Udvikling, der ikke
havde Gennembruddets Karakter, om Kriser og Tvivl o. s. v.
Gennem Analyse af de fremkomne Udtalelser har man saa søgt
at naa til Beskrivelse af visse Hovedformer for religiøs Udvikling ---
naturligvis indenfor de givne historiske Forudsætninger
(nemlig den amerikanske Metodisme). Om meget store Tal kunde
der ikke her være Tale, og derfor ikke om streng statistisk
Metode. Man har kastet Nettet ud paa et enkelt Sted i de religiøse
Fænomeners Hav og beskrevet de Arter, man fandt ved
sit Dræt. Der er, som det Følgende vil vise, virkeligt vundet
nyt Stof for sammenlignende Psykologi ad denne Vej.

Den Første i Rækken af disse Forskere er \textsc{James Leuba}
med sin Afhandling \textit{Psychology of religious phenomena} (American
Journal of Psychology. April 1896). Leuba havde fra sin Barndom
af bevæget sig i Kredse, hvor religiøse Gennembrud vare
hyppige, og senere har han af psykologisk Interesse udvidet sin
Erfaring ved Samtaler med Mennesker, der havde oplevet saadanne,
ved at besøge Møder og ved at udsende Spørgsmaalslister
% --- p. 44 ---
\opage{44}angaaende Forholdene før, under og efter Omvendelsen. „Spørgsmaalslisterne“,
siger han, „bleve sendte til Personer, som mente
at have „erfaret Religion“, til Præster og til Ledere af indre
Mission, og gennem dem ogsaa til en stor Klasse, det ikke havde
været muligt at naa direkte til. De bekendtgjordes desuden i
„The Presbyterian“ i Filadelfia, i „The Christian“ i London og
i „The Outlook“ . . . . En lille Procentdel af de modtagne Svar
indeholdt de Talemaader, som ere gængse ved revivals og i
Missionshuse og smagte saa afgjort af religiøs Affektation (religious
cant) og saa lidet af personlig Erfaring, at de maatte
skydes til Side; men de fleste af Svarene bare indre Vidnesbyrd
om alvorlig Bestræbelse for at give nøjagtig Beskrivelse af en
Erfaring, man tillagde umaadelig Betydning . . . . . Udbyttet af
vore Spørgsmaal begrænsedes meget ved visse Folks Mangel paa
Tiltro til Nytten af at anvende Psykologi paa religiøse Fænomener,
ved de Vanskeligheder, sædvanlige Mennesker, der ikke
ere øvede i Selviagttagelse, føle ved at skulle gøre Rede for
deres egen Erfaring, ved nogle kirkelige Autoriteters fjendtlige
Holdning, og endeligt ved den Afsky, Opfordringen til at granske
religiøse Erfaringers Natur undertiden vakte.“ Angaaende dette
sidste Punkt skriver Leuba i en senere Afhandling Følgende:
„Den Grundtone af dyb, undertiden endog tragisk Alvor, der
gaar igennem de fleste af mine Dokumenter, vilde endog stemme
en ironisk Skeptiker til Ærefrygt. Den virkelige Værdi, man
tilskriver Religionen, aabenbarer sig, naar de intellektuelle Antagelser,
af hvilke det religiøse Liv menes at afhænge, trues . . . .
En Dame af mit Bekendtskab, hvem jeg gav en Spørgsmaalsliste,
fór tilbage i stor Rædsel, da hun havde læst de første
Spørgsmaal. Hendes uklare Begrundelse af sin Vægring ved at
svare godtgjorde, at hun ikke turde se sin Tro lige i Ansigtet;
blot Tanken derom indjog hende Skræk. En anden Dame, der
med paafaldende Heftighed havde afslaaet en personlig Anmodning
om Svar, skrev et Aar senere af egen Drift for at forsvare
sin Adfærd. Hendes Tro var paa den Tid, da Spørgsmaalet kom,
i en betænkelig Tilstand, og hun kæmpede med Fortvivlelsens
Energi for, at hendes Himmel ikke skulde miste sin Gud. Hendes
vigtigste Forsvarsmiddel var at forjage fra sin Bevidsthed
% --- p. 45 ---
\opage{45}enhversomhelst Ting, der kunde true hendes Tro. Den kritiske
Undersøgelse, jeg foreslog, saarede hende paa det Dybeste. Disse
gribende Exempler paa en Klasse af vore Medmennesker, der
foretrække at stikke Hovedet i en Busk fremfor at lade de haarde
Fakta berøve deres Verden den Sjæl, der giver den Ynde, vækker
til betydningsfuld Eftertanke om Sammenhængen mellem
Tro og Trang.“

I en Række senere Afhandlinger har Leuba fortsat sine Undersøgelser
\textit{(The personifying passion in youth.} The Monist.
Juli 1900. --- \textit{Introduction to a psychological study of religion.} ibid.
January 1901. --- \textit{The content of religious consciousness.} ibid. July
1901. --- \textit{Les tendances religieuses chez les mystiques chrétiens.} Revue
philosophique. Juillet--Novembre 1902). Hvad der karakteriserer
Leuba er hans Interesse for de individuelle Tilfælde i deres
Mangfoldighed og Ejendommelighed. Han har tilvejebragt en
Samling af religiøse Selvkarakteristiker, der vise, hvormange
Nuancer og Overgange der findes paa dette Omraade, og hvor
lidt vore sædvanlige Rubriker slaa til. ---

I Anlæg og i Betydning beslægtede med Leubas Undersøgelser
er en nyligt udgiven Afhandling af \textsc{Lucien Arréat}:
\textit{Religion in France.} (The Monist. January 1903). Ved sin Undersøgelse
af fransk Katolicisme giver han et interessant Modstykke
til Leubas Undersøgelse af amerikansk Metodisme. ---

\textsc{Edwin Starbuck} \textit{(The psychology of religion.} An empirical
study of the growth of religious consciousness. London
1899) raader over et meget omfattende statistisk Materiale, som
han har underkastet en gennemgaaende Analyse. Han opfatter
Omvendelsen som en Proces, gennem hvilken den Unge skal
gøre det socialt Overleverede til fuld Alvor og Virkelighed i
sig, og derfor som et sociologisk Fænomen, der kan paavises hos
alle Folk og paa alle aandelige Omraader. Af særlig Interesse
er hans Paavisning af forskellige Arter af Omvendelse og hans
Sammenligning af Omvendelsesprocessen med andre Former for
religiøs Udvikling.

\textsc{George Coe} \textit{(The spiritual life.} Studies in the science of
religion. New-York 1900) har ad lignende Vej som Leuba og
Starbuck særligt undersøgt Temperamentets Indflydelse paa den
% --- p. 46 ---
\opage{46}religiøse Udviklings Karakter og dets Forhold til andre medvirkende
Elementer.

\textsc{George Santayana} har i sit Skrift \textit{Interpretations of
Poetry and Religion} (London 1900) særligt fremdraget det poetiske
Element i Religionen og er derved bleven ført til at stille sig
i Modsætning baade til Ortodoxien og til den gængse Kritik.
For ham er Religionen den højeste Poesi, medens han tillige
søger at vise, hvorledes Poesien paa sit Højdepunkt gaar over
til Religion. ---

Gennemgaaende er de amerikanske Psykologers Standpunkt
overfor det religiøse Liv et sympatisk. De hævde energisk
Psykologiens Selvstændighed overfor al Teologi og Spekulation,
men de fordybe sig med Begejstring og Interesse i de religiøse
Fænomener. Deres eget religiøse Standpunkt synes ikke at
være udpræget antikonfessionelt, selv om de mene at kunne
paavise Nødvendigheden af, at den konfessionelle Begrænsning
viger for den værdifulde Mangfoldighed af Veje og Typer,
uhildet Iagttagelse opdager. Saavidt muligt vil jeg i den
følgende Drøftelse af deres Undersøgelser ogsaa søge at bestemme,
til hvilken Type de selv høre, for derved at bestemme
den personlige Ligning og angive en Begrænsning for deres
Undersøgelsers Betydning.

\underafsnit{a. Gennembrud og Kontinuitet i den religiøse Udvikling.}{a. Gennembrud og Kontinuitet i den religiøse Udvikling}

\textsc{Starbuck} fandt i det af ham indsamlede Materiale to Typer
af religiøs Udvikling. Den ene bevægede sig gennem en mere
eller mindre pludselig Afgørelse frem til en ny Livsretning; den
anden foregik mere kontinuerligt.

\textit{Gennembrudstypen} frembyder igen Forskelligheder af Interesse.


Alderen, i hvilken Gennembruddet finder Sted, kan være
forskellig. Det er hyppigst i de unge Aar og indtræder noget
tidligere for Kvinder end for Mænd; senere end i 40 Aars Alder
forekom det sjeldent (indenfor de 1265 Tilfælde, hvor Alderen
kunde konstateres). --- Anledningen til Gennembruddet var forskellig.
En stor Del af Tilfældene skyldtes Vækkelsesmøder
(revivals), Nogle fandt Sted under sædvanlig Gudstjeneste,
% --- p. 47 ---
\opage{47}Ved andre igen kunde ingen særlig ydre Anledning paavises. ---
Hvad Motiverne angaar, forekom Frygt langt hyppigere end
Haab, og altruistiske Motiver (Idealitet og Menneskekærlighed)
vare sjeldne. Karakteristisk er det, at Syndsfølelse var sjeldnere
ved revivals end ved andre Anledninger; Starbuck mener,
at ved revivals faar man ikke Tid til at fordybe sig i sig selv
og gruble over sin Færd og sin Fortid.

Vigtigst ere Forskellighederne med Hensyn til selve Gennembruddets
Karakter. Der fremtræder her to Former, som dog
kun tilnærmelsesvis kunne paavises i deres Renhed. I den ene
Form er det en Vane, der brydes, en lastefuld Tilstand, fra
hvilken man undslipper; forud gaar da en dyb Følelse af Synd.
I den anden Form er det et nyt Lys, der bryder ind; der udfyldes
et Savn og naas en positiv Fuldendelse; forud gaar da
en Følelse af Ufuldstændighed. Hin Form har en negativ, denne
en positiv Karakter. Den negative Form behøver fuldstændig
Overvindelse og Fortrængen af Fortiden for at blive til; den
positive Form behøver kun en lille Impuls for at træde i Kraft.
Starbuck bemærker, at indenfor Grænsen af hans Materiale er
den positive Type (spiritual illumination, forud for hvilken der
findes a sense of incompleteness) hyppigere end den negative (escape
from sin, forud for hvilken der gaar a feeling of wretchedness),
og han fremhæver det Betænkelige ved den i Kirken herskende
Tendens til at gøre den negative Type til den ene rette og
normale. De, der kæmpe for fuldere Lys, maa behandles anderledes
end de, der kæmpe for at befries fra Mørke. Ogsaa i det
af Leuba indsamlede Materiale forekommer den positive Type
hyppigst. Starbuck og Leuba drøfte ikke Aarsagen hertil. Det
er ikke ganske sikkert, at den positive Type i Virkeligheden er
den hyppigste; maaske er den negative Type mere sky og utilbøjelig
til at udtale sig, saa at deri Grunden ligger til, at der
kan konstateres færre Tilfælde af den.

Af meget stor Interesse er Vexelvirkningen af det Bevidste
og det Ubevidste ved Gennembruddet. Der er her mange forskellige
Grader af Bevidsthed, hvilket sandsynligvis hænger
sammen med, at der er forskellige Grader i de indre Hindringer,
der skulle overvindes. En Forudsætning for Gennembruddet
% --- p. 48 ---
\opage{48}er det --- indenfor begge Former --- at der er foregaaet en Væxt
under Bevidsthedens Tærskel (within the sub-liminal consciousness).
Nu er det ikke sagt, at den bevidste Stræbens Retning
falder sammen med den ubevidste Væxts Retning, --- og da saa
alligevel Villiens Tilslutning til det ubevidst Anlagte er nødvendig,
for at Gennembruddet kan indtræde, er det intet Under,
at der her foreligger saa mange psykologiske Gaader. Hvad der
skal komme i Stand, er jo, at det, der hidtil kun har staaet
som Genstand, som fjernt og ophøjet Ideal, skal blive selve det,
der lever og rører sig i Individet. Hvad der hidtil kun var
Objekt, skal blive Subjekt.

Er Gennembruddet indtraadt, opstaar der en Følelse af
Nyhed, i Regelen ogsaa af Glæde og Fred, og hyppigt af Selvforglemmelse
og Menneskekærlighed. Hvad der hidtil holdt
Selvet tilbage indenfor snevre Skranker, er fortrængt, og en
Udvidelse kan ske.

Psykologisk set er Omvendelsen i begge dens Former ikke
noget enestaaende Fænomen. Der gives en Række analoge Erfaringer,
saaledes naar Uvished og Nedstemthed pludseligt afløses
af frejdig Afgørelse, eller naar en lang Anstrengelse for
at tilegne sig en Færdighed som med ét Slag afløses af Herredømme
over Vanskelighederne, eller naar en Følelsestilstand
melder sig pludseligt ved Kontrast eller forberedt ved indre,
organiske Processer. Især forekomme saadanne Fænomener i
den Alder, da den Unge skal indleve sig i Slægtens Vaner og
Traditioner og gøre dem til sine egne --- skal udvide sit Selv
til at omfatte større Interesser, nedbryde Skranken mellem den
individuelle og den sociale Villie. Fysiologisk set stilles denne
Opgave paa en Tid, hvor Væxten nærmer sig sin Fuldendelse,
og hvor nye Opgaver og Funktioner begynde at kræve deres
Ret. Det er en Bevægelsens og Modsætningernes Tid; baade i
Organismen og i Bevidstheden træder Modsætningen mellem det
Gamle og det Nye skarpt frem. Men ved religiøs Paavirkning
forstærkes denne Modsætning ofte paa unaturlig og usund Maade.
Man anerkender, som allerede bemærket, ikke altid Forskellen
mellem de to Typer, men behandler dem ens. Starbuck anfører
et Exempel paa, hvorledes det kan føre til Fortvivlelse,
% --- p. 49 ---
\opage{49}naar Alle skulle tvinges ind i ét bestemt Udviklingsskema og
nødes til at gøre Erfaringer, som for dem ere umulige. ---

I Modsætning til Gennembrudstypen i dens to Former stiller
Starbuck \textit{Kontinuitetstypen.} Der er her ingen Bevidsthed om
afgørende Forandring, om Modsætning mellem Gammelt og Nyt.
Der sker en uvilkaarlig Opvaagnen, som vel har en vis Lighed
med den positive Gennembrudsform, men mangler Erfaringen
om pludselig Forandring. Udviklingen foregaar kontinuerligt
fra Barndommen gennem de unge Aar. Dog kan der komme en
Periode af „Storm og Trængsel“ (storm and stress), der hænger
sammen med Overgangsaarenes organiske Forandringer og kan
medføre Angst og Spænding. Det er den Periode, hvor den
Unge skal indleve sig i en større Verden og gøre Slægtens Erhvervelser
til sine egne. Mangfoldige Fordringer og nye Muligheder
fremstille sig. I intellektuel Form (og noget senere) fremtræder
den saaledes vakte Uro som Tvivl. Det gælder om med
klarere Bevidsthed at fastholde, hvad man hidtil havde antaget
paa Autoritet og med barnlig Troskyldighed. En saadan Periode
var meget hyppig hos dem, Starbuck lærte nærmere at kende;
den kunde begynde, forløbe og afslutte uden bratte Overgange.
Ofte førte den over til en Tid, hvor Individet stod fremmed
overfor det Overleverede og betragtede dette som aldeles værdiløst.
Denne Fremmedhedsperiode (period of alienation) skyldtes
ofte Slappelse efter Uroen eller Tvivlen; men den kunde ogsaa
have sin Grund i, at Individet nu vilde hævde sig selv overfor de
Former, Samfundet vilde føre det ind i, eller i Lysten ved ubunden
Brug af Evner og Kræfter. 1 Stedet for den religiøse Tradition % PRINTED AS IS: 1 for I
træder i en saadan Periode ofte æstetiske, intellektuelle
eller andre Interesser i Forgrunden. Ved Slutningen af Ungdomsperioden
indtræder der som Regel en Rekonstruktion: Individet
har vundet en bestemt Tro, som ikke altid stemmer med
den overleverede, men i Regelen er en Blanding af selvstændig
Produktion og Genoptagelse af det Gamle. Starbuck er forundret
over at finde denne Periode saa bestemt udpræget. Han
forklarer sig det ved, at Livsforholdene kræve et bestemt Standpunkt
i praktisk Henseende, saa at der maaske endog i Slægten
har udviklet sig en Tendens til saadan Rekonstruktion efter de
% --- p. 50 ---
\opage{50}bevægede Ungdomsaar. --- Gennem alle disse Perioder forløber
Udviklingen indenfor denne Type saa kontinuerligt, at det ofte
kun er gennem Sammenligning af fjernt fra hinanden liggende
Punkter, at man bliver opmærksom paa dem.

Karakteristisk for Udviklingens Afslutning er --- og dette
gælder begge Hovedtyper --- en vis Ændring i de religiøse Forestillinger
og i hele Opfattelsen af Religionen. Selv om de konfessionelle
Forestillinger væsentligt fastholdes, kan der dog i
Regelen paavises en gennemgaaende Forskydning fra det Ydre
til det Indre, fra det Dogmatiske til det Etiske, fra det isoleret
Personlige til Bevidstheden om Delagtighed i et mere omfattende
Liv. Religionen opfattes mere og mere som et indre Liv,
der er Del af en universel Tingenes Orden.

\textit{Sammenlignes nu til Slutning Gennembrudstypen med Kontinuitetstypen,}
fremtræder hin selvfølgeligt med langt større Uro
og med „Storm og Trængsel“ i mere udpræget Grad end denne.
Gennembrudstypen er mere modtagelig for ydre Paavirkning;
Modsætningen mellem det Gamle og det Nye gør sig stærkere
gældende i den, og Følelseslivet er heftigere. Gennembrudstypen
holder sig desuden langt mere end Kontinuitetstypen til de bestemte
konfessionelle Former, til de organiserede Institutioner.
Sine indre Erfaringer udtrykker den helst i bibelske og konfessionelle
Udtryk og har ingen Trang til at danne selvstændige
Former og Udtryk eller finde nye Analogier. En saadan Trang
er derimod meget hyppig indenfor Kontinuitetstypen, hvor ogsaa
Opfattelsen af Religionen som et indre Liv forekommer hyppigere.
--- Starbucks almindelige Resultat er, at det religiøse
Livs Udvikling under konfessionelle Former kan belyses fra almindelige
psykologiske og fysiologiske Synspunkter. Der udfoldes
i Religionen Anlæg og Tilskyndelser af almindelig menneskelig
Natur, og ogsaa de forskellige Former, i hvilke den
religiøse Udvikling foregaar, kendes fra andre sjælelige Omraader.
Selv betragter han Religionen „i dens højeste Form“ som en
Udstraaling, et Sammenspil og en Idealisering af Tendenser, der
allerede forud ere til Stede i den menneskelige Natur. De Former
og Udtryk, Slægten har fundet herfor, føres den Unge i
Overgangsaarene for Alvor ind i, og Udviklingens Gang og Kri%
% --- p. 51 ---
\opage{51}ser betinges ved de Vanskeligheder, her ere at overvinde. Han
er tilbøjelig til at mene, at de Ulemper, som komme frem under
denne Udvikling, langt mere skyldes de Ældres Mangel paa
Forstaaelse af Menneskenaturens Krav end Særheder og Mangler
hos de Unge. --- Med denne sidste Bemærkning kan sammenstilles
en Udtalelse af Professor \textsc{Peabody} i hans Afhandling
\textit{The religion of a college student.} (Forum. June 1901. p. 447): „For
én ung Mand, der fjernes fra Religionen gennem Paavirkning af
Videnskab, Filosofi eller Kunst, er der ganske sikkert ti, det
gaar saaledes paa Grund af uforstandig eller upraktisk Religionsundervisning . . . .
Det er ikke Universitetet, der skal forny
Overensstemmelsen med Kirken; det er Kirken, der maa aabne
sine Øjne for den vidunderlige Udvidelse af den intellektuelle
Horizont, som enhver Studerende i vore Dage har Adgang til
at erfare“.


\underafsnit{b. Temperamentets Indflydelse.}{b. Temperamentets Indflydelse}

Ogsaa \textsc{George Coe} fremhæver Betydningen af et psykologisk
Studium af Religionen i dens nuværende Tilstand. Og
det er ogsaa Ungdomstiden, han særligt har henvendt sin Opmærksomhed
paa. Denne Tid er, paastaar han, langt mindre
kendt og forstaaet end Barndommen og den modne Alder. Hvad
der foregaar i denne Tid, er en Overgang fra at være optagen
af det Enkelte til at fatte det Universelle, fra at overvældes af
det Ydre til at fordybe sig i det Indre, fra at lade sig lede af
Autoritet til at bestemmes ved Grunde og fra at stanse ved det
Givne til at stræbe mod Idealet. Der spørges mere indtrængende
om Livets Betydning, og de sociale Interesser begynde
at vaagne. Alt dette gør, at der føles Trang til en mere personlig
Tilegnelse af Religionens Indhold, end der er Mulighed for i
Barndommen. Omvendelse eller en tilsvarende Inderliggørelsesproces
er derfor et naturligt Led i Ungdomstidens Udvikling, og
Coe finder ogsaa faktisk, at en saadan Forandring er overordentligt
hyppig (især ved 16 Aars Alder). Han minder om, at det
hos Folkeslag rundt om paa Jorden er almindeligt at henlægge
en religiøs Indvielse til dette Tidspunkt. Ofte sendes den unge
Mand ud i Ensomheden for at vente paa en Aabenbaring fra
% --- p. 52 ---
\opage{52}Stammens Gud og for ved dens Hjælp at blive klar over sit
Kald. I Tilslutning til Starbuck hævder han, at det hos unge
Amerikanere indenfor Kristendommen ingenlunde blot er Frygt
eller Anger, der ere Motiver ved den religiøse Udvikling, men
ogsaa mere ubestemte Tilstande: Uro, Tomhed, Følelse af Ufuldstændighed,
Higen mod det Uendelige, --- selv om de Fleste udtrykke
deres Erfaringer i det konfessionelle Sprog. Det viser
sig fremdeles gennemgaaende, at rent intellektuelle Problemer
og Motiver ere underordnede; Grundene findes bagefter; man
søger først at leve det Liv, der udtrykker det dybeste Selv, og
siden ser man maaske til, hvorledes man ved Fornuftens Hjælp
kan vise Berettigelsen af det. Og det er ikke nogen enkelt Erfaring,
der fører til Maalet. Det Meste af, hvad Ungdommen
fordrer afgjort paa én Gang, kan man kun leve sig ind i efterhaanden,
alt som Livet udfolder sig med sit Arbejde, sine Skuffelser
og sine Sejre.

Coe vil studere disse Fænomener uden teologiske og metafysiske
Fordomme. Fristelsens Psykologi f. Ex. er og bliver den
samme, hvorledes man saa forklarer sig Oprindelsen til de Tilskyndelser,
der medføre den, --- om man mener, at de skrive sig
fra Djævelen, eller at de ere Rester af den Dyriskhed, hvorfra
Menneskeheden har udviklet sig. Det Samme gælder om Omvendelsens
Psykologi, som Coe særligt har beskæftiget sig med.
Det er hans Opgave at paavise de vigtigste Faktorer, som medvirke
ved dette psykiske Fænomen. Han gaar her den Vej at
anstille en nøjagtig Undersøgelse af en forholdsvis lille Række
af individuelle Tilfælde. Foruden Udsendelse af Spørgsmaalslister
anvendte han indgaaende Samtale med Enkelte, indhentede
Andres Vidnesbyrd om dem, anvendte Hypnose, søgte ad alle
mulige Veje at naa til en individuel Karakteristik. Antallet af
Individer, han undersøgte, var 77 (52 Mænd, 25 Kvinder); de
vare næsten alle Studerende, Mændene gennemgaaende 24 Aar
gamle, Kvinderne 22 Aar; de Fleste vare opdragne i den metodistiske
Kirke, der jo lægger saa stor Vægt paa personlig religiøs
Udvikling.

Det viste sig da, at hos saadanne unge Mennesker fremtræder
der store Forskelligheder med Hensyn til Omvendelsen. Der
% --- p. 53 ---
\opage{53}kan i denne Henseende skelnes mellem tre Klasser af Individer:
A) saadanne, som ventede en Forvandling (transformation) og
virkeligt erfarede den; B) saadanne, som ventede en Forvandling,
men ikke erfarede den; C) saadanne, som ikke tilfredsstilledes
ved, hvad de erfarede, eller som dog stræbte at opnaa en
fyldigere Erfaring.

Ved nærmere Undersøgelse viser det sig nu, at Klassen A
overvejende var Følelsesmennesker (sensibility distinctly predominant),
medens Klassen B overvejende var Forstandsmennesker
(intellect distinctly predominant). Og medens Klassen A
nærmest mindede om de saakaldte melankolske og sangvinske
Temperamenter (the slow-intense and the prompt-weak varieties),
mindede Klassen B mere om det saakaldte koleriske Temperament
(the prompt-intense variety). Klassen C frembød Forskelligheder,
der snart mindede om A og snart om B. Temperamentet
viser sig saaledes at være en væsentlig Faktor.

Hermed forbinder sig nu en anden Iagttagelse. I Klassen
B forekom saa godt som ikke automatiske Fænomener i Sammenhæng
med de religiøse Tilstande, hverken paa Forestillingslivets
Omraade (stærke Drømme, Hallucinationer), eller paa Bevægelsesorganernes
Omraade (ubetvingelig Latter, Sammenfaren,
Klappen i Hænderne, Højttalen med sig selv o. s. v.), medens
saadanne Fænomener langt hyppigere forekom i Klasse A (og i
C). --- Naar Klasserne desuden undersøgtes med Hensyn til
Modtagelighed for Suggestion og Hypnose, viste det sig, at A’s
Suggestibilitet overvejende er passiv, idet Individerne her ganske
rette sig efter den givne Indskydelse og følge den i det Enkelte,
medens Klassen B’s Suggestibilitet kan kaldes spontan,
idet Individerne enten give sig selv deres Indskydelser eller
selvstændigt føje Noget til den oprindelige Indskydelse.

Foruden den vakte Forventning spiller altsaa Temperament,
Tendens til Automatisme og Suggestibilitet en Rolle ved Omvendelsen.
Karakteren af den religiøse Erfaring bestemmes ved
disse Faktorer. ---

Ligesom Starbuck fremhæver Coe meget stærkt Faren ved
at anse en eneste Type for ene berettiget. Vækkes der en Forventning
om at opleve en mere eller mindre pludselig Sinds%
% --- p. 54 ---
\opage{54}omvæltning, og denne Forventnings Opfyldelse efter Individets
hele psykologiske Natur er umulig, saa kan det føre til Fortvivlelse.
Coe mener, at det egentligt kun er det melankolske
og det sangvinske Temperament, Kirken har anerkendt i sin
Kultus, sine Hymner og sin Moral\textsuperscript{11}). I den Kristustype, som
raader i Kirken, lægger dette sig særligt for Dagen; det er
Sorgernes Mand, det lidende Offerlam. Den virkelige Jesus led,
fordi han trods al Modstand fastholdt sin Tanke til det Yderste;
han var en aktiv, ikke en passiv Individualitet; der var et
kolerisk Element til Stede, der er forsvundet fra det traditionelle
Billede. --- Kirken maa være ligesaa bred, som Menneskenaturen
er det. Men den har vænnet sig til at se paa Tingene med
kvindelige Øjne; derfor er det især Mændene, der vende sig
fra den.

Coe finder, trods sin øjensynlige Sympati for Religionen,
ingen Betænkelighed ved en gennemført psykologisk Drøftelse
af de religiøse Tilstande. Hvis vi --- siger han --- ville forstaa
religiøse Erfaringer, maa vi ikke forklare dem ved Guds urandsagelige % PRINTED AS IS: urand-|sagelige
Veje, men ved paaviselige Ejendommeligheder i Menneskets
Natur; vi skulle ikke spekulere over den guddommelige
Naade, men studere den menneskelige Aands skjulteLiv. Ligesom % PRINTED AS IS: skjulteLiv
andre Religionsfilosofer i vore Dage finder han den sidste
Maalestok for religiøse Værdier i det Etiske. Det kommer an
paa, hvad der ved religiøs Udvikling vindes for et dybere, indholdsrigere,
aktivt Liv.

\underafsnit{c. Forskelligheder i Forholdet til Konfessionen.}{c. Forskelligheder i Forholdet til Konfessionen}

Der foreligger her, som allerede omtalt, to Undersøgelser,
en af Leuba om amerikansk Metodisme, en anden af Arréat om
fransk Katolicisme.

\textsc{Leuba} har dels studeret det religiøse Liv, som det fremtræder
hos Samtidige (i hans Hjemland Nordamerika), dels
Mystikernes Religiøsitet. Begge Studier føre ham til at gøre en
bestemt Forskel mellem Religionshistorikerens og Religionspsykologens
Synspunkt. Historikerens Emne er ikke saa meget
selve de religiøse Erfaringer, som de i Kultus, Myte og Dogme
foreliggende Produkter af disse Erfaringer, saaledes som de
% --- p. 55 ---
\opage{55}forme sig indenfor et Samfund. For at naa til den umiddelbare
religiøse Erfaring maa man henvende sig til den individuelle
religiøse Bevidsthed. Af den bæres og betinges tilsidst al
Religion, og ved de højere Folkereligioner bliver dette Forhold
særligt tydeligt. Det er ikke Guderne, men Forholdet til
Guderne, der interesserer Religionspsykologien. Dennes Opgave
er at give en religiøs Dynamik, d. v. s. en Lære om det religiøse
Livs Motiver. (Udtrykket „religiøs Dynamik“ havde Coe
brugt før Leuba). For at faa Materiale hertil fra Samtidige
har Leuba gennem udsendte Spørgsmaalslister indhentet Udtalelser
fra en Række Personer, der stode paa meget forskellige
Standpunkter, men alle indenfor den angloprotestantiske
Verden. Spørgsmaalene angik Grundene til deres Deltagelse
eller Ikke-Deltagelse i Kultus (offentlig eller privat Kultus), ---
de Omstændigheder og Lokaliteter, der øvede særlig religiøs
Indflydelse paa dem ---, de mest karakteristiske religiøse Erfaringer,
de kendte, --- de Forskelle, de gjorde mellem religiøse
og ikke religiøse Tilstande og Handlinger. Det er lykkedes
ham at samle en Række karakteristiske Udtalelser. Om Statistik
er der ikke Tale; men Leuba mener at kunne paavise visse
Ejendommeligheder ved det religiøse Liv i den Kreds, han har
med at gøre.

Et Hovedpunkt er for ham det Resultat, at hvad der ytrer
sig i Religionen, ikke er en Stræben efter Erkendelse, men en
Stræben efter Væren. Religionen udspringer, efter hans Opfattelse,
af en Følelse af Brudstykkeagtighed (unwholeness) i
Forbindelse med en Følelse af indre Delthed og Splittethed.
Hin Følelse medfører Trang til at blive delagtig i et større,
mere omfattende Liv, medens Delthedsfølelsen medfører Trang
til at ledes af et eneste Formaal (unification of motives). Syndfølelsen
er en Form for denne sidste Følelse. Til de to nævnte
Motiver slutte sig endnu Altruisme og Afhængighedsfølelse;
den første er beslægtet med Trangen til Enhed, den anden er
betinget ved den indre Modstand, der er at overvinde under
Bestræbelsen for at naa til Enhed. De store Religionsstiftere
agtede ikke paa teoretiske Problemer, men stræbte efter ideal
Fuldkommenhed og indre Enhed. Og der ytrer sig nu til Dags
% --- p. 56 ---
\opage{56}mere og mere en Trang til at gaa tilbage til de umiddelbare
religiøse Erfaringer, af hvilke Religionen stedse fra først af er
voxet frem. Skønt man endnu ikke har haft Mod eller Forstand
nok til at gøre dette fuldt ud, saa ere dog de første
Skridt i denne Retning (to return to plain religious empiricism)
gjorte, især indenfor Protestantismen. I Stilhed føres der
--- ved Siden af Kampen om Tilværelsen mellem Individerne indbyrdes
--- en rent indre Kamp, der i Strenghed og tragisk Alvor
ofte overgaar den ydre. Den føres paa forskellig Maade indenfor
de forskellige Religioner og finder Udtryk i forskellige
Former; men stedse hænger den sammen med Trangen til et
større og tillige mere udelt Liv. Den religiøse Trang er egentligt
ingen ganske speciel Trang, men har sit dybe og varige
psykologiske og fysiologiske Grundlag.

Ved Sammenligning mellem Religionen, som den fremtræder
i nyere Tid i Modsætning til tidligere Former for den, finder
Leuba en Række af Forskelligheder.

I primitive Samfund er Religionen et sammensat Produkt
af alle Menneskets fundamentale Fornødenheder og Bestræbelser,
baade teoretiske og praktiske. De teoretiske Motiver have
nu udskilt sig og finde deres Tilfredsstillelse i Videnskaben; det
er ikke paa dem, at Religionens Betydning i nyere Tid beror,
og allerede i Oldtiden staar Buddhismen som Vidnesbyrd om,
at disse Motiver ikke ere de afgørende. Selv om alle hidtil til
Religionen knyttede intellektuelle Elementer falde bort, vil Religionens
psykologiske Væsen, saaledes som det ovenfor er beskrevet,
kunne bestaa.

En anden betydningsfuld Forskel mellem mere primitiv og
mere moderne Religion finder Leuba deri, at medens forhen
Følelserne af Lydighed og Ærefrygt vare fremherskende, ere
disse nu --- i hvert Tilfælde i den angelsaksiske Verden --- ikke
saa hyppige, som det antages. Ikke den Lydighed, Katekismen
indskærper, men den klare Forstaaelse af, hvad der er ret og
godt, er nu Motiv til Handlen, og Enheds- og Frihedsfølelsen
er ved at fortrænge Underkastelsen. Hermed staar det Træk i
Forbindelse, at religiøs Sandhed ikke længere betragtes som
Noget, der kan og skal begrundes ad objektiv Vej; det Af%
% --- p. 57 ---
\opage{57}gørende er, hvad der tilfredsstiller den personlige Trang og Higen.
Mange af de Udspurgte udtalte dette med største Uforbeholdenhed.
Følelsesmotiverne havde afgjort Overvægten over alle intellektuelle
Interesser, og der var maaske endog fuldstændig
Ligegyldighed overfor den rent objektive Betydning af de religiøse
Forestillinger, selv hos Individer, som ellers paa speciellere
Omraader viste sig i Besiddelse af kritisk Sans og Evne.
Man \textit{kender} ikke Gud; man \textit{forstaar} ham ikke; men man \textit{bruger}
ham. Ikke Gud, men Liv, mere Liv, højere, rigere, mere fyldestgørende
Liv er Religionens Formaal. I Fortiden viser dette
sig klart hos de saakaldte Mystikere, og her slutter Leubas
Studier over Mystiken sig til hans Studier over samtidige Fænomener.
Det mystiske Liv maa ikke fortolkes og bedømmes
efter de Begreber og Dogmer, Mystikerne anvende. Dets Ejendommelighed
er kun en mere extrem og koncentreret Form for
Ejendommeligheden ved al Religion: en Higen efter højere og
fuldstændigere Tilfredsstillelse, efter Fred og Enhed i Villien.
Mystikernes Higen er mere lidenskabelig end den sædvanlige
Religiøsitets, og Koncentrationen mod et eneste Maal er stærkere.
Maal og Retning ere ens, men de gaa frem med langt
større Energi.

Endnu et karakteristisk Træk i den nyere Tids Religiøsitet
finder Leuba deri, at der fremtræder saa mange og betydelige
individuelle Forskelligheder, selv indenfor en og samme Civilisation.
Under Indflydelse af den uendeligt større Horizont,
Livet har faaet siden Konciliernes Tid, er ogsaa Religionen
undergaaet Forandringer, og der viser sig Forskelligheder ikke
blot i de ydre Ting ved Religionen, men i Individernes inderste
Sjæleliv. Fornødenhederne ere forskellige, og Forhaabningerne
ere forskellige. Og dog gaar man i Teologien ud fra, at en og
samme Føde passer for Alle. ---

Indenfor sin Samling af Selvbekendelser paaviser Leuba en
hel Række af Nuancer ligefra et Standpunkt, hvor det Religiøse
og det Konfessionelle gaa i den Grad i Et, at Individet ikke
formaar at udtrykke sig om sine indre Oplevelser uden i bibelske
og konfessionelle Udtryk, til saadanne Standpunkter, hvor der
aldeles ikke kan forbindes nogensomhelst Mening med den Art
% --- p. 58 ---
\opage{58}Udtryk. Mellem begge Yderpunkter ligge de interessanteste
Tilfælde; ti hverken af dogmatisk Position eller af dogmatisk
Negation er der nogen synderlig Psykologi at lære. Kun nogle
enkelte Træk kunne her anføres fra de interessante Selvkarakteristiker,
Leuba har samlet.

En videnskabeligt dannet ung Kvinde, Tilhænger af den
presbyterianske Kirke, skelner bestemt mellem, hvad der skyldes
hendes konfessionelle Opdragelse, og hvad der skyldes Trangen til
Noget, som har evig Værdi (a desire for things of permanent
value). Disse to Ting udelukke ikke hinanden, men de falde
heller ikke sammen. Naar hun deltager i sin Kirkes Kultus, er
det dels af Vane, men dels ogsaa fordi hun trænger til at træde
i Forhold til et højere Væsen og befries for selviske Tanker.
Forskellen mellem religiøse og ikke-religiøse Tanker og Følelser
er for hende en Gradsforskel og beror paa, at hine fremtræde
som betydningsfuldere og stærkere. Men den ideale Trang tilfredsstilles
aldrig fuldt. „Medens de gængse religiøse Paastande
erklære Religionen for den eneste Ting i Verden, der kan stille
tilfreds, beror“, siger hun, „for mig dens Værdi paa, at den
stadigt driver mig frem til nye Anstrengelser.“ --- Andre Korrespondenter
fremhæve Menneskekærlighed, barnligt Sind eller
Taknemmelighed som det Væsentlige i, hvad de kalde Religion.
--- En Læge, som deltager i Kirkens Kultus, skønt hans
egen private Andagt for ham har en mere afgjort religiøs Karakter,
udtaler bl. A. Følgende: „Jeg er langt fra at have den
antropomorfistiske Forestilling om Gud, som jeg havde i mine
unge Aar. Ved Gud forstaar jeg den Magt eller det Princip,
der bevirker al Udvikling i Universet henimod en højere Tingenes
Orden, og som ved sine Love hæver Livets Stade fra det
Ubevidste til det Bevidste, fra det Dyriske til det Fornuftige,
fra det Forgængelige til det Uforgængelige. Hos mig selv kalder
jeg alle Drifter, Tendenser og Anstrengelser guddommelige, som
føre til at udvikle hos mig større Modtagelighed for Skønhed,
Sandhed og Retfærdighed, højere Fornuft og Evne til moralsk
Fremgang“. --- En beslægtet Udtalelse fra en anden Korrespondent
er følgende: „Jeg har opgivet ethvert Forsøg paa at forestille
mig Gud og føler ikke mere nogen Trang til et Billede
% --- p. 59 ---
\opage{59}eller en Iklædning af det guddommelige Princip, som det er en
Nødvendighed for mig at tro paa. Mit religiøse Forhold til det,
man kalder Gud, kan jeg beskrive fra to Synspunkter: 1) Jeg
føler i min inderste Bevidsthed Trang til at være Et med et
stort Princip, Trang til Enhed med Noget, som jeg ikke kan
definere, men som omgiver mig, som jeg saa at sige er dykket
ned i, og som jeg ønsker at være i Harmoni med, et altomfattende
Nærværende .... 2) Jeg har næsten uafbrudt en dyb
Bevidsthed om at staa langt under et indre Ideal, en Type, der
er anlagt i mig, og som jeg stedse har for Øje. Jeg føler, at
jeg \textit{maa} virkeliggøre denne Type, der stadigt bor i mig og søger
at lede mig. Erfaring har lært mig, at jeg bliver lykkelig ved
at følge den, ulykkelig ved at fjerne mig fra den. Til disse to
fundamentale Instinkter eller Intuitioner har jeg reduceret alle
mine religiøse Forestillinger. Jeg er en Discipel af Kristus og
stræber at efterligne ham; mit Ideal er at blive Et med det
Guddommelige, ligesom han var det. Men der er mange af de
store kristelige Samfunds Antagelser, jeg ikke deler.“ --- En
Kvinde, der lever med Mand og Barn paa Landet i stadig Kamp
med Fattigdom, erklærer, at for hende er Religion Sandhed, og
ved Sandhed forstaar hun den Lov, der gælder for vort Væsen.
„Naar jeg siger Gud, mener jeg det Gode, mit Væsens Lov ...
Jeg tror, at Liv er Udvikling, Væxt, . . . og at vi have saa
megen Himmel og saa meget Helvede her i Livet, som vi selv
ville . . . . Det Gode er vort Ideal; det Onde er vor Lærer, som
skal faa os til at forstaa det Gode . . . . Mennesker kunne ikke
deles i gode og onde, men i udviklede og uudviklede . . . . Jeg
kender ingen religiøs Handling uden selve det at leve . . . .
Maaske var jeg nærmest ved religiøs Erfaring, da det første
Gang blev klart for mig, at jeg var fri, fri for Tro, Lære, Form
og Sentimentalitet, fri til at voxe, til at blive Alt, hvad der
ligger i min Natur . . . . Hvis De ikke kan finde nogen Religion
i dette, vil De neppe finde nogen Religion hos mig . . . . Jeg
er imod Læren om Kristi Efterfølgelse, fordi den strider mod
selvstændig Tænkning, og fordi sjælelig og aandelig Udvikling
beror paa selvstændig Tænkning. Vi kunde ligesaa godt lade
en Anden spise for os som tænke for os. Begejstring kan en
% --- p. 60 ---
\opage{60}Aand indgive en anden, men ikke Existens . . . . Kristendommen
har været en stor Faktor i Udviklingen af den menneskelige
Aand, men Dogmernes Kristendom er man nu næsten voxet
fra . . . . . . ..... Vi ere ikke ved at miste vor Religion, men vi ere
ved at finde os selv. Vi betvinge Frygten og lære at have Tillid
til den guddommelige Lov, der raader for Livet“.

Alle disse Exempler høre til den Gruppe, der mener at være
positivt religiøs, skønt det ofte, især for det sidste Exempels
Vedkommende, er tvivlsomt, med hvilken Ret. Men de ere
karakteristiske Udtryk for Trangen til personligt at gøre sig
Rede for, hvad man --- den Enkelte --- kan finde af Værdi for
sit Liv i den givne positive Religion. Gennem nogle Overgangstilfælde
af mere ubestemt Form komme vi i Leubas Række over
til afgjort ikke-religiøse Former, som --- hvor langt man saa
strækker Begrebet Religion --- ikke kunne komme ind under
det. Dels mangler der her religiøs Erfaring i enhver Betydning
af Ordet, dels er der bestemte intellektuelle og moralske Grunde,
som føre bort fra Religionen --- ganske vist fordi Religion og
Kirkevæsen antages at falde sammen. Saaledes bemærker en
Korrespondent: „Jeg benegter Kirkernes moralske Overlegenhed.
De tage aldrig Ledelsen af en moralsk Reform. De vare paa
den gale Side i Slavespørgsmaalet. De ere fejge, smaatskaarne
og trællesindede i vor Tids store industrielle Spørgsmaal.“
--- Leuba mener, at Antallet af ikke-religiøse Personer er langt
større, end man i Almindelighed tror. Naar der ikke er mange
Udtalelser fra saadanne i hans Samling, er det efter hans Mening
dels, fordi hans Spørgsmaalslister ikke have naat til dem,
dels, fordi de ikke altid ere klare over deres virkelige Standpunkt. ---


Leuba har ikke selv afsluttet Analysen af sit Materiale;
hans hidtil foreliggende Afhandlinger ere kun Resultater af en
første Behandling, og han ender egentligt med et ubesvaret
Spørgsmaal. I det religiøse Grundmotiv, de religiøse Tendenser
finder han Intet, der er særligt for Religionen, men særegne
Former for Selvopholdelses- og Selvudviklingstrangen. Denne
Opfattelse finder han afgjort bekræftet ved sine Korrespondenters
Udtalelser. Han afviser enhver Opfattelse af Religionen,
% --- p. 61 ---
\opage{61}hvorefter denne skal gaa ud paa Erkendelse, hvad enten denne
Erkendelse skal give en Aabenbaring, eller konstatere et Mysterium.
(The Monist. July 1901. p. 571---573). Men samtidigt med
at han i sin sidste Afhandling (Revue philos. Novbr. 1902. p. 487)
viser tilbage til disse Udtalelser, udtaler han: „Hvad er da det
Element, som adskiller det religiøse Liv fra det verdslige (ti der
er dog en Forskel mellem dem)? Det er det Middel eller de
Midler, Mennesket anvender for at opnaa sine Ønsker, og ikke
selve disse Ønsker, der adskille Religionen fra det øvrige Liv.
Psykologisk udtrykt: det er ikke Villien (i Schopenhauers Betydning
af Ordet), men Erkendelsen, der skiller dem. Naar man
indser, at Forskellen mellem de religiøse Mennesker og dem,
som ikke ere det (forsaavidt der er saadanne), ikke bestaar i
Stræben, Ønske og Tendens, altsaa ikke i selve Livets Kilder,
men kun i de Midler, der anvendes for at naa samme Maal,
fremtræder Religionen i et andet Lys end det, i hvilket den
sædvanligvis staar i Folks Øjne.“ --- En nærmere Forklaring af
disse Antydninger ville forhaabentligt Leubas følgende Arbejder
bringe. Men naar han bestemt afviser den intellektualistiske
Opfattelse, at det skulde være Religionens \textit{Formaal} at give Erkendelser
af nogensomhelst Art, og dog mener, at Forskellen
mellem Religionen og andre aandelige Livsformer \textit{beror paa} Erkendelse,
saa kan Meningen vel kun være den, at Forskellen
beror paa den Maade, paa hvilken Erkendelsen i Religionen
bruges som \textit{Middel.} Og Middel maa her betyde Fortolkningsmiddel,
altsaa Symbol. Jeg tror ogsaa, at Leuba af sit Materiale
vilde kunne uddrage et Resultat i denne Retning. Der er i al
Religion en Trang til at forme det Oplevede i Tanker. I den
konfessionelle Religion tilfredsstilles denne Trang umiddelbart
ved de overleverede Forestillinger; i Overgangsformerne søges
der efter andre Udtryk til at fortolke selve disse overleverede
Forestillingers Betydning for Livet. Til Grund ligger stedse en
Erfaring om Forholdet mellem Ideal og Virkelighed, en Erfaring,
der --- hvor Religion idethele opstaar --- gaar over til Tro og
Tillid overfor Magter, der antages at virke i Retning af et
harmonisk Forhold mellem Ideal og Virkelighed. Jeg finder
min Hypotese om Religionens Væsen som Tro paa Værdier%
% --- p. 62 ---
\opage{62}nes Bestaaen bekræftet ved Leubas og hans Kollegers Undersøgelser.


Derimod forekommer det mig, at Leuba lægger for liden
Vægt paa herskende eller overleverede Forestillingers Indflydelse.
Det er skæbnesvangert for en Følelses hele Karakter at blive
knyttet til en bestemt Kreds af Forestillinger; den faar herved
en Retning og et Indhold, der afgør dens Kvalitet, og naar
Forbindelsen mellem den og de Forestillinger, den hidtil har
været knyttet til, skal sønderrives, kan det ske, eller i hvert
Tilfælde true med, at Følelsen forsvinder med det Samme. Man
kan ganske vist med Leuba skelne mellem Formaal, som kunne
siges at være fælles for Religionen og andre aandelige Livsformer,
og Midler, der ere forskellige for hin og for disse; men
det maa da vel fastholdes, at selve det, at et Maal kan naas ved
forskellige Midler, tyder paa en Forskellighed i Maalet selv.
En ganske udvortes Sammenhæng mellem Maal og Midler er
der paa saadanne Omraader, som det her drejer sig om, ikke.
Det religiøse Problems Brod ligger netop paa dette Punkt.
Historien lærer jo ogsaa, hvor ondt Mennesker have ved at
finde sig i Kvalitetsforskelligheder indenfor Religionen, og
Striden om Grænserne for selve Begrebet Religion beror ligeledes
paa det Spørgsmaal, om vi ved en Kvalitetsforskel komme
udenfor Religionens hele Modalitet (for at bruge Sansepsykologiens
Terminologi). Det bliver herved af særlig Betydning,
at der gives Naturer, som kun formaa at finde Udtryk for
deres indre Liv i \textit{fælles} Symboler, og som, hvis de nødes til at
opgive dem, føle sig aldeles hjemløse.

Dog maa det gentages, at Leuba ikke endnu har fuldendt
Bearbejdelsen af sit Materiale. Det vil være af Interesse at
høre de Ord, hvormed han afslutter sine foreløbige Meddelelser.
„Om Forskellighederne mellem Religion og Ikke-Religion kun
ere overfladiske, eller om de ere fundamentale, og da i hvilken
Grad, --- om Opdragelsen i Hjemmet og andre Omstændigheder
i og for sig kunne forklare disse Forskelligheder, eller om de
have dybere Rødder i særegne, medfødte Anlæg, --- om den ene
er bestemt til at fortrænge den anden eller blot er et forberedende
Stadium for den eller ikke, --- det er et Spørgsmaal,
% --- p. 63 ---
\opage{63}som maa forbigaas i Taushed i denne foreløbige Afhandling.
Kun det ønsker jeg at betone her: selv mellem Individer af
samme Samfund og omtrent af samme Intelligens og Kultur
gives der Forskelligheder, som ere tilstrækkelige til at gøre
den Ene til et Menneske, for hvem Religion er Livets Alt, og
lade den Anden staa aldeles ligegyldig og utilgængelig overfor
den, selv om den fortolkes i videste Forstand.“ ---

\textsc{Lucien Arréat} har --- ad lignende Veje som de amerikanske
Psykologer --- søgt at finde forskellige Typer indenfor
den franske Katolicisme. Det mærkes strax, at Forholdene ere
andre. Forskellen mellem Katolicisme og Protestantisme hænger
utvivlsomt sammen med Raceforskelligheder. Medens man fra
anden Side (Charles Renouvier) har betragtet det som det
franske Folks Syndefald, at det ikke sluttede sig til Reformationen,
hævder Arréat, at denne stred imod den franske Races
Ejendommelighed. Protestantismen krænkede Skønhedssansen
ved sin Ødelæggelse af Billeder og Bygninger, og hvad den berøvede
Fantasien, gav den ikke Erstatning for paa Tankens
Omraade. Den protestantiske Individualisme fordrer af den
Enkelte en selvstændig Drøftelse og Forstaaelse, der strider
mod den latinske Races Utilbøjelighed til teologiske Diskussioner.
Katoliken ser Læren som fastsat en Gang for alle, og
hans Inderlighed bryder frem i det personlige Følelsesforhold,
den begejstrede Tilslutning til den faststaaende evige Sandhed,
som Kirken garanterer. Arréat har derfor heller ikke i sit
Materiale fundet nogen Antydning af den Overgang fra
Lydighedsfølelse til Frihedsfølelse i Gudsforholdet, som Leuba
mente at finde i Angloprotestantismen; den katolske Religiøsitet
er væsentligt Underkastelse. --- Efter sine Iagttagelser
beskriver Arréat fire forskellige Klasser indenfor den franske
Katolicisme.

Den første Klasse, der især findes i udpræget Form i Sydfrankrig,
indbefatter den største Mængde af Katoliker og har
Repræsentanter i alle Klasser af Samfundet. Fromheden bestaar
her i Overholdelse af de religiøse Skikke, der opfattes paa
en Maade, som nærmer sig til Afgudsdyrkelse, hvad der især
viser sig i den stærke Tilstrømning til de Steder, hvor der
% --- p. 64 ---
\opage{64}sker Mirakler. Her virker den Racekarakter, hvis Tilbøjelighed
til Overtro allerede Cæsar omtalte.

De sande Katoliker udgøre den anden Klasse. Gennem
Deltagelse i Kirkens Naademidler styrke de Tryghedsfølelsen
ved at være ombord paa St. Peters Skib, der midt i Tidernes
Uro føres af en ufejlbar Styrmand. Den katolske Religion prises
for sin Ophøjethed og sin betagende Skønhed, for sit Fuldkommenhedsideal
og for sin virksomme Menneskekærlighed.
Uden Trang til Mystik eller til selvstændig Tænkning ytrer der
sig her en aandelig Ligevægt, som dog undertiden er forbunden
med en udtrykkelig Forvisning om, at det kun er Kirkens
vældige Myndighed, der holder Tvivl og Fare borte.

Til en tredie Klasse regnes de, der støtte deres kirkelige
Tro ved mere eller mindre bevidst Tænkning. Mange Gejstlige
høre herhen, og til denne Klasse mener Arréat ogsaa at kunne
henregne Mange, som vægrede sig ved at udtale sig. Man
protesterer her imod, at Religionen skulde være en Følelsessag:
den grunder sig paa rationel Erkendelse, skønt Naaden
bringer den sidste Afgørelse. Den katolske Lære staar som en
Helhed, af hvilken Intet kan udskilles, og som umiddelbart
betager Sjælen. Dog udtales det undertiden, at Kirken kun
er Troens ydre og forgængelige Klædebon, der vil kunne omdannes
til nye Former; det egentligt Religiøse sættes i Trangen
til det Uendelige. Ved saadanne Udtalelser gøres Overgangen
til den sidste Klasse.

Hos mange Mænd og Kvinder, der erklære sig selv for
gode Katoliker, kan en vis Opløsningstilstand spores. De vælge
f. Ex. mellem Kirkens Lærdomme og Skikke og holde sig til,
hvad de have Brug for. Ofte er det kun en udtrykkelig Villiesakt,
der holder dem tilbage i Kirken; man betoner de sociale
Motiver, Lejligheden til at udøve Menneskekærlighed f. Eks.,
som det, der knytter til Kirken. Eller man opfatter det Religiøse
som ideal Stræben, hylder en praktisk Idealisme uden
Metafysik og benytter kun Kirkens Kultus som Stemningsmiddel.
Man udtaler den inderligste Sympati med de Troende,
fordi de ikke, som mange Fritænkere, give Afkald paa Religionens
ophøjede Poesi eller ere blaserede overfor de vældige
% --- p. 65 ---
\opage{65}Problemer, Ordet Gud peger hen paa. Et sidste Skridt gøres,
naar man gaar over til Positivismen eller ad anden Vej søger
et helt nyt Udtryk for sin Livsanskuelse. Dette er dog usædvanligt.
Enten vender man tilbage til Barndommens Tro efter
en Uroens Tid, eller man synker i aandelig Henseende ned paa
et lavere Trin.

Det er --- udtaler Arréat til Slutning --- umuligt for Øjeblikket
at sige, hvorhen den nuværende Krise paa det religiøse
Omraade vil føre, eller hvilken Energi den religiøse Følelse vil
kunne besidde, naar den ikke støttes til Dogmetro eller poetiske
Myter: „Jeg ser, hvad man vil udrydde af det moderne Sjæleliv;
jeg ser ikke, hvad der ydes til dets Berigelse.“

\underafsnit{d. Poesi og Religion.}{d. Poesi og Religion}

\textsc{George Santayana} har i sit ovenfor nævnte Skrift drøftet
Forholdet mellem Poesi, Religion og Videnskab paa en Maade,
der giver ham en Plads i Religionspsykologernes Række. Det
er for ham Fantasien, den billeddannende Evne, der ligger til
Grund for alle hine tre aandelige Livsytringer. De Billeder,
vi kunne bruge i Praksis, og ved hvis Hjælp vi kunne forudsige
fremtidige Erfaringer, kalde vi Forstandsforestillinger,
medens Billeder, der ikke kunne udnyttes paa denne Maade,
vedblive at hedde Fantasiforestillinger. Forstand er en Fantasi,
der gør Erfaringsvidenskab mulig. Vi vælge under Erfaringens
Indflydelse mellem de uvilkaarligt opdukkende Billeder. Men
ved de Fantasibilleder, der ikke kunne gaa over til Forstandsforestillinger,
maa igen skelnes mellem to Arter. Naar Fantasibillederne,
skønt de ikke udtrykke empirisk Realitet, dog staa
som Udtryk for Livsværdier, Idealer og Formaal, faa de Religionens
Karakter, medens de staa som rent poetiske Former,
naar de ingen saadan Livsbetydning have. Religion og Poesi
ere i deres inderste Væsen beslægtede, ja identiske, og blot forskellige
ved deres forskellige Forhold til Livet. Poesien gaar
paa sit Højdepunkt over til Religion, naar den udtrykker vore
Erfaringers Værdi og giver store, ledende og begejstrende Antecipationer.
Omvendt gaar Religionen over til Poesi, naar den
opgiver Forudsætningen om at indeholde metafysiske Sandheder.
% --- p. 66 ---
\opage{66}Det gælder, for ret at forstaa dette, at skelne mellem Værdiernes
og Kendsgerningernes Verden: Værdi beror ikke paa selve det
Givne, men paa dets Betydning.

Vor Villie er uendelig, men vor Erfaring er begrænset.
Poesien fører os ud over denne Begrænsning og gør derved en
højere Udvikling og Øvelse af Villien mulig, idet dennes højeste
Maal nu ved Poesiens Hjælp kan fremtræde med større Klarhed,
end den blotte Erfaring tillader. En Stræben, en Lidenskab
forstaa vi kun ret, naar vi se det Ophøjede i dens hele og fulde
Udfoldelse. Men fordi Poesien bliver Religion, hører den ikke
op at være Poesi. Paa dette Punkt forsynde ikke blot de Ortodoxe,
men ogsaa deres Kritikere sig. Allerede den græske
Filosofis Kritik af Polyteismen var uretfærdig, idet den ikke tog
Hensyn til Mytologiens nøje Sammenhæng med Livet. En stor
Religion er et Produkt af Menneskeaanden i dens Helhed, og i
Forhold til det staar Kritiken, der ledes af den specielle, intellektuelle
Interesse, som begrænset og smaatskaaren. I Kristendommen
traadte historisk Idealisme i Stedet for den antike
Religions Naturalisme; den Harmoni, den gamle Religion havde
søgt at fremstille i Rummet og i Naturen, søgte den nye at
fremstille i Tiden og i Historien. Sit Fodfæste i Historien fik
Kristendommen derved, at den havde sin Rod i en national Tro,
der havde formet sig under Jødefolkets Prøvelser; men ud herfra
formede den et stort poetisk Udtryk for Erfaringer, Menneskene
kendte fra deres Sjæls Inderste. Det var Menneskehedens
store Selvbiografi, der her blev gjort til Verdenslivets Midtpunkt.
Den Betydning, en Religion har og bevarer, er ikke nødvendigvis
den, som den havde for dens Grundlæggere; disse betages
saaledes af deres Syner, at de betragte disses Indhold som ydre
Kendsgerninger. Religion er ikke forskellig fra Overtro i psykologisk
Henseende; ikke paa Oprindelsen, men paa Værdien
beror Forskellen. Selv om den etiske Værdi ikke var Grunden
til Religionernes Tilblivelse, saa er den dog den sidste Maalestok
for deres Betydning. ---

Hele denne Tankegang, som Santayana udfører i det Enkelte
gennem interessante Karakteristiker af religiøse og poetiske Værker,
faar en særegen Interesse, naar man sammenholder den med
% --- p. 67 ---
\opage{67}nogle af Arréats, Leubas og Coes Iagttagelser, og naar man tillige
mindes de idealistiske eller mystiske Tendenser, der røre
sig i den nyeste Tids Literatur. Hvad jeg savner hos Santayana
er en mere indtrængende psykologisk Analyse af Forholdet
mellem Erfaring og Tro; men iøvrigt har det for mig været af
stor Interesse at finde en med min egen beslægtet Betragtningsmaade
i et Skrift, der er udkommet omtrent samtidigt med mit
eget religionsfilosofiske Forsøg. Jeg tillader mig til Sammenligning
at anføre en enkelt Udtalelse fra min Bog (Religionsfilosofi
p. 334 f.): „Maaske forholder det sig saaledes, at Poesi
er et fuldkomnere Udtryk for det Højeste, end noget videnskabeligt
Begreb vilde kunne være. Ved Poesi maa da her ikke
tænkes paa løse Stemninger og Fantasier, men paa den uvilkaarlige
og livfulde Form, det virkeligt Oplevede iklæder sig i
stærkt bevægede Øjeblikke . . . . Det er en Poesi, der hverken
staar i Modsætning til Villie eller til Tanke, selv om den særligt
bliver til, hvor Tanke og Villie støde paa Skranker ... Al høj
Kunst forudsætter en saadan Livspoesi og er dens Omsætning i
klare Former.“

\underafsnit{e. Sammenligning mellem forskellige Typerækker.}{e. Sammenligning mellem forskellige Typerækker}

I min Religionsfilosofi har jeg, især støttet til et biografisk
og historisk Materiale, fremstillet en Række typiske Forskelligheder
i religiøs Tro (p. 111---120) og, i en anden Sammenhæng,
en Række Forskelligheder i Henseende til de oprindelige Dispositioner,
der betinge det religiøse Livs Ejendommelighed hos
de Enkelte (p. 251---256). Sammenligner man disse Rækker med
de Typer, som de amerikanske Religionspsykologer have opstillet,
vil man især finde én Modsætning, som i alle Fremstillinger
spiller en Hovedrolle, nemlig den, jeg har beskrevet som
en Modsætning mellem Tro som Koncentration og Hvile efter
Udligning af Disharmonier og Tro som højeste Form for Livsudfoldelse.
Hertil svarer indenfor Dispositionerne Modsætningen
mellem disharmoniske og expansive Naturer. Baade hos Leuba,
Starbuck og Coe er dette den mest fremtrædende Modsætning.
Men det viste sig tillige, at de lægge ganske særlig Vægt paa
den expansive Type og ere tilbøjelige til at betragte den dis%
% --- p. 68 ---
\opage{68}harmoniske eller diskontinuerlige Type som et specielt Tilfælde
af den. Der er naturligvis altid nogen Modstand at overvinde,
nogle Kriser at gennemgaa, og det er kun, naar Modstand
og Kriser blive særligt stærke, at vi efter disse Forskeres
Mening komme over fra den expansive til den disharmoniske
Type. Om Grunden til, at den expansive Type saa hyppigt
forekommer i deres Materiale, har jeg ovenfor udtalt en Formodning.
Man kunde maaske føje til, at de nævnte Forfattere (og
det Samme gælder aabenbart Santayana) selv personligt afgjort
synes at høre til den expansive Type, forsaavidt man kan lære
deres psykiske Habitus at kende af den Maade, hvorpaa de opfatte
og vurdere Andres. Jeg for mit Vedkommende vilde være
tilbøjelig til at anse Forskellen mellem de to Typer for større
og mere fundamental, end hine Forfattere antage. Maaske ligger
der i denne Modsætning Spiren til en aandelig Kamp, der atter
og atter vil komme til at staa mellem aandelige Livsretninger.
Dybest inde forstaa de tilsidst ikke hinanden. I det danske
Aandsliv have vi et klassisk Exempel paa denne Modsætning i
Forholdet mellem Grundtvigianismen og den indre Mission.

Skulde jeg reducere min Typerække, vilde denne Modsætning
dog blive staaende. Jeg har idethele syv Typer, nemlig
foruden de to allerede nævnte (Troen som 1) Hvile og som
2) Udfoldelse): 3) Kontemplation, 4) tillidsfuld Voven, 5) Resignation,
6) krampagtig Religiøsitet og 7) Kulsviertro (fides implicita).
Af disse vilde 3, 4 og 7 gaa ind under 2, medens 5 og 6
vilde høre ind under 1. Hvad Forskellighederne i Dispositioner
angaar, opstiller jeg paa Følelses- og Villieslivets Omraade
følgende Række af Modsætninger: mellem de Disharmoniske
og de Expansive, mellem de Passive og de Aktive, mellem de
Affektive og de Kontinuerlige og mellem de Idiopatiske og de
Sympatiske, --- og paa Erkendelsens Omraade Modsætningen
mellem analytiske og syntetiske Naturer og mellem diskursive
og intuitive Naturer. Det er let at se, at der er en Analogi
mellem alle disse Modsætningsforhold, og at de alle kunne indordnes
under det første, saa at det bliver Modsætningen mellem
de Disharmoniske og de Expansive, der ogsaa fra denne Side
set kommer til at staa i første Linie.

% --- p. 69 ---
\opage{69}I den Overbevisning, at der for Fremtiden i langt højere
Grad end hidtil ved Forstaaelse og Bedømmelse af religiøse Fænomener
maa lægges Vægt paa de individuelle Forskelligheder,
stemmer jeg ganske overens med mine amerikanske Kolleger.

\whiteline

\afsnit{4. Komparativ Religionspsykologi.}{4. Komparativ Religionspsykologi}

Et sammenfattende Arbejde paa Religionspsykologiens Omraade
foreligger i den amerikanske Filosof \textsc{William James’} Bog
\textit{The varieties of religious experience.} A study in human nature.
London, New-York and Bombay. (1902). James hører til vor
Tids mest fremragende Tænkere. Han forener omfattende Kundskab
med stor Iagttagelsesevne, skarp Kritik med idealistisk
Begejstring, Fordomsfrihed med alvorlig Overbevisning. Det
betydeligste af hans Skrifter er \textit{Principles of psychology} (1890),
der paa flere Punkter af Psykologien gør nye Synspunkter
gældende og overalt drøfter Spørgsmaalene paa en aandfuld og
vækkende Maade. I en Samling af Afhandlinger, som er udkommen
under Titelen \textit{The will to believe} and other essays in
popular philosophy (1897), har han i noget friere Form udtalt
sin ejendommelige Forening af Erfaringsfilosofi og udpræget
idealistisk Livsanskuelse, og allerede her kommer han ind paa
religionsfilosofiske Spørgsmaal. Det nye Værk er ganske helliget
det religiøse Problem, særligt den rent psykologiske Belysning
af det, dog saaledes, at ogsaa de andre Synspunkter, fra hvilke
Problemet maa behandles, komme frem. James benytter de
yngre amerikanske Forfatteres Specialarbejder, men støtter sig
desuden til en omfattende Literatur af biografisk, især selvbiografisk
Art. Han ser klart, at Emnet er for mangfoldigt til
at kunne udtømmes i en enkelt Fremstilling, og derfor --- tillige
for at slippe for en Ordstrid om, hvad Religion er --- erklærer
han, at han til Brug ved denne sin Undersøgelse vil
holde sig til en bestemt Betydning af Ordet Religion eller til
en bestemt Side ved Religionen, og han definerer saa Religion
som: individuelle Menneskers Følelse, Handling og Erfaring,
forsaavidt disse Mennesker i deres Ensomhed opfatte sig selv
som staaende i Forhold til det Guddommelige, hvorledes de
% --- p. 70 ---
\opage{70}end opfatte dette. Denne personlige Religion stiller han i
Modsætning til den institutionelle Religion med dens mere eller
mindre udviklede Teologi eller Mytologi og dens kirkelige Organisation.
Han hævder, at \textit{institutional religion} (med Alt, hvad
dertil hører) stedse er afledet i Forhold til \textit{personal religion;} det
følger allerede af, at den bygger paa Overlevering. Hvis man
vilde holde sig til den institutionelle Religion alene, vilde
Undersøgelsen blot komme til at angaa Gudsdyrkelse, Ofring,
Ceremoniel og kirkelig Organisation, og man kunde da definere
Religion som den Kunst at vinde Gudernes Gunst. I den personlige
Religion er det derimod Menneskets egne indre Tilstande,
dets Samvittighed, dets Værdi, dets Hjælpeløshed eller
dets Ufuldkommenhed, Vægten ligger paa. Holder man sig til
denne Side af Religionen, faar Begrebet Religion et større Omfang,
end naar man anser det Institutionelle for Hovedsagen.

Den personlige Religion er den Maade, hvorpaa Mennesket
reagerer mod Livet som Helhed (a man’s total reaction upon
life), i Modsætning til de øjeblikkelige og specielle Reaktioner.
De Erfaringer, vi have gjort om Livet og om Verden, virke paa
vort individuelle Temperament og gøre os enten energiske eller
ligegyldige, enten fromme eller spottende, enten nedstemte eller
begejstrede overfor Livet som Helhed. Vor „Reaktion“, der ofte
foregaar uvilkaarligt og halvt ubevidst, udtrykker sig i det fuldstændigste
Svar, vi kunne give paa det Spørgsmaal: Hvad er
Karakteren af denne Tilværelse, vi høre til? --- Dog vil James
kun kalde en saadan Reaktion religiøs, naar den bæres af en
alvorlig og ophøjet Stemning, en Stemning, der paa sit Højdepunkt
fører til Opgivelse af egne Ønsker, til villig Underkastelse
under Virkelighedens Smerte, fordi man føler sig baaren af en
Kraft, der er større end den, man selv raader over. Der er i al
udpræget religiøs Følelse et tragisk Element, som beror paa, at
en højere Lykkefølelse har fortrængt en lavere. Den religiøse
Følelse indeholder eller forudsætter saa mange Modsætninger, at
den er en af de rigeste menneskelige Følelser. Verden er rigere
ved, at der er en Djævel, naar der er en St. Michael, som kan
sætte sin Fod paa hans Nakke. Dette Særkende træder naturligvis
mest frem i de extreme Former for Religion, Former, der
% --- p. 71 ---
\opage{71}for den jevne, dagligdags Opfattelse staa som abnorme; men i
alle Former for religiøs Bevidsthed vil det til en vis Grad forekomme.. % PRINTED AS IS: the doubled stop


Denne Beskrivelse angiver den almindelige Ramme, indenfor
hvilken James bevæger sig i sin Undersøgelse. Til nærmere
Karakteristik af hans Arbejde skal jeg først vise, hvorledes de
tre Synspunkter, der --- som omtalt i Indledningen --- maa komme
frem ved enhver omfattende Drøftelse af det religiøse Problem,
ogsaa fremtræde hos James, skønt han hovedsageligt har holdt
sig til den rent psykologiske Betragtningsmaade. Dernæst skal
jeg vise, hvorledes han opfatter den religiøse Erfarings Væsen,
hvilke Typer for Religion han finder, og til hvilken Grundtanke
Religionen efter hans Opfattelse kan føres tilbage.

\underafsnit{a. De tre Synspunkter.}{a. De tre Synspunkter}

James betegner sit Standpunkt, baade i almindelig filosofisk
Henseende og i religionsfilosofisk Henseende, som Erfaringens
Standpunkt. Han vil være gennemført Empiriker. Alle intellektuelle
Operationer --- hvad enten de bestaa i Konstruktion, i
Sammenligning eller i Kritik --- forudsætte umiddelbar Erfaring.
Han slutter sig til den gamle engelske Skole, den, der begyndte
med John Locke, og som stedse krævede bestemt Regnskab for
de Erfaringer, paa hvilke vi bygge vore Begreber. Denne Skole
bør efter hans Mening langt snarere end Kant have Æren for
den kritiske Metode i Filosofien, den eneste Metode, der gør
Filosofi til et Studium, som kan drives af alvorlige Mennesker.
Men det drejer sig ikke blot om de Erfaringer, vi have havt,
ogsaa om dem, vi maa vente eller konsekvent selv fremkalde,
hvis vore Tanker skulle være gyldige. Vore Tankers Værdi
beror paa de Slutninger, vi kunne drage af dem. Hvis alle
Sætninger, vi opstille, vare praktisk ligegyldige, d. v. s. hvis der
slet Intet fulgte af dem, kunde vi saa skelne mellem sande og
falske Tanker? Forgæves har Dogmatismen søgt efter Kriterier
for Sandhed, der kunde fritage for at appellere til Fremtiden. I
Modsætning til Dogmatismen staar her Pragmatismen, som
prøver hver Sætning ved at forfølge dens Konsekvenser. Denne
Opfattelse (i hvilken James slutter sig til den amerikanske Tæn%
% --- p. 72 ---
\opage{72}ker \textsc{Charles Peirce}, fra hvem ogsaa Navnet Pragmatisme stammer)
minder om den biologisk-økonomiske Erkendelsesteori, der i
Tyskland er bleven udviklet af Richard Avenarius og Ernst Mach.

James’ empiriske Standpunkt fører ham til Skepsis overfor
spekulative og teologiske Synsmaader. Han hævder særligt Mangfoldigheden
i Tilværelsen og er kritisk stemt overfor Monismen
i dennes forskellige Former. Pluralismen, Mangfoldighedslæren
vil, tror han, mere og mere komme til Ære --- baade i Videnskaben
og i Religionen. Ud fra Erfaringen kan man intet Universelt,
ingen absolut og altomfattende Enhed naa. Og forskellige
Tydninger kunde gælde for forskellige Egne af Tilværelsen.
Religiøst set blive her især Erfaringerne om det Onde af Betydning:
har Jobs Bog ikke en Gang for alle godtgjort Umuligheden
af at løse det Problem, disse Erfaringer føre til at stille?
Her vil den monistiske Hypotese stedse frembyde store Vanskeligheder,
og vilde det saa ikke være den simpleste Udvej at antage,
at Verden fra først af har indbefattet Elementer, der ikke
vare i Harmoni med Helheden? Sædvanlige Mennesker have
stedse været mere eller mindre polyteistiske, og det er da ogsaa
klart, at man ikke af nogen Erfaring kan udlede Guds Uendelighed.
Hvad vi kunne erfare, kan i det Højeste være det, at
vi staa i indre Sammenhæng med Noget, der er større end vi
selv; men at dette er et Uendeligt, kan ikke godtgøres.

Hvad Tænkningen kan udrette paa det religiøse Omraade,
kan kun være det: at ordne, bestemme og fortolke Kendsgerningerne,
ikke at frembringe dem. Filosofien er her stedse sekundær.
Dog kan den derfor godt være til Nytte. Ti den religiøse
Bevidsthed udtrykker stedse sine Erfaringer i Former, der
stamme fra den intellektuelle Overlevering, under hvis Indflydelse
den staar. Derfor kan kun kritisk Sammenligning finde det virkeligt
Umiddelbare og Væsentlige og udsondre det fra det Lokale
og Tilfældige. Tillige er det Filosofiens Opgave at udskille
Antagelser, der ere inkonsekvente eller afgjort stride mod videnskabelig
Erfaring. Hvad der bestaar gennem denne Skærsild,
kan saa betragtes som Hypoteser, hvis Berettigelse og Betydning
kunne blive Genstand for nærmere Drøftelse, efterat de ere
blevne reducerede til deres simpleste Former.

% --- p. 73 ---
\opage{73}Saavidt den erkendelsesteoretiske Betragtning. Saa energisk
James gør den gældende, saa er det dog den psykologiske Betragtningsmaade,
han har størst Interesse for, og som giver hans
Arbejde dets største Betydning.

Religionspsykologien yder den almindelige Psykologi vigtige
Bidrag. Der er mange psykologiske Fænomener, der særligt
godt kunne studeres paa det religiøse Livs Omraade. James
fremhæver saaledes, at de religiøse Fænomener godtgøre, hvor
ringe en Del af vort aandelige Liv der kan forklares klart og
tydeligt. Bevidstheden fortoner sig her gennem mange Grader
over i det Ubevidste, eller, som James hellere vil sige, i det
Underbevidste. Der rører sig umiddelbare Tilskyndelser og Anelser,
som ofte, uden at vi selv mærke det, give os de første
Præmisser, paa hvilke klart bevidste Tanker grunde sig. Bevidste
Argumenter bevæge sig ofte kun paa Overfladen af vort Væsen,
og en uvilkaarlig og umiddelbar Forvisning er det Dybeste i os
(the deep thing in us). Fremdeles lærer Studiet af det religiøse
Liv os store Forskelligheder mellem Mennesker at kende, baade
i Henseende til Arten af Følelseslivet og i Henseende til Graden
af Modtagelighed, --- Forskelligheder, vi ellers neppe vilde blive
saa opmærksomme paa. Vi lære, hvorledes Udviklingen henimod
et og samme Maal kan foregaa højst forskelligt, snart kontinuerligt,
snart springende. Allerede rent teoretisk eller naturhistorisk
set har et psykologisk Studium af det religiøse Liv sin Betydning.

Den psykologiske Betragtningsmaade bliver saa igen af Betydning
for Forstaaelsen af det religiøse Problem. De religiøse
Processer ses som specielle Former for Processer, som ogsaa
kendes paa andre psykiske Omraader, og adlyde Love, der ogsaa
ellers kendes. James viser her hen til Starbucks Undersøgelser
af Omvendelsen i dens Sammenhæng med den organiske og
psykiske Udvikling i Overgangsaarene. Ganske særlig Vægt
lægger James paa det, der bevæger sig under Bevidsthedens
„Tærskel“ eller udenfor dens „Margin“. Han finder det meget
vanskeligt her at drage en bestemt Grænse og forkaster den
atomistiske Psykologi, for hvem Bevidstheden er sammensat af
enkelte, hvert for sig faste og bestemte Elementer. I hvert
Tilfælde gøre de subliminale eller ultramarginale Indflydelser sig
% --- p. 74 ---
\opage{74}stadigt gældende, og James er for sit Vedkommende tilbøjelig
til at betragte disse Indflydelser som de Veje, ad hvilke en
højere Tingenes Orden virker i vort Indre. Disse indre Tilkendegivelser
udgøre vor nærmeste, vor egentlige Verden; saa inderligt
tilhøre vi ikke den ydre Verden. Her finder James et
Enhedspunkt mellem Religion og Psykologi, skønt han erklærer
enhver nærmere Bestemmelse deraf for en Fortolkning, der ikke
længere er rent psykologisk. Fortolkningen er heller ikke det
Afgørende; den er stedse afledet, og en og samme Oplevelse kan
være Genstand for forskellige religiøse Fortolkninger\textsuperscript{12}). Hvad de
fleste Mennesker paa dette Punkt især mangle, er Kritik og
Forsigtighed, ikke Tro (faith). De ere for tilbøjelige til at lade
dogmatisk Antagelse (belief) følge hovedkuls efter enhver livfuld
Forestilling, især naar denne støttes af uvilkaarlig Interesse. ---

Det tredie Synspunkt, det etiske, kommer frem ved Værdsættelsen
af de religiøse Fænomener. Ligesaa bestemt, som James
skelner mellem Oplevelse og Fortolkning, ligesaa bestemt skelner
han mellem Beskrivelse eller Forklaring af det Oplevede og
Vurdering af det. Selv det, der fra et rent medicinsk Synspunkt
maa kaldes abnormt, kan ved sit Indhold have stor Værdi. En
Tilstands Værdi beror ikke paa, hvorledes den er bleven til.
Naar vi kalde visse Tilstande højere end andre, er det ikke paa
Grund af, hvad vi mene at vide om deres organiske Betingelser.
Hvis der er noget Værdifuldt i Sta Teresas Tanker, saa er det
ligegyldigt, om hun var hysterisk eller ikke. Vi vurdere en
Kraft efter dens Virkning, ikke efter dens Udspring.

Ved Vurderingen kan intet spekulativt eller teologisk System
lægges til Grund. Ogsaa her maa vi være Empirikere. Vi bedømme
de religiøse Fænomener efter deres Frugter, og det have
Menneskene egentligt bestandigt gjort. En Gudsdyrkelse hører
op, naar den ingen Virkning øver paa Sindet, og naar den i sin
hele Karakter kommer i Strid med Noget, hvis Værdi vi have
erfaret saa grundigt, at vi ikke kunne fornegte det. De Guder,
Menneskene holde fast ved, ere de, som de kunne bruge, og hvis
Bud bekræfte de Fordringer, de stille til sig selv og til Andre.
Vi bruge stedse menneskelige Maalestokke (human standards).
Vi spørge, hvorvidt det religiøse Liv godtgør sig som en ideal
% --- p. 75 ---
\opage{75}Form for menneskeligt Liv. Det er naturligvis en Maalestok,
der varierer med Tiderne, men anden Maalestok have vi nu
engang ikke. Kun maa Prøvelsen atter og atter gentages, da
den maaske under nye Forhold kunde føre til nye Resultater.
Og Vurderingen vil ofte falde forskelligt ud for forskellige Individer.
Naar psykologisk Undersøgelse godtgør, at der er store
Forskelligheder mellem Mennesker i deres inderste Natur, maa
man ogsaa antage, at de trænge til forskellig aandelig Næring,
saa at de komme til at staa forskelligt i religiøs Henseende.

Overfor Religionen idethele staar James afgjort sympatisk.
De bedste Frugter af religiøs Erfaring ere efter hans Overbevisning
de bedste Ting, Historien har at fremvise. Han føler
her det indre Liv røre sig med Alvor og Energi, virke med en
Inderlighed og en Koncentration, i Modsætning til Alt, hvad
der kan hæmme og splitte. Det er, som man var hævet op i
en højere og renere Atmosfære. --- James henviser paa dette
Punkt til Sainte Beuve’s Vurdering af de religiøse Personligheder,
han har studeret i sit berømte Værk om Port-Royal. ---

Vi finde saaledes hos James de tre Hovedsynspunkter ved
filosofisk Drøftelse af religiøse Spørgsmaal. Han bebuder, at
han i et senere Værk vil imødegaa de Forsøg, der endnu i vore
Dage ere gjorte paa en mere spekulativ Behandling af Religionsfilosofien,
--- Forsøg, der mærkværdigt nok forekomme hyppigere
i England og Amerika end paa Kontinentet. Foreløbigt staar hans
Arbejde som et betydningsfuldt Indlæg til Gunst for, at det
netop er hine tre Synspunkter, der ere at anlægge i Religionsfilosofien.
--- Desværre kan jeg i min Karakteristik af James
Standpunkt ikke lade den Aabenhed, Friskhed og sproglige
Kunst, han raader over, komme til deres fulde Ret.

\underafsnit{b. Den religiøse Erfarings Natur.}{b. Den religiøse Erfarings Natur}

Religiøs Følelse er ifølge James ikke nogen speciel eller
elementær Følelse. Enhver Følelse kan under visse Betingelser
faa en religiøs Karakter, nemlig naar den medvirker ved det
Opgør af de menneskelige Livserfaringer, James kalder „den
totale Reaktion overfor Livet“. En noget nærmere Bestemmelse
af, hvad Religion er, giver den Udtalelse, at Religion er
% --- p. 76 ---
\opage{76}en af de Veje, ad hvilke der kan komme Enhed i vort Liv.
Med Enheden følger en Følelse af Fred og Lykke, som overgaar
Alt, hvad der ellers kan erfares af den Art. Denne Enhed
kan naas ved pludseligt Gennembrud eller gennem en lang Udvikling;
her lægge de individuelle Forskelligheder sig for Dagen.
Men Hovedsagen er, at Individet, der før har følt sig delt og
hæmmet, lavt og ulykkeligt, nu har Følelsen af Harmoni, Frihed
og Opløftelse. Tanker, der før stode som ideale Muligheder
ude i Sjælens Periferi, ere nu blevne Sjælens centrale Ejendom.
Det er, som om der er tilstrømmet Individet ny Energi fra
Kilder, der ligge udenfor dets Bevidsthed, og som det ikke
vilde kunne udløse direkte ved bevidst Anstrengelse. Denne
underordnede Betydning af det Bevidste og Vilkaarlige (the
hopeless inferiority of voluntary to instinctive action) i Forhold
til det Ubevidste og Uvilkaarlige er et karakteristisk Træk i al
Religiøsitet. Sygelige Former opstaa, naar der mangler intellektuel
Kraft og Interesse til at forme og tolke saadanne
Erfaringer.

Religiøs Erfaring har en mystisk Karakter, dels paa Grund
af Uudsigeligheden, idet Enhver her kun indirekte kan forstaa
de Andre, dels paa Grund af Umiddelbarheden, idet den diskursive
Tænken ingen Rolle spiller, dels paa Grund af Tilstandenes
korte Varighed og Vanskeligheden ved at genkalde
sig dem, dels endeligt paa Grund af det passive Forhold, Individet
staar i til disse Erfaringer. Det ikke-religiøse Bevidsthedsliv
kender lignende Tilstande. Den religiøse Karakter
fremkommer, som allerede bemærket, derved, at Menneskets
Forhold til Livet som Helhed gøres op i saadanne Tilstande;
derved faa disse en Enhed og Energi, der ellers mangler. Den
særligt saakaldte Mystik er kun en bestemt og ofte extrem Form
for, hvad der fremtræder i al religiøs Erfaring.

\underafsnit{c. De vigtigste Typer.}{c. De vigtigste Typer}

James erklærer, at det vigtigste Resultat, han er kommen
til ved sine Undersøgelser af den religiøse Bevidsthed, er de
store Forskelligheder, den frembyder. En betydningsfuld Del
% --- p. 77 ---
\opage{77}af hans Arbejde gaar derfor naturligt ud paa at føre disse Forskelligheder
tilbage til visse Typer.

Der er særligt én Modsætning, han meget stærkt fremhæver,
nemlig den mellem det sunde og det syge Sind, eller
--- for at bruge et Udtryk, han har fra Francis Newman --- mellem
de én Gang Fødte og de to Gange Fødte (the once-born and the
twice-born). De syge Sjæle behøve en Forandring, en Krise, en
Omvendelse for at naa til Enhed og Fred, medens der er sunde
Sjæle, som ved umiddelbar Udfoldelse og Hengivelse naa til Maalet.
De Sundes Religion betegner James som Naturalisme, de Syges
som Salvationisme. Men indenfor hver af disse Grupper kunne
igen højst forskellige Personligheder og Retninger fremtræde.

De sunde Sjæles Religiøsitet er en Entusiasme, der fastholder
den lyse Side ved Tilværelsen og maaske endog erklærer
de mørke Sider for Illusioner. Selv om de ikke mangle Blik
for det Hæmmende og Formørkende, ere de dog forvissede om,
at det kan overvindes; saa stor er deres umiddelbare Tillid til
den indre Energi og til den Harmoni, der raader i Universet.
Denne Retning fremtræder mere i Græciteten end i Kristendommen,
mere i Katolicismen end i Protestantismen, mere hos
de latinske Racer end hos de germanske. I nyere Tid fremtræder
den hos Spinoza og Rousseau, i den liberale Retning
indenfor Kristendommen, især i Unitarismen (Emerson, Parker).
I Amerika træder den i skarp Modsætning til Metodismen og
Revivalbevægelsen. I nyeste Tid er den saakaldte mind-cure
movement et karakteristisk Exempel. Den helbreder baade
Sjæl og Legeme ved stadigt at indskyde og fastholde den Tanke,
at Smerten er en Illusion, og at Alt i Virkeligheden er lyst og
herligt. Frygt og Bekymring banlyses. Den virker gennem
Suggestioner, ikke mindst Selvsuggestioner. Ved stadigt at
raabe Hurra for Universet bestyrker man sig i Troen paa, at
Universet fortjener et Hurra.

Den syge Sjæl ligger under for Smerten og Disharmonien.
Bevidsthedens Tærskel overfor den Art Erfaringer er en Gang
for alle sænket, medens den hos den sunde Sjæl er hævet.
Verden staar som tom og forfængelig; der næres Frygt for,
% --- p. 78 ---
\opage{78}hvad Livets Gang kan bringe; Sjælen føler sig syndig og fordærvet.
Tomhedsfølelse, Frygtfølelse og Syndfølelse ere de
tre Hovedformer for sjælelig Sygdom. Hos Tolstoi raader Tomhedsfølelsen,
hos Bunyan Frygtfølelsen og Syndfølelsen. Der
er her en oprindelig Disharmoni i Temperamentet, som kræver
en anden Fødsel. Og selv om denne naas, ville Naturer af
denne Art neppe naa til den overstrømmende Livsfølelse, som
de Sunde kunne have.

Der findes naturligvis i Virkeligheden en Mængde Overgangsformer
mellem de to Typer. Men i deres extreme Former
er Modsætningen mellem dem saa stor, at de neppe forstaa hinanden.
De ville hver for sig betragte deres Livsanskuelse som
den ene rette. Saaledes hævdede Wesley Omvendelse som eneste
Vej til Fred, medens Emerson tillagde den Natur, der kun behøvede
én Fødsel, højere Værdi end den, der maatte kæmpe
haardt med Djævlene i sit Indre. James selv mener, at de
Livserfaringer, de to Gange Fødte kunne bygge paa, ere dybere
og fuldstændigere end den første Types. Ingen Profet kan bringe
et afgørende Budskab til Menneskene, naar han ikke siger Ting,
der kunne klinge som Realiteter for Mænd som Bunyan og
Tolstoi. De fuldstændigste Religioner maa være de, der baade
kende og overvinde Livets mørkeste Elementer. Derpaa beror
Buddhismens og Kristendommens Overlegenhed. ---

I Forhold til denne store Modsætning ere de andre Modsætninger,
der kunne paavises, for James af mindre Betydning.
Saaledes Modsætningen mellem springende og kontinuerlig Udvikling,
--- Modsætningen mellem Naturer, hos hvem de Elementer,
der skulle overvindes, først maa have mistet deres
Livskraft, og andre Naturer, hos hvem det nye Liv kan trænge
igennem, før det gamle er helt udtømt, --- og Modsætningen
mellem en Udvikling, der kræver bevidst Anstrengelse (selv om
denne ikke direkte fører til Maalet), og en Udvikling, der
gennemgaaende forløber uvilkaarligt. ---

Lejlighedsvis giver James foruden denne almindelige Typelære
ogsaa Antydninger af sin egen rent individuelle Tro. Han
har i en længe raadende, nervøst slappet Tilstand søgt Hjælp
hos mind-cure Bevægelsen og mener at være bleven helbredet
% --- p. 79 ---
\opage{79}ved dens Hjælp, hvilket ikke hindrer ham i at stille sig kritisk
overfor den. Naar han udtaler, at Mystiken kun giver i extrem
Form, hvad al religiøs Erfaring indeholder, føjer han til, at
han selv ikke personligt kender mystiske Tilstande. Sin Tro
karakteriserer han ved, at skønt han ikke kan slutte sig til den
gængse Kristendom eller til den spekulative Teisme, er han
overbevist om, at der i de Tilstande, i hvilke Mennesket føler
sig i den inderligste Enhed med det Højeste, det kender, kommer
ny Kraft ind i Verden og sættes nye Udgangspunkter; forsaavidt
er hans Tro supranaturalistisk, antager guddommelige
Indvirkninger indenfor den naturlige Erfarings Omraade. Han
føjer til, at han meget vel véd, at den herskende Mening
i akademiske Kredse gaar i modsat Retning. --- Udødelighedstroen
spiller kun en underordnet Rolle for ham (skønt han i et
særligt lille Skrift har søgt at vise, at den ikke kan gendrives).
Han siger herom: Naar der blot „i al Evighed“ sørges for vore
Idealer, ser jeg ikke, hvorfor vi ikke skulde være villige til at
overgive Omsorgen for dem i andre Hænder end vore.

\underafsnit{d. Den religiøse Grundtanke.}{d. Den religiøse Grundtanke}

Skønt Religionen efter James har sin Grund i Følelseslivet
og dets Erfaringer, og skønt den fremtræder i højst forskellige
Former, er der dog en Grundtanke eller en Hypotese, som den
stedse hviler paa. Det er den, at der kan naas en Tilstand af
Harmoni, Enhed og Fred, selv om det maa ske gennem Kampe
og Kriser, og at dette Maal naas ved Tilstrømmen af Energi
fra en mere omfattende Tingenes Orden, der psykologisk set ytrer
sig i den ubevidste Natur i os. Hvorledes vi i det Enkelte
fortolke, hvad der her sker, er ikke det Væsentlige; her ville
stedse rent individuelle Suppleringer og Udførelser (over-beliefs)
komme til. Hovedsagen er, at for den religiøse Bevidsthed er
Tragedien, saa ofte den end melder sig paa Livets Vej, stedse
kun foreløbig og partiel. Skibbrud og Opløsning ere ikke Tilværelsens
sidste Ord. Al Religion hviler saaledes tilsidst paa
en kosmisk Hypotese. Derimod er Religionens Bestaaen ikke
afhængig af nogensomhelst speciel dogmatisk Antagelse. Naturligvis
--- hvis Religion gøres til Et med eller ganske afhængig
% --- p. 80 ---
\opage{80}af en animistisk Verdensopfattelse, saa vil dens Tid være
forbi, naar Animismen uddør; ligesaa, naar man gør den til
Magi eller Fetischisme. Det bliver for James let en Strid om
Ord, her føres. Det Afgørende for ham er, at de indre Erfaringer,
hvorpaa Religionen væsentligt bygges, ere de mest
konkrete og umiddelbare Erfaringer, vi have. I Forhold til
dem har Alt, hvad vi kunne erfare om de ydre Verdensforhold,
en abstrakt og generel Karakter. Der er ingen Grund til at
antage, at saadanne Erfaringers Tid skulde være forbi, selv om
de maa udskilles fra Forbindelsen med Elementer, de hidtil
have været tilbøjelige til at smelte sammen med.

\underafsnit{e. Bemærkninger om James’ Religionsfilosofi.}{e. Bemærkninger om James’ Religionsfilosofi}

Studiet af James’ Bog, der udkom et Aar efter mit eget
religionsfilosofiske Forsøg, har været af stor Interesse for mig.
Ikke blot er det længe siden, jeg har læst en Bog, der saaledes
har lært mig at se med nye og friske Øjne paa Livet og paa
Menneskene, som denne Bog af den amerikanske Filosof har
formaaet. Men Bogen fik en særegen Betydning for mig derved,
at hele Retningen og Behandlingsmaaden, trods al Forskellighed
paa ikke uvæsentlige Punkter (f. Ex. i det, James
kalder sin „Supranaturalisme“), i saa høj Grad er beslægtet.
Jeg har derfor en personlig Grund til at henlede Opmærksomheden
paa den. Ved sin aandfulde Form og sine mange
interessante Enkeltheder vil den iøvrigt lønne et Studium, selv
om man staar Bogens Retning mere fjernt end jeg\textsuperscript{13}). Kun
nogle faa Bemærkninger vil jeg her fremføre for at betegne
mit Standpunkt overfor den.

Det er i mine Øjne med fuld Ret, at James skelner mellem
personlig og institutionel Religion. Det er ganske vist en Forskel,
som er lettere at gøre ved de højere end ved de lavere
Religioner, og det er maaske derfor, Religionshistorikere undertiden
krympe sig ved at indrømme den. Men den psykologiske
Interesse maa nu engang efter Sagens Natur nærmest angaa det
personlige Element i Religionen, medens den objektive eller
„institutionelle“ Side kun har Betydning som Symptom. Og
hvad der er Hovedsagen: det religiøse Problem angaar egentligt
% --- p. 81 ---
\opage{81}kun den subjektive Side. Spørgsmaalet om Religionens Stilling ---
nu og i Fremtiden --- i det menneskelige Aandsliv gaar jo dog
nærmest ud paa, om Menneskets Følelse og Trang vil kunne
vedblive at fastholde de hidtil givne Former og Symboler, og
om der, hvis disse skulde vise sig ufyldestgørende, vil kunne
findes eller skabes nye. Den, der fæster sit Blik paa Kontinuiteten
i den aandelige Udvikling, maa afgjort lægge Hovedvægten
paa det subjektive, personlige Omraade, saa meget dette end
tilsløres ved de ydre Formers historiske Overmagt.

Dog forekommer det mig, at James gaar for let henover de
Vanskeligheder, Kontinuitetsproblemet her frembyder. Det er
ingenlunde selvfølgeligt, at det personlige Element i Religionen
kan bevares, dersom de historiske Former falde bort. Som det
vil ses af min Religionsfilosofi, staar jeg her saaledes, at jeg ser
en \textit{Mulighed} for en Bevarelse af den Inderlighed og Koncentration
i Livsopfattelse, der har havt en saa væsentlig Støtte i
de historiske Former for Religion, skønt jeg ikke kan føre noget
tvingende Bevis for Nødvendigheden deraf, --- end ikke for, at
der kan gives Ækvivalenter derfor. Vi staa overfor den Mulighed,
at der her vil gaa Værdi tabt. Historien viser os et stadigt
Vexelforhold mellem det Personlige og det Institutionelle. Skal
dette Vexelforhold bestaa i Fremtiden, saa er det kun muligt,
dersom den menneskelige Aand under Udviklingens Gang bevarer
Evnen til at finde eller skabe store Symboler for de dybeste
Livserfaringer, Symboler, der kunne faa fælles Betydning for
større eller mindre Grupper af Mennesker. Det kan haabes, men
ikke bevises. Og dog er det et af de vigtigste Spørgsmaal med
Hensyn til vor Slægts Fremtid. Man maa se paa Livets største
Konflikter med meget let Sind og meget uklart Blik, naar man
ikke føler Brodden i dette Problem. --- Hvorvidt man under de
nye Forhold vil kunne vedblive at bruge \textit{Ordet} Religion, er et
ligegyldigt Spørgsmaal, som James behandler med den Ironi,
det fortjener. ---

Sammenlignes James’ Beskrivelse af religiøs Erfaring med
den, jeg har givet, fremtræder der en Overensstemmelse, der
viser, at vi have de samme Fænomener for Øje. Religiøs Erfaring
er for ham karakteriseret ved at være Menneskets „totale
% --- p. 82 ---
\opage{82}Reaktion“ overfor Livet; den er en Erfaring af Fred, Enhed
og Energi, ofte vunden gennem en Krise, der førte til Fortrængen
af Mørke og Disharmoni. Hvad jeg her savner, er en
bestemt Fremhæven af, at den Disharmoni, den religiøse Bevidsthed
under alle Former kæmper for at opløse, er en Disharmoni
mellem Værdi og Virkelighed. Ved Forholdet mellem
de Værdier, Mennesket kender, og den Virkelighed, han kender,
er hans Religion bestemt. Det Værdifulde kan for ham være
selve det fysiske Liv i dets Bestaaen og Udfoldelse, eller det
kan være ideelle Goder. Stedse bestemmes hans Totalfølelse
overfor Livet ved, om Værdiernes Himmel er formørket eller
ikke. Det er ud fra dette Synspunkt, at jeg er kommen til min
Hypotese om, at Værdiens Bestaaen er den religiøse Grundtanke
eller det religiøse Axiom. James lægger Vægten paa den indre
Enhedsfølelse og paa Erfaringen om Tilstrømning af indre Energi,
uden at han tager tilstrækkeligt Hensyn til, at denne Enhed og
Energiforhøjelse for Mennesket betyder, at Disharmonien mellem
Værdi og Virkelighed er hævet. Dog tror jeg ikke, at James
vil være uenig med mig i min Hypotese. Som allerede ovenfor
omtalt, antyder han jo, at Religion vel først og fremmest
bygges paa subjektiv Erfaring, men at den mere eller mindre
bevidst hviler paa den Forudsætning, at det Tragiske i Verden
kun kan være partielt og forbigaaende, --- at Skibbrud og Opløsning
ikke kunne være Tilværelsens sidste Ord. Den religiøse
Bevidsthed kommer altsaa ifølge sin Natur til at forudsætte
Noget om kosmiske Forhold, om Tilværelsens inderste Natur.
Det er paa dette Punkt, Religionen stedse har havt Tilbøjelighed
til at blive Metafysik, ligesom omvendt Metafysiken, uden
altid at være sig det bevidst, er gaaet ud fra religiøse Forudsætninger
i sine Konstruktioner. Ifølge James er ethvert Forsøg
paa bestemtere og konkretere Formulering af hine sidste Forudsætninger
„over-belief“, d. v. s. fører fra Tænkningens og Erfaringens
Omraade over i Fantasiens og Mysticismens. James’
egen „Supranaturalisme“ er aabenbart et Exempel herpaa. Om
de nye Begyndelser, han antager, har han ikke Ret til at sige
Mere, end at deres Aarsag foreløbigt ikke kan angives; at de
% --- p. 83 ---
\opage{83}skulde have en overnaturlig Aarsag, er ikke motiveret efter hans
egen Opfattelse. ---

Der er endnu et væsentligt Punkt, i hvilket jeg stemmer
overens med James, nemlig i den store Vægt, han lægger paa
de individuelle Forskelligheder i Henseende til religiøs Erfaring
og Tro. Hans mesterlige Skildring af Modsætningen mellem de
sunde og de syge Sjæles Maade at tage Tingene paa, er et stort
Indlæg til Fordel for, hvad jeg kalder Personlighedsprincipet.
Denne Modsætning svarer til den allerede ovenfor omtalte Modsætning
mellem expansive og disharmoniske Naturer. Men James
har ført Karakteristiken igennem med en saadan Kraft og Klarhed,
at Spørgsmaalet om, hvorvidt forskellige Personligheder
virkeligt kunne have „samme“ Religion, melder sig med endnu
større Nødvendighed end før. ---
% --- p. 84 ---
\kapitel{}{Noter.}{Noter}\opage{84}
\streg

1) S. 1. --- Smlgn. min Afhandling \textit{Filosofien og Livet.} (Tilskueren 1901).

2) S. 4. --- L. P. \textsc{Larsen}: \textit{Høffdings Filosofi blandt Indiens Studenter.}
(Studenterhjemmets Julebog 1901). --- Jeg kan ikke undlade at bemærke, at den
pompøse Titel, denne Afhandling har faaet, vækker falske Forestillinger baade
om „Høffdings Filosofi“, om „Indiens Studenter“ og --- om selve Afhandlingens
Indhold. Derimod passer det ret godt til Afhandlingens Karakter, at den har
fundet sin Plads i et opbyggeligt Skrift. Forfatteren siger om mig, at min Interesse
er mere religionsfilosofisk end religiøs; det kan jeg kun betragte som en
Kompliment, da jeg netop har villet skrive en Religionsfilosofi. Til Gengæld
kan jeg nok sige, at hans Interesse for det Opbyggelige er større end hans
Interesse for det Filosofiske. Det er maaske derfor, at han heller ikke læser
filosofiske Bøger med fornøden Opmærksomhed. Saaledes siger han, at jeg
p. 214 i min Religionsfilosofi „protesterer mod Skabelsesdogmet“. Der findes
ingen saadan Protest i min Bog. Paa det anførte Sted drøfter jeg den Vanskelighed,
en stor religiøs Genius som Augustinus har fundet i Skabelsen af en
endelig Verden, fordi der derved vilde ske en Værdiforringelse, og jeg ser deri
et Vidnesbyrd om, at Troen paa Værdiens Bestaaen er den religiøse Grundtanke.
Hvis Afhandlingens Forfatter ikke anerkender hin Vanskelighed, maa
han afgøre Sagen med Augustinus. Foreløbigt holder jeg paa den gamle Kirkefader,
som baade kunde \textit{erfare og tænke.} --- Naar jeg her har maattet afvise
en teologisk Kritikers Forstaaelse af min Bog, er det mig en saa meget des
større Glæde at henvise til to teologiske Kritikere, en dansk og en tysk, der
begge paa klar og nøjagtig Maade gøre Rede for Indholdet af det Skrift, de
ville bedømme, ligesom ogsaa deres kritiske Bemærkninger paa flere Punkter
have været lærerige for mig: F. C. \textsc{Krarup} i Dansk Tidsskrift April 1901, og
\textsc{Friedrich Lipsius} i Protestantische Monatshefte. 6. Jahrg. Heft 7.

3) S. 6. --- I første Udgave af sit \textit{System der Philosophie} (p. 444) havde
\textsc{Wilhelm Wundt} udtalt, at Filosofien kunde bevise Troens Nødvendighed
(skønt den ikke kunde omsætte Troen til Viden). I anden Udgave er denne
Sætning udeladt, uden at der er givet nogen Grund derfor. Muligvis er Udeladelsen
foranlediget ved en stigende Bevidsthed om den rent empiriske Betragtningsmaade,
vi her maa stanse ved.

% --- p. 85 ---
\opage{85}4) S. 10. --- \textsc{Dowden}: \textit{Puritan and Anglican.} London 1900. p. 234.
--- Smlgn. min \textit{Religionsfilosofi.} p. 278.

5) S. 15. --- For Tydeligheds Skyld hensætter jeg Siebecks Udtalelse i
ordret Oversættelse: „Naar Personligheden med Frihed og Bevidsthed afviser
det overleverede eller uvilkaarligt udviklede Trosindhold, idet den ligefrem erklærer
Troen for Noget, der ikke skal være, og holder den videnskabelige Fremgangsmaade
for den ene berettigede, da ligger herved en Illusion til Grund.
Personligheden bliver sig ikke bevidst, at denne Akt, der er en Frihedsakt, kun
bliver mulig ved et Moment, der hører til Troens Væsen .... Ti Troen paa
et højeste Gode og paa en Værdi, som ved det kan blive Personligheden til
Del, ligger stedse forud til Grund for Afgørelsen, kun at det højeste Gode i
dette Tilfælde udelukkende sættes i den videnskabelige Erkendelses Bestaaen
og Sammenhæng. Men denne sidste Antagelse er ikke selv først Resultatet af
en videnskabelig Bevisførelse; ti uden den vilde der aldrig være opstaaet nogen
Tilskyndelse til videnskabelig Virksomhed“. H. \textsc{Siebeck}: \textit{Religionsphilosophie.}
p. 200.

6) S. 19. --- Smlgn. hermed \textsc{Schleiermacher’s} Opfattelse af Bønnen \textit{(Der
christliche Glaube} § 146---147). En Teori om Bønnen, der forudsætter et Vexelforhold
mellem Skaberen og det Skabte, betegner ifølge Schleiermacher en Overgang
til det Magiske. Der er en indre Sammenhæng mellem Bøn og Bønhørelse;
begge ere i Grunden Et, og ligeledes ere i den rette Bøn Passivitet og
Aktivitet forenede. Hovedsagen er, at selve Bønnen ligesaa vel hører til „den
guddommelige Verdensorden“ som dens Bønhørelse eller Ikke-Bønhørelse. (ibid.
§ 47). --- En amerikansk Læge har nyligt udtalt sig saaledes om Bønnen:
„Maaske er selve Bønnens Akt den Vej, ad hvilken jeg hjælpes af den guddommelige
Magt: naar jeg taler til Gud, er det Gud, der taler i mig. Forunderligt
er det, at skønt jeg ikke tror paa Mirakler, hverken i Fortid eller i
Nutid, tror jeg dog stadigt, at Gud maa have en eller anden Vej til at virkeliggøre
mit Ideal af guddommelig Menneskelighed“. (Anført hos \textsc{Leuba}: \textit{The content}
\textit{of religious consciousness.} The Monist. July 1901. p. 553).

7) S. 21. --- I en vis Forstand ere ogsaa alle videnskabelige Begreber tilsidst
symbolske. Om Forskellen mellem den videnskabelige Symbolisme og den
religiøse (eller metafysiske) se \textit{Religionsfilosofi} p. 70---72. --- \textit{Filosofiske Problemer}
(Universitetsprogram 1902). p. 37 f.

8) S. 37. --- Smlgn. \textit{Det psykologiske Grundlag for logiske Domme.} p.
47 f. --- \textit{Filosofiske Problemer.} p. 75 f.

9) S. 38. --- Smlgn. hermed den yngre \textsc{Ritschl’s} Skrift \textit{Ueber Werturteile}
(Freiburg 1895). Kap. 5 (Die geschichtlich gegebenen geistigen Werte im Unterschied
und im Verhältniss zu der subjectiven Wertbeurteilung).

10) S. 40. --- Se \textit{Religionsfilosofi.} p. 111---115; 251---257. Smlgn. ogsaa
\textsc{Leuba’s} Kritik af Murisier i Revue philosophique. Juillet 1902 p. 33---35.

11) S. 54. --- Maaske der her kunde være at skelne mellem Kirkens Ledere
og dens Troende. Dens Ledere have ofte været afgjort koleriske og aktive
Naturer; men de Troende vare som saadanne henviste til overvejende Passivitet,
og det var Sorgens og Haabets Stemninger, der gennem Kultus dyrkedes hos dem.

12) S. 74. --- Paa dette Punkt slutter James sig til Leuba. „Leuba is
undoubtedly right in contending that the conceptual \textit{belief} . . . although so often
efficacious and antecedent, is really accessory and non-essential, and that the
joyous conviction can also come by far other channels than this conception.
It is to the joyous conviction itself, the assurance that all is well with one, that
% --- p. 86 ---
\opage{86}he would give the name of \textit{faith} par excellence.“ \textit{(Varieties of religious belief.}
p. 246 f.). --- Smlgn. hermed \textit{The Will to believe.} p. X: „What mankind at
large most lacks is criticism and caution, not faith. Its cardinal weakness is
to let belief follow recklessly upon lively conception, especially when the conception
has instinctive liking at its back“. James føjer til her, at hvis han
talte til Frelsehæren eller lignende Kredse, vilde han ikke betone Villien og
Retten til at tro; saadanne Kredse behøvede snarere „Videnskabens friske Nordvestvind“
til at feje deres Sygelighed og Barbari bort. Men han taler til akademiske
Kredse, hvis Fornødenheder ere andre: „Paralysis of their native capacity
for faith and timorous abulia in the religious field are their special forms
of mental weakness“.

13) S. 80. --- Af filosofiske Kritiker af James’ Bog kender jeg kun to, en
af Genferpsykologen \textsc{Th. Flournoy} i Revue philosophique (Novbr. 1902) og
en af J. H. \textsc{Muirhead} (Professor i Birmingham) i International Journal of
Ethics (January 1903). Flournoy slutter sig med Begejstring til James; Muirhead
anerkender ham som Psykolog, men er misfornøjet paa Metafysikens Vegne.

\end{document}
